\documentclass[11pt,a4paper]{article}
\usepackage[utf8]{inputenc}
\usepackage[T1]{fontenc}
\usepackage{lmodern,amsmath,amssymb,textcomp}
\usepackage[margin=24mm]{geometry}
\usepackage{longtable,booktabs,array,enumitem}
\usepackage[hidelinks]{hyperref}
\setlength{\parindent}{0pt}
\setlength{\parskip}{6pt}
\emergencystretch=3em
\begin{document}
{\LARGE\bfseries Formalizing Established Mathematics at OpenMath 2026: Reusable Libraries, Exact Certificates, and Research Boundaries\par}\medskip
{\small Prepared for author review from the preserved competition packets, 5 October 2026
Editorial compilation by Alejandro Zarzuelo Urdiales\par}\medskip
{\small Judging and evaluation record: the submissions discussed in this document have been judged and evaluated by Alejandro Zarzuelo Urdiales for OpenMath 2026. This dated review records the stated verification, classification and unresolved holds; it is a preliminary evaluation rather than a signed award decision.\par}

{\small Event provenance: OpenMath 2026 ran 27 September-2 October 2026 with worldwide online participation. Detailed organizer listings specify Harvard Innovation Labs for the 27 September in-person event; the OpenMath programme advertises a hybrid kickoff at MIT. Sundai identifies its Harvard/MIT community. Organizer sources: \url{https://rsihouse.ai/openmath} ; \url{https://luma.com/yzp9abvr} ; \url{https://www.sundai.club/events/boston/recursive-learning-hack-with-harvard-innovation-labs}\par}

Central preserved source archive: \url{https://github.com/alejandrozu/openmath-2026-judging}

\clearpage
\section{Contents}
\tableofcontents
\section{Abstract}
The mathematical value of an AI-assisted competition extends beyond new numerical bounds and newly resolved questions. Its preserved artifacts also include formalizations of established theorems, interfaces between different definitions of arithmetic progression, exact computational certificates, and incomplete research programmes whose limitations are informative. This paper records that part of the OpenMath 2026 corpus. It describes ten selected Sakana formalization families; HTPeo's perfect-matching and sixteen Erdős families, distinguishing substantive developments from smaller helpers; the arithmetic-progression developments of Leon Koerbs and Cameron Fen, Wilson Wu, and Jamie Steeg; and Leanification's formalization of the Campos-Samotij container theorem. It also examines Collatz descent certificates, finite Busy Beaver runs, known matrix and Kobon constructions, and Ramsey and Grothendieck certificates whose claimed original-mathematics status is unresolved. Established mathematics is attributed to its prior authors. Potential formal-library novelty is a separate question, requiring exact-source compilation, statement correspondence and prior-formalization comparison. An appendix keeps unfinished unitary-synthesis and geometric research separate from proved known results. The document is an author-review draft, not a declaration of journal acceptance or first-formalization priority.

\section{1. Purpose, Scope And The Meaning Of Formalization Value}
A proof assistant can make an established theorem useful in a new way. A reusable formal statement may let future developments invoke a theorem with all its assumptions visible, compose its conclusion with other libraries, or check a computational certificate without trusting its search procedure. These are valuable outcomes even when the informal mathematics was already known. Their value does not depend on presenting them as solutions of open problems. The appropriate scientific questions are which mathematical statement was formalized, how faithfully it matches the source, what reusable infrastructure was added, and what was already available in the dependency libraries.

The competition record contains several layers of evidence. A theorem statement and source file identify a candidate formalization. A supplied successful build log records what an author's environment reported. A fresh build at the exact submitted source and dependency pins is stronger reproduction evidence. An endpoint axiom audit identifies admitted dependencies, while mathematical review establishes whether the machine-checked statement expresses the intended mathematics. These layers should remain visible in publication. A compiler success does not by itself establish novelty, appropriate attribution, or equivalence between a finite encoding and the ordinary mathematical object.

This document uses the preserved judging repository as its central source: \url{https://github.com/alejandrozu/openmath-2026-judging} . The audited archive revision is d0ed487f487fa7ec814ff705ab55249983fe8707. The original packets and statement cards, rather than a leaderboard number, govern the descriptions. Fresh replay results are recorded in the companion verification ledger supplied with the publication package. Where this paper refers to earlier replay receipts, they remain receipt evidence unless that ledger records an independent rerun. No result is upgraded merely because its source contains the words 'proved' or 'no sorry'.

One mathematical family is described once, even if it contains many helper lemmas, alternate endpoints or specializations. Independently produced formalizations can deserve separate contributor recognition, but a collective paper should not pretend that equivalent interfaces represent several different mathematical discoveries. The source formalization authors and the prior informal theorem authors are separate attribution categories. AI systems are disclosed as tools; human authors remain responsible for the submitted mathematical statements and their publication.

The chair has approved the October 4 competition amendments and has retained the DMS progress value of 0.15. Those administrative choices are reported in the separate short judging document. They do not turn established mathematics into novel mathematics. Leanification is included here as a mention and a substantive formalization case study, following the chair's latest instruction, without a competitive placement. The long original-results paper is a companion volume, so the material here need not compete rhetorically with novel contributions.

\section{2. Sakana: Ten Selected Formalization Families}
The Sakana packet is a joint submission of Rachel Teo, Wenyi Wang and Hamdi Abderrahmene. The supplied record does not justify a per-result division of their human contributions. Its known-mathematics material ranges from Bernoulli denominators and quaternion equations to digit carries and geometric counting. These are ten theorem families, not ten completely solved parent open problems. Their potential publication form is a coherent formalization catalogue with reusable definitions and proof interfaces. Individual sections may also support focused library notes if exact prior-formalization comparisons establish an appropriate new delta.

\subsection{2.1. A046969 and Bernoulli denominators}
For n>0, define a(n) as the denominator of B\_(2n)/(2n(2n\ensuremath{-}1)), with a(0)=0. The selected theorem states that if p>3 and both p and 2p\ensuremath{-}1 are prime, then a(p)=12(2p\ensuremath{-}1). In particular, a(p)/12 is prime. This is denominator mathematics with a classical source chain, rather than an unrestricted new theorem about Bernoulli-number prime values. The mathematical priority comparison includes Powell-Frame and the classical Bernoulli denominator theory recorded in the packet.

The useful formal task is to connect rational denominators, Bernoulli-number identities, primality hypotheses and an exact natural-number formula. A small-looking final theorem may require careful treatment of signs, rational reduction and the exceptional small primes. Its reuse lies in a verified denominator interface and exact hypotheses. The existing Mathlib Bernoulli infrastructure should receive credit, and the new contribution should be measured against that infrastructure rather than the number of final declarations. Sources: \url{https://oeis.org/A046969} ; \url{https://dlmf.nist.gov/24.10} ; judging card E02-a 046969-i.

\subsection{2.2. A two-witness quaternion scalar equation}
The packet formalizes an explicit existential definition of scalar Hamilton quaternions. Over a linearly ordered commutative ring with the required strict-order compatibility, q is scalar precisely when there are unrestricted quaternions u,v satisfying ((uv\ensuremath{-}vu)\textasciicircum{}2+4)\textasciicircum{}2\ensuremath{-}(qu\ensuremath{-}uq)\textasciicircum{}2\ensuremath{-}(qv\ensuremath{-}vq)\textasciicircum{}2=0. The integer Lipschitz order and rational Hamilton quaternions are specializations of the same generic result. The equation is the central object; neither the number of witnesses nor the scalar domain should be silently altered.

This development makes a known Diophantine encoding explicit in a formal algebraic setting. It can be useful to future work on existential definability and noncommutative arithmetic, since it isolates a concrete polynomial equation and its interpretation. It does not solve Hilbert's tenth problem over the rationals. The formal novelty inquiry should compare the selected equation and generic proof with prior quaternion algebra, norm and commutation APIs, which are already available. Source: \url{https://aimath.org/WWN/hilberts10th/hilberts10th.pdf} ; judging card E02-opdp67-quaternion.

\subsection{2.3. A135508, an lcm recurrence and prime-index criteria}
Let x(0)=0, x(1)=1, and x(n+1)=2x(n)+lcm(x(n),n+1) for n\ensuremath{\geq}1. Write a(n)=x(n+1)/x(n)\ensuremath{-}2 for n>0. For a prime p>3, the packet proves that a(p\ensuremath{-}1)=p exactly when gcd(x(p\ensuremath{-}3),p\ensuremath{-}2)>1. It also gives sufficient residue conditions: for p\ensuremath{\geq}7, either p mod 3=2 or p mod 5=2 implies a(p\ensuremath{-}1)=p. A more general acquired-divisor condition transfers a divisor already present at an earlier recurrence stage to the desired prime index.

The selected scope is a criterion and sufficient classes, not an unconditional characterization for every prime or a complete parent recurrence conjecture. Its formalization value lies in controlling an integer recurrence whose defining lcm creates persistent divisibility information. Recurrence inequalities, exact divisions and gcd conditions can become reusable lemmas. Classification is conservative because earlier conditional work and elementary specializations overlap. If the precise unconditional criterion is later shown to be new mathematics, it should be reclassified once rather than simultaneously counted as known and original. Sources: \url{https://oeis.org/A135508} ; \url{https://arxiv.org/html/2510.18891v3} ; judging card E02-a 135508.

\subsection{2.4. Apéry prime-power polynomial irreducibility}
For B\_n(X)=\ensuremath{\Sigma}\_(k=0)\textasciicircum{}n binom(n,k)\textasciicircum{}2 binom(n+k,k)X\textasciicircum{}k in Q[X], the selected theorem gives irreducibility of B\_(p\textasciicircum{}r) whenever p is an odd prime, r\ensuremath{\geq}1 and p\ensuremath{^2} does not divide 2\textasciicircum{}(p\ensuremath{-}1)\ensuremath{-}1. This is a non-Wieferich sufficient hypothesis. It should not be replaced by a claim covering every odd prime, nor by an unrestricted higher-valuation formula.

An earlier exposed informal sketch overlaps the odd-family mathematics, while a previously completed binary-power chain is not part of the new selected contribution. The meaningful formal delta is the odd-prime-power route under its correct sufficient assumption. This is a good example of why informal priority and formal priority differ: a valid sketch can defeat a first-mathematical-result claim without already supplying a checked endpoint in Lean. Publication should explain the coefficient and valuation interfaces, credit the exposed argument, and identify which new formal modules are reusable. Sources: \url{https://oeis.org/A005258} ; the pinned Epoch LeanOpenProblems-results comparisons recorded in E02-apery-odd-formalization.

\subsection{2.5. Base-three reverse-and-add carry obstruction}
The selected statement says that if x plus its canonical base-three digit reversal has palindromic base-three digits, then either x already has palindromic digits or every mirrored pair of input digits has sum less than 3. It is a carry obstruction for one reverse-and-add operation. It does not prove that arbitrary repeated reverse-and-add dynamics reach a palindrome, and it does not transfer automatically to decimal arithmetic.

The interesting formal content is the bridge between natural-number arithmetic and finite digit lists, including canonical digit representation, length changes, reversal and mirrored carries. These are delicate boundaries for automated proofs: leading zeroes and a final carry can change the apparent symmetry. A reusable list-level carry analysis is more valuable than simply listing a few palindromic examples. The packet records the selected endpoint and supporting digit lemmas under the base-three research formalization family; an exhaustive first-formalization claim is not established.

\subsection{2.6. Theta half-sum cancellation}
For an odd prime k and odd naturals a,b with 1\ensuremath{\leq}a,b<k and ab\ensuremath{\equiv}\ensuremath{-}1 mod k, define H(k,a)=\ensuremath{\Sigma}\_(j=1)\textasciicircum{}((k\ensuremath{-}1)/2)(\ensuremath{-}1)\textasciicircum{}floor(aj/k). The selected theorem gives H(k,a)+H(k,b)=0. The hypothesis is negative inverse congruence, and the summation domain is exactly the lower half of the nonzero residues. A change to an ordinary inverse or a full residue sum would be a different statement.

This is useful as a formal finite-sum identity with number-theoretic symmetry. The source decomposes parity counts and modular inverses, so the development can support later reciprocity and cancellation arguments. Such an identity may have a very short informal proof but still require substantial indexing and floor-arithmetic infrastructure. Its role in this volume is a faithful formalization of known-method mathematics; the exact formal-library novelty remains to be compared with older modular and finite-sum developments.

\subsection{2.7. Reciprocal-square numerator arithmetic}
For m\ensuremath{\geq}1, the rational sum \ensuremath{\Sigma}\_(k=1)\textasciicircum{}m1/k\ensuremath{^2} has nonpositive 2-adic and 3-adic valuations. The packet also supplies a bridge from nonpositive valuation to exclusion of a prime from the reduced numerator, and a prime-index statement gcd(p,A007406(p\ensuremath{-}1))=p for every prime p>3. The selected scope is these valuations and prime-index arithmetic. It is not the full composite-index or stronger Wolstenholme target.

Existing formal harmonic-number valuation theory overlaps some infrastructure but does not automatically supply the reciprocal-square theorem. The main selected content includes an equal-minimum 3-adic noncancellation argument and the numerator/divisibility bridge. These are reusable interfaces for rational sums, valuations and reduced numerator conventions. Publication should state which preexisting p-adic facts are imported and which finite-sum lemmas were added. The family counts once even when its valuation and numerator endpoints are in separate files. Source: judging card E02-m2-research-wolstenholme and the packet's selected proof modules.

\subsection{2.8. A060841 reciprocal determinant integrality}
For the n\ensuremath{\times}n matrix with entries 1/lcm(i+1,j+1), the inverse determinant is an integer exactly when n belongs to \{1,...,34,36,38\}. The informal classification was already public. Its formalization involves a determinant-to-product correspondence and a valuation-based classification; the packet identifies overlap in preexisting product and valuation helpers.

The reusable contribution is the verified connection from the actual rational matrix determinant to the arithmetic criterion. It should not be presented as a new numerical classification merely because the final finite list is explicit. This is a suitable library case study because it shows how linear algebra, divisibility and finite exceptions interact. The endpoint also needs its nonzero determinant and inverse conventions documented so that totalized inversion cannot conceal a vacuous statement. Source: \url{https://github.com/umaia1234/agentic-conjectures/blob/b89af1ced140d79c4d04ac84e39f00181bb6540a/problems/oeis-a060841/PROOF.md} ; judging card E02-m2-research-a060841.

\subsection{2.9. Four divisors just above a square root}
The packet formalizes an Erdős-Rosenfeld construction giving infinitely many natural n with at least four distinct divisors in the open interval (sqrt(n),sqrt(n)+40n\textasciicircum{}(1/4)). Its preserved integer version uses floor and ceiling endpoints equivalent to that real interval. The theorem is an infinitude statement about a specified short interval, not an unrestricted classification of nearby divisors.

The formal engineering problem is to preserve strict inequalities, distinctness and divisibility while moving between a real scale and exact integers. Such a bridge is useful for many short-interval arithmetic problems. The original construction must be credited to its prior source rather than rediscovered under a new name. The event contribution is the faithful selected formal proof and its reusable arithmetic infrastructure, subject to comparison with existing developments. Source: judging card E02-m2-research-erdos887.

\subsection{2.10. Regular-polygon distance multiplicity}
For even n\ensuremath{\geq}5, every distance occurring between distinct vertices of the regular n-gon is attained by at most n unordered vertex pairs. The statement concerns regular polygons, not arbitrary convex point sets or the full Erdős 94 problem. The use of unordered pairs is essential: counting ordered pairs changes the numerical bound by a factor of two.

Generic two-circle-intersection uniqueness is already in Mathlib. The selected family adds the connection to regular-polygon geometry and a complete counting argument. It can serve as a compact verified example of how symmetries, finite sets and geometric intersection bounds compose. The scientific description should therefore identify the actual regular-polygon model and counting conventions instead of treating a one-line final inequality as evidence of either triviality or broad novelty. Source: judging card E02-m2-research-regular-polygon.

\section{3. Htpeo: Selected Known Theorems And Smaller Helpers}
The HTPeo packet names Woohyuk Kang, Hyunjin Lee, Sanghyeon Lee and Youchan Oh. Its selected known-theorem formalization sessions and packet assembly are attributed to Woohyuk in the pinned TEAM.md. The collective record includes one perfect-matching family and sixteen Erdős families. Ten were proposed as substantive, while seven were held for usefulness review; within the ten, the powerful-number examples for r=7,8 may themselves be more suitable as examples than as a standalone family. This paper records all seventeen without deciding their competitive counts from the number of files. Six explicitly unclaimed wrappers remain excluded.

\subsection{3.1. Petersen and Schönberger perfect matchings}
The selected formal theorem concerns connected, bridgeless, cubic simple graphs and perfect matchings with the specified edge-inclusion property. Cubicity, bridge-freeness and the simple-graph convention are part of the statement. The formal predicate Bridgeless checks absence of graph bridges, so the theorem is not merely a declaration that a convenient matching exists.

This is established graph theory, useful both on its own and as infrastructure for the separate DMS work. The correct publication practice is to credit the prior matching theorem and document the added graph and matching interface. Its reuse within DMS does not create a second mathematical result or justify counting an imported dependency again. A future library integration should make the finite vertex, connectivity and regularity assumptions straightforward to invoke. Sources: pinned HTPeo packet, entries E08-G-PM and entries/dms-star 6.

\subsection{3.2. Powerful numbers between consecutive squares: Erdős 942}
A powerful number has every prime factor occurring with exponent at least two. The selected formal result states that the count h(n) of the relevant powerful numbers between consecutive squares has unbounded limsup. It is an unboundedness theorem, not a claim that every such interval contains many powerful numbers.

The usefulness lies in the precise counting function and the translation of arbitrarily large counts to an extended-natural limsup. Future arithmetic developments can reuse that interface without reconstructing the asymptotic quantifiers. The original theorem is known; formal novelty requires checking earlier powerful-number libraries and exact endpoint overlap. Source: HTPeo E08-E942.

\subsection{3.3. A greedy Sidon lower bound: Erdős 44}
For N\ensuremath{\geq}1, the selected theorem constructs a Sidon subset A of \{1,...,N\} with N\ensuremath{\leq}|A|\ensuremath{^3}. It records the classical scale produced by a greedy exclusion argument. A Sidon set forbids repeated additive representations under the source convention, which must remain aligned with the intended mathematical notion.

This is a natural reusable finite-combinatorics formalization. The proof pattern can support later greedy constructions by relating maximality, forbidden choices and cardinality. Its merit is a verified combinatorial argument and useful API, not a new asymptotic Sidon bound. Source: HTPeo E08-E44.

\subsection{3.4. Completeness of powers of 2 and 3: Erdős 123}
The exact definition is: A is d-complete when every sufficiently large natural n is a sum of a finite set s contained in A, with distinct x,y in s satisfying neither x divides y nor y divides x. The summands form a divisibility antichain. The selected A=\{2\textasciicircum{}a 3\textasciicircum{}b:a,b nonnegative integers\} satisfies this property; the proof constructs representations for every n, including the empty sum for 0. This is established mathematics, here implemented as reusable antichain-sum infrastructure.

Its reusable content is the link between generated sets and the precise representation predicate. This can be valuable in formal additive number theory because finite-subset and set-product conventions are easy to conflate. The statement does not establish completeness for arbitrary multiplicatively generated sets. Source: HTPeo E08-E123.

\subsection{3.5. Impossible countable chromatic restrictions: Erdős 918}
The selected set-theoretic graph result rules out the specified graph of cardinality aleph 2 and chromatic cardinal aleph 2 when all induced subgraphs on aleph 1-sized subsets are required to have the exact restricted chromatic cardinal in the source statement. The word 'exact' matters: upper bounds and equalities in infinite chromatic number are different hypotheses.

This development is useful for cardinal graph interfaces, induced subgraphs and chromatic-cardinal arguments. Its source contains a boundary convention at n=0 that must be disclosed wherever the corresponding generalized endpoint is used. It is known mathematics, not a new theorem resolving arbitrary uncountable chromatic phenomena. Source: HTPeo E08-E918.

\subsection{3.6. Multiplicative closure and excluded prime powers: Erdős 292}
The selected family proves the source-defined set A is closed under multiplication, contains 2m when m\ensuremath{\in}A and m>1, and has the specified prime-power exclusion. The same set definition must govern all clauses. These are structural statements about one arithmetic object, so their internal lemmas and corollaries should remain one family.

The reusable benefit is a package of closure and obstruction properties, allowing later proofs to move beyond direct evaluation of the defining predicate. Here A consists of possible largest denominators in distinct-unit-fraction expansions of 1: a finite S contained in\{1,...,n\}, containing n, satisfies sum\_\{m in S\}1/m=1. Multiplicative closure replaces denominator n in one expansion by ns from an expansion with maximum m, producing maximum mn. Doubling uses 1/2+sum\_\{s in S\}1/(2s)=1. Prime-power maxima p\textasciicircum{}k are excluded because 1/p\textasciicircum{}k has uniquely lowest p-adic valuation among the distinct smaller positive denominators, preventing cancellation to 1. These structural facts are the claimed scope; Martin's density-one theorem and sharp complement asymptotic are not claimed by this packet. Source: HTPeo E08-E292.

\subsection{3.7. A radius-one Littlewood-Offord counterexample: Erdős 395}
The selected endpoint refutes the proposed uniform c/n lower bound for the source's radius-one signed-sum probability of unit complex vectors. Its scope is the exact radius and closed-ball convention in the formal count. It should not be described as refuting every reverse Littlewood-Offord inequality.

This is useful as a verified negative example and as infrastructure connecting combinatorial sign assignments to probabilities in the complex plane. A counterexample formalization can prevent later theorem statements from omitting a necessary radius or nondegeneracy hypothesis. Its mathematical outcome is already known. Source: HTPeo E08-E395.

\subsection{3.8. A binomial gcd inequality: Erdős 698}
For 1\ensuremath{\leq}i<j\ensuremath{\leq}n/2, the selected Erdős-Szekeres inequality gives binom(n,i)/binom(j,i)\ensuremath{\leq}gcd(binom(n,i),binom(n,j)). It mixes exact integer divisibility with a real-valued ratio. The source casts and nonzero denominator assumptions therefore belong in the formal statement.

This is a substantive reusable arithmetic family: binomial identities, gcd divisibility and an inequality endpoint can be composed in later combinatorial-number-theory arguments. It is not a new estimate merely because its formal proof did not already appear under the exact declaration name in Mathlib. Source: HTPeo E08-E698.

\subsection{3.9. Powerful-number sum examples: Erdős 939}
The r=7 and r=8 endpoints show that the relevant powerful-number sum sets are nonempty by known explicit examples. These may be useful as verified test cases, particularly for making sure a formal definition matches the intended sum problem. They do not give an unbounded family or settle the unresolved ranges.

The correct level of recognition may be smaller than a substantive standalone formalization family. If the proof consists principally of checking a literal example with existing arithmetic, its chief scientific use is regression testing and statement validation. It belongs in the catalogue, with its exact parameters visible, without making the count of examples a proxy for new mathematics. Source: HTPeo E08-E939.

\subsection{3.10. Nonunique polynomial-range representations: Erdős 477}
The square-range endpoint states that for any A\ensuremath{\subseteq}Z, some z does not have a unique representation as a plus a square from A and the polynomial's integer range. A selected quadratic extension treats ax\ensuremath{^2}+bx+c with a\ensuremath{\neq}0, b\ensuremath{\neq}0 and a dividing b. Its quantifiers distinguish failure of unique representation from a claim that no representation exists.

The formalization creates interfaces for polynomial evaluation ranges, Cartesian-product representations and uniqueness. This is reusable additive algebra, especially because a uniqueness predicate can be easier to misstate than an existence predicate. The variants remain a single family. Source: HTPeo E08-E477.

\subsection{3.11. Seven smaller helpers retained in the catalogue}
Erdős 295 supplies distinct increasing denominators, all at least an arbitrary N, whose reciprocal sum is 1. This is an Egyptian-fraction existence helper. Its mathematical utility is a clean finite decomposition that can later be inserted into constructions with a large-denominator requirement. Whether it warrants standalone formalization recognition depends on what genuinely new reusable API it adds.

Erdős 703 gives the known t=0 value T(n,0)=2\textasciicircum{}(n\ensuremath{-}1) for the source's intersecting-family function, with n\ensuremath{\geq}1. Erdős 748 gives the known sum-free-set lower bound 2\textasciicircum{}(n/2)\ensuremath{\leq}f(n). These are boundary and baseline theorems, useful for checking an extremal definition and for later induction or asymptotic comparisons. They should not be presented as resolution of all parameters of their parent problems.

Erdős 1136 formalizes the multiples-of-three density baseline: that set avoids the specified powers-of-two pattern and has density 1/3. Erdős 358 identifies the number of consecutive-sum representations with the number of odd divisors. Both offer useful arithmetic/density interfaces, while the underlying constructions and identities are classical.

Erdős 619 extends a connected triangle-free graph to a maximal triangle-free supergraph with the specified diameter bound. Erdős 1148 records a known weaker difference-of-squares bound. The first can be useful as an existence/reduction lemma, while the second supplies a baseline against which later improvements can be tested. These seven families are included because the corpus contains them, but inclusion in a scientific catalogue does not by itself award points. Sources: HTPeo E08-E295, E703, E748, E1136, E358, E619 and E1148.

\section{4. Arithmetic-Progression Infrastructure Across Three Entries}
Erdős 3 asks whether every set of positive integers with divergent reciprocal sum contains nonconstant arithmetic progressions of every finite length. None of the packets described in this section establishes the full statement. Their contribution is to make important equivalences, known cases and analytic tools explicit. They should be published with hypotheses printed beside the conclusions, so a conditional implication cannot be mistaken for a closed solution.

\subsection{4.1. Leon Koerbs and Cameron Fen: the dyadic equivalence}
For each K\ensuremath{\geq}3, their closed endpoint identifies an exact equivalence between summability of \ensuremath{\Sigma}\_j r\_K(2\textasciicircum{}j)/2\textasciicircum{}j and the assertion that every reciprocal-divergent set contains a K-term progression. Here r\_K(N) is the maximal cardinality of a K-AP-free subset of \{1,...,N\}. The forward direction relates counting in dyadic blocks to reciprocal sums. The converse constructs separated narrow windows beginning at 4K\ensuremath{\cdot}8\textasciicircum{}j, of width 8\textasciicircum{}j, and controls cross-window progressions.

The theorem characterizes the missing density input; it does not prove that input for every K. The source also contains earlier parent attempts ending in labelled proof holes. Those holes must not be advertised as solved, but they do not invalidate a separately closed named endpoint outside their dependency chain. This is useful reusable equivalence infrastructure because future density improvements can be connected to the reciprocal formulation without repeating the entire transfer argument. Source: \url{https://github.com/Koerbser/erdos-3-autolab/blob/551e388b0a7a933e68cadbab099342af326c1dc2/lean/climb1-library/Converse.lean} .

\subsection{4.2. Their four-term generalized von Neumann inequality}
For the cyclic group ZMod n, with 2 and 3 invertible, and the first three functions bounded by 1, the selected theorem gives ||apCount 4(f1,f2,f3,f4)||\textasciicircum{}8\ensuremath{\leq}n\textasciicircum{}12 Re(gowers 3(f4)). The final function need not be bounded because the inequality is homogeneous in it. The sums in the packet are unnormalized; the factor n\textasciicircum{}12 is therefore part of the theorem, not a decorative normalization choice.

This is meaningful analytic groundwork based on repeated Cauchy-Schwarz and reindexing. Independent finite tests at n=5 and 7, including unbounded fourth functions, support the normalization and model interpretation, but finite tests do not prove the universal result. Its potential use is substantial: a later proof can invoke a correctly normalized inequality instead of developing Gowers machinery from scratch. Source: \url{https://github.com/Koerbser/erdos-3-autolab/blob/551e388b0a7a933e68cadbab099342af326c1dc2/lean/climb1-library/VonNeumann4.lean} .

\subsection{4.3. Their known three-term case and imported deep theorem}
The three-term endpoint proves the known implication from divergent reciprocal sum to a nonconstant three-term progression. It invokes an external formalization of the Kelley-Meka density estimate, then supplies a cardinality/predicate bridge, a dyadic summation argument and the reciprocal transfer. The deep analytic estimate is a load-bearing imported result and must be credited to Kelley and Meka for the mathematics and to the relevant prior formalization authors.

Forty-four required plby modules were omitted from the entrant ZIP and identified by an immutable upstream pin. The record also notes redistribution permission has not been obtained for those files. A publishable packet must preserve that licensing boundary while making reproduction possible through an authorized route. Removing a sorry-backed extension from another dependency does not automatically certify the entire final closure. The exact endpoint axioms must be checked after all imports are assembled. Source: \url{https://github.com/Koerbser/erdos-3-autolab/blob/551e388b0a7a933e68cadbab099342af326c1dc2/lean/three-term/Solution.lean} and ATTRIBUTION.md.

\subsection{4.4. Wilson Wu's equivalence and conditional bridges}
Wilson's source gives the per-length dyadic equivalence for k\ensuremath{\geq}4, using a different window construction based on [4\textasciicircum{}j,2\ensuremath{\cdot}4\textasciicircum{}j). The threshold matters: cross-window three-term progressions can occur, so this particular gluing route should not be described as covering k=3. The closed equivalence is useful even though neither side is proved unconditionally.

The three-term modules define Kelley-Meka and logarithmic-barrier density assumptions, then prove that those assumptions imply the known three-term reciprocal conclusion. A clean axiom list for this implication does not discharge its analytic parameter. Other modules connect extremal cardinality to the canonical Roth-number vocabulary and supply AP special cases and prefix interfaces. These can form one reusable interface family if their new formal delta is clear; density-envelope variants and elementary corollaries should not be counted separately. On 5 October, an independent exact-source Lean 4.33.1 replay compiled the separate Std-only Erdos3SpecialCases module and its requested 14-print audit: five endpoints are axiom-free and nine use only standard kernel axioms. Its literal ContainsAP predicate requires a natural start and a strictly positive natural step, and its all-length predicate quantifies every finite natural length. The checked statements cover monotonicity, prefixes, cofinite and affine-tail cases, positive multiples, two-term progressions in unbounded sets, and the unbounded powers-of-two obstruction to three-term progressions. They prove known elementary mathematics and reusable statement infrastructure, not the reciprocal-divergence conjecture. The isolated Windows replay of the recovered original dependency bytes reached a completed checkpoint at 07:05:52 UTC on 5 October 2026, with 187/199 own source modules checked: 183 granular modules and four whole-resource modules. All 187 successes are recovered FC dependency-library modules. Six further FC modules and the six FC-dependent entrant modules remain pending, together with their six organizer audits and 42 current selected outputs. The earlier exact Std one-source/14-print witness remains separately accepted; the final 200-body/56-output evidence composition has not been qualified. The retained PerfectPower waiting label is an unattempted source, with no active E65 Lean process. This is partial dependency execution and known-mathematics verification, with no full-project PASS, new originality score or unconditional reciprocal-divergence theorem. Source: \url{https://github.com/wilsonwu-ai/sundai-erdos-3/tree/492fd601f9c76e44dcb87f683a5ea6ccb7aaefc9} .

One endpoint uses an 'answer(sorry)' token in a proposition macro. In the copied FormalConjecturesUtil.Answer elaborator, the default option google.answer.alwaysTrue maps a Prop-valued answer(sorry) to True; this is not an admitted proof of the endpoint. Its addDecl creates a safe definition, rather than an axiom declaration. A placeholder used at a non-Prop type can instead retain a sorryAx dependency. These cases must be distinguished by inspecting the elaborated term and expected type. Equally, a clean endpoint axiom list does not by itself authenticate the intended problem semantics: replacing an unspecified proposition with True can change what the displayed statement means. In M2 and skeleton developments, the elaborated theorem must therefore be compared explicitly with the ordinary mathematical target and its hypotheses. Neither the spelling of a macro argument nor a token scan suffices to certify or reject the endpoint. Compiler success, dependency axioms and statement correspondence are complementary checks.

The original dependency environment is now better preserved. A source-only recovery verified the full original container layer against its digest and retained 1,527 source, manifest and license files (4,897,328 bytes), without running the image or retaining Linux binaries or compiled proof artifacts. Its compiler is Lean 4.33.1, and its Mathlib revision is 0df444a360eaa60ab8c11dca51a86af692955474; all eight other dependency pins match the prepared official dependencies. This resolves the former compiler and ordinary dependency-pin uncertainty. It does not identify the deleted whole FormalConjectures Git revision. The original reachable FormalConjectures source closure has 193 modules, whereas the separately prepared compatible utility plan has 233: 173 of those compatible modules are byte-identical to files in the entire recovered source tree; 168 of these matches belong to the original 193-module reachable closure, and the other five lie outside that closure. Consequently, an eventual compatible-source replay must remain distinct from an original-source replay. Independent checking of the recovered bytes has therefore begun in the separate Windows source route, with the completed partial checkpoint recorded above. The original Linux container was not executed by this review. FC sources use explicit pp.unicode.fun=true, autoImplicit=false and relaxedAutoImplicit=false; historical author invocation options remain unverified. The checkpoint establishes neither the unrun selected FC-dependent theorems nor axiom cleanliness of all unselected dependency theorems.

\begin{samepage}
The preserved source ZIP has SHA-256 014333c7edffa94686ef4be2e30e871c0df03ac8feceab1efab53df9793b31e5 and is included in the verification materials.

\end{samepage}
\subsection{4.5. Jamie Steeg's AP conclusion-representation interface}
The selected H7 D4 endpoint proves arithmetic-progression truncation and an equivalence between frequent-atTop existence and existence at every exact finite length. It uses the imported canonical Set.IsAPOfLength definition. The cardinality argument establishes injectivity, thereby preserving nontriviality rather than allowing a constant progression to make the statement vacuous.

D3 repeats the first endpoint and is not another family. D5 contains related finite-deletion, tail and residue-fibre reciprocal lemmas, but it is not part of the submitted H7 selected manifest; it is documented as supporting source rather than silently expanding the claim. The formalization is a narrow conclusion interface. It does not prove a new analytic estimate or solve Erdős 3. The human authority declaration specifies Jamie Steeg as an individual, self-resourced entrant and asserts no organizational affiliation; GCL/GCT references are project provenance only. Source: \url{https://github.com/grandchallenge/MATHSOLVE/blob/2fd9fd643bf3e9e98d121dc2e7af1cf1be411ae2/work_packages/OPENMATH_2026/OM26_H7_ERDOS_3/H7_D4.lean} .

\subsection{4.6. Collective reuse and attribution}
The three entries overlap in vocabulary but differ materially in scope. E11's k\ensuremath{\geq}3 dyadic equivalence is more than an AP-prefix interface; Wilson's conditional transfer remains conditional; Jamie's conclusion equivalence supplies neither density estimate. A collective library can retain separate author provenance while presenting shared mathematical interfaces once. The most useful integration would expose canonical definitions, isolate external deep inputs and make the exact hypotheses obvious in documentation.

Their finite four-AP-free harmonic searches are a separate matter. Leon and Cameron's certified reciprocal sum 4.314 does not exceed Walker's located 4.43975 or the stronger candidate in the novel-results volume. It can be published as a bounded computational benchmark if its finite-N target is clearly specified, but it is not a new unrestricted lower bound and does not supply a divergent reciprocal set.

\section{5. Leanification: A Container-Theorem Case Study And Mention}
Leanification formalizes Theorem B of Marcelo Campos and Wojciech Samotij's Towards an Optimal Hypergraph Container Lemma, together with Proposition 2.2 and Lemmas 4.1-4.3 used in its proof. It does not formalize Theorem A or the whole paper. The informal mathematics belongs to Campos and Samotij. The formalization lead is Luca Seiki Pereira Fujii; Mateus and Tiago helped check the statement according to Luca's final correspondence. Augusto and Felipe Leria are event roster members whom that correspondence explicitly says did not contribute to the submitted work. Team recognition should preserve this distinction instead of automatically turning all registrants into manuscript authors.

The formal theorem concerns a finite hypergraph, with 0<p\ensuremath{\leq}\ensuremath{\delta}<1, and produces a family of fingerprints and associated containers satisfying independent-set coverage, a small fingerprint condition and a lower bound on the probability that a random sampled set remains independent. Independence is defined by absence of a contained hyperedge. The probability is a genuine finite Bernoulli measure using p\textasciicircum{}|A|(1\ensuremath{-}p)\textasciicircum{}|X\textbackslash{}A|, not an arbitrary function whose inequality would be automatic.

The proof includes a nonempty decreasing-family condition in Proposition 2.2, needed to avoid zero conditioning probability and log 0. The assembly supplies that nonemptiness where the lemma is used: the independent fingerprint gives the empty set as a member of the relevant family. Thus the final meaningful cases are not weakened by an unproved extra premise. If the hypergraph contains the empty edge, it has no independent sets, so the corresponding covering clause is legitimately vacuous. A finite linear order selects a deterministic eligible vertex; every finite vertex type admits such an order, so this is a modelling choice rather than a restriction on the combinatorics.

An independent finite semantic check tested 278 hypergraphs on up to 3 vertices, four rational parameter pairs and 1,560 independent-input algorithm runs. All passed covering, fingerprint-size, probability-bound and replay tests. Another 108 decreasing-family checks passed the selected probability inequality. These tests are useful negative-control and interpretation evidence. The proof for arbitrary finite vertex sets still comes from the formal theorem and its exact closure, not from the small cases.

The author's logs were made at commit 816f7e3..., while the preserved archive is adcaf45.... A fresh exact-pin replay or a source diff must close that provenance gap before treating the older build log as independent verification of the final artifact. The source also spells one reviewer 'Matheus', while the final email says 'Mateus'; the preferred full name must be confirmed before publication. Under the chair's latest decision this contribution receives a mention, with the technical substance retained in this paper. Sources: \url{https://doi.org/10.1007/s00493-026-00214-1} ; \url{https://arxiv.org/abs/2408.06617} ; \url{https://github.com/lazyluca/hyprgraph_containers/tree/adcaf45c06fa70407f1abb2df80e1bd801d970b3} .

\section{6. Exact Reconstructions And Certificate Methods}
\subsection{6.1. Collatz descent coverage}
Chaewon Yoon's512-rule packet and HTPeo's234-rule packet cover the same 1,765 odd residue classes at modulus 4096 under the audited rule format. Independent enumeration shows 1,765 is the maximum coverable set in that format, while the classical coefficient stopping-time comparison covers 1,822 odd classes. Jamie's modulus 256, depth-four format has maximum 94 of 128 odd classes, compared with 109 in the classical stopping-time comparison.

These results do not establish a new Collatz theorem. They are exact finite descent certificates reproducing a known mechanism and the pre-event format ceiling. More rules need not mean more mathematical coverage, and better performance within a restrictive checker is not automatically a better theorem than the classical comparison. Their potential value is a transparent certificate language, verified residue arithmetic and a benchmark for search or formal proof generation. A method paper would need to explain compression, checker semantics, negative controls and the relation to existing stopping-time theory, instead of presenting the coverage percentage as a solution of the Collatz conjecture.

\subsection{6.2. Busy Beaver finite runs}
The Sakana six-state machine halts after 1,235,211 steps, with 1,519 ones and span 2,028 in the independently reproduced blank-tape run. HTPeo's machine halts after 249,881 steps, with 554 ones and span 735. These are finite machine certificates, not global Busy Beaver 6 records or a determination of BB6. The current HTPeo packet includes a blank-start counter certificate; an older description alleging only a nonblank initial-state proof is stale.

The reusable interest is in the state-transition model, blank-start correspondence, induction or counter certificate, and exact distinction between steps, tape ones and visited span. A certificate that compresses a million-step computation into a kernel-checked argument can be a useful formal-methods example even when the machine is not mathematically extremal. Publication should reproduce the actual transition table and endpoint convention. Numerical halting alone does not prove structural or global maximality.

\subsection{6.3. Known matrix constructions}
Rank 23, support 139 multiplication tensors appear in several baseline packets. They should be credited to their earlier construction sources. Chandragupt's separate support 143 certificate is exact but does not improve either 139 or the new 138 scheme described in the novel-results volume. Its potential interest is as a differently structured seed and an exact test case for coefficient/gauge transformations, not a new sparsity record.

Exact tensor verification checks 729 coefficient identities for 3\ensuremath{\times}3 matrix multiplication. The search that finds the coefficients can be untrusted if an independent rational certificate proves those identities. Rank, coefficient support, addition count and practical runtime are distinct metrics. A verified rank 23 tensor with more coefficients can still help explore equivalence classes, but unsuccessful searches for support 137 do not prove an absolute lower bound. Baseline reconstruction should therefore be presented as reproducible algebraic infrastructure rather than a further award-generating record.

\subsection{6.4. Kobon baseline arrangements}
The 18-line,93-triangle reconstruction in HTPeo and other accounts reproduces a known lower bound. Several other leaderboard entries match or fall below frozen construction baselines. Their value may lie in a reusable exact arrangement checker, robust rational perturbations, SAT/search tools or reproducible seed data. They do not each become new Kobon theorems because separate accounts found them.

Rohith's39-line,470-triangle result is also subsumed by a pre-event simple straight-line certificate in the chair's public repository. The independent geometric recount confirms 470 under the same ordinary bounded-triangular-face notion. Higher-order catalogue entries need comparison against the full stronger pre-event bound, rather than an older weak formula. Rohith's61-line,1,190-triangle candidate remains a separate novelty question and is reserved for the unresolved-candidate record. Sources: \url{https://github.com/alejandrozu/kobon-proof/commit/99fdc8ec1ef8b1fb22c3da32b011b7361762e958} ; the pinned Rohith packet and companion exact-geometry audit.

\section{7. Correct Certificates With Unresolved Original-Mathematics Status}
\subsection{7.1. Ramsey multiplicity}
Luke Van Seters and Madhan Jothimani's newly recovered packet also preserves the older published 768-vertex adjacency construction. This is a known baseline, separate from their 192-class and 3,840-class original-candidate constructions. An independent data-only recount of the frozen UInt 64 rows verifies a simple symmetric graph with 148,608 red edges. It finds 9,581,824 red triangles and 216,079,680 red four-cliques; the complementary blue graph has 145,920 edges, 9,149,952 triangles and 207,075,456 four-cliques. The balanced blue-diagonal blow-up numerator is therefore 10,487,165,184, giving 4,551,721/150,994,944, approximately 0.0301448570. McKay's recorded numerator 10,486,266,368 with the same denominator is strictly smaller, so this published certificate is not a new numerical improvement.

The ordinary counting identity is transparent. A four-tuple with all entries in the same vertex class contributes n. Each unordered blue edge contributes eight tuples with equality pattern 3+1 and six with pattern 2+2. Each blue triangle contributes 36 tuples with pattern 2+1+1. Finally, each red or blue four-clique contributes 24 tuples with four distinct classes. These exhaust the equality partitions, so the numerator is n+14E\_blue+36T\_blue+24(K4\_red+K4\_blue). The submitted source implements this formula using packed bitsets but explicitly leaves its universal equivalence to a literal tuple sum unproved in Lean. The independent recount uses ordinary arbitrary-precision integer bitsets and separately compares all 64 two-color graphs on four vertices against their literal ordered tuple sums. Those finite checks corroborate the ordinary argument without replacing the missing universal machine proof. The exact frozen Published768.Certificate source has now freshly compiled, including actual execution of its native arithmetic certificate, in 221.671 seconds. The complete repository's 51 formal source files also passed, and its 35 requested axiom outputs include no selected admission or unrecognized axiom; native arithmetic certificates retain the explicitly declared compiler/runtime trust boundary. Four additional standalone research diagnostics separately passed source type-checking; none of their counting main functions was executed, and this is not an extra theorem certificate. The native Count adapter's successful image hashes and three cache slots were observed, but its original WinError 299 observer error and unknown timing remain preserved: continuous observation was incomplete, and root accepted the positive identity/initialization evidence with that limitation. This successful source replay certifies the submitted formula and arithmetic endpoint scope, while the source's omitted universal packed-counter/literal-tuple correspondence remains a separate ordinary-mathematics qualification. This baseline is useful as a regression test for future counters, graph encodings and blow-up interfaces. Sources: \url{https://github.com/lukevs/k4-ramsey/blob/2cf216ba885fbb1b2d1b2d4920efeaa6ba438842/lean/K4Ramsey/Counting/Multiplicity.lean} ; \url{https://github.com/lukevs/k4-ramsey/blob/2cf216ba885fbb1b2d1b2d4920efeaa6ba438842/lean/K4Ramsey/Constructions/Published768/Data.lean} . The dated independent data-count receipt and script are retained in the companion Verification materials.

HTPeo and Sakana supplied exact weighted constructions with K4 multiplicity upper bounds approximately0.030139911990 and 0.030140932308. Both exact counts were independently reproduced. Their finite weighted density and blow-up lifting are mathematical objects worth preserving even when the numerical novelty claim is held.

Technion's January 21,2026 official seminar announcement already gives c(4)<0.030139, numerically stronger than both packets in the same standard density normalization. The announcement is a concrete priority objection, although it is not the entire earlier certificate or proof. The proper next scientific step is to retrieve the dated prior construction and compare the methods and hypotheses. A new method, a distinct formalization or improved reproducibility may remain valuable; a new numerical world-record claim is currently unsupported. Source: \url{https://math.technion.ac.il/events/noam-feinstein/} .

The release should preserve the weighted templates, exact rational numerators and denominators, count certificates and lifting theorem. It should also state explicitly that beating a frozen competition baseline is different from beating all earlier mathematical work. Missing certificates from other leaderboard accounts should remain missing-packet records rather than be assigned a guessed construction or formal status.

\subsection{7.2. Restricted Grothendieck certificates}
Sakana's selected theorem bounds the vector objective Q(A) by sqrt(3/2) times the scalar sign objective C(A) for every 3\ensuremath{\times}n sign matrix, and gives sqrt 2 inequalities for four fixed archived matrices. These statements do not establish a new bound on the unrestricted real Grothendieck constant. The dimensions, sign restriction and four literal finite cases are part of their scope. The independent pinned Lean 4.33.1 replay compiled both exact source modules and produced all five selected endpoint prints with standard kernel axioms or no axioms. The verified statement scope remains restricted, and the mathematical priority hold is unchanged.

The independent audit derived a stronger exact 6/5 bound for the 3\ensuremath{\times}n sign-restricted class, with equality for J\ensuremath{-}2I. This post-event audit argument is not credited to the entrant and is not a claim of first-discovery priority. It shows why the weaker submitted inequality should not automatically receive original-mathematics credit merely because it is correctly formalized. The selected proof can still have value as a reusable norm/certificate interface after formal-library comparison. Known CHSH sqrt 2 witnesses elsewhere in the corpus are likewise established baseline mathematics.

\subsection{7.3. Formalization of weaker or conditional statements}
A correct weak theorem can be a useful engineering milestone, especially when it validates a mathematical encoding. Its publication should show how it composes with stronger known results and why the formal interface is reusable. It should not be made to look stronger by omitting a dimension, hypothesis or exceptional family. Conditional reductions are especially important: they organize future work by exposing a precise missing theorem, but their conclusions remain conditional until that theorem is proved.

Some baseline DMS covering and pullback constructions are known consequences of earlier work. They should remain attributed infrastructure supporting the original-family theorem in the companion volume. The fact that the same reduction is used inside a novel proof does not convert the reduction into another new mathematical result. Conversely, imported known mathematics should not diminish the credit for a new theorem built faithfully on it.

\section{8. Unfinished Research: A Separate Boundary Record}
This section is deliberately separate from proved established mathematics. 'Not accepted as a novel result' does not mean 'proved nonnovel'. Incomplete source, unclosed assumptions, absent certificates and prior overlap are distinct evidence states. The following research directions may be scientifically interesting, but their missing steps cannot be supplied by polished exposition alone.

\subsection{8.1. Matt Bowring's unitary synthesis programme}
Matt's original manuscript proposed a recursive eigenphase construction intended to improve the unitary-synthesis leading coefficient from 22/48 to 21/48. A new 4October intake supplies a stronger formal development at bowrango/lean-unitary commit7014b63ee2cfd01c641a7f539c73b664a2aa3d66. Its generic matrix endpoint proves existence of a pivot g=Rz(alpha)Ry(beta) whose polar-factor multiplexor admits an eigenphase assignment with zero final Walsh angle, assuming invertible blocks in dimensions divisible by 4 and distinct eigenvalue crossing angles. It closes a substantive normalized-eigenphase continuity and intermediate-value argument. The exact endpoint quantifies an unrestricted real angle alpha and an angle beta in [0, pi]. The independent pinned Lean 4.30.0-rc2 replay completed on 5 October: all 26 authored source modules compiled, and all six selected generic-pivot endpoint axiom outputs contain only the standard kernel axioms. A separate inspection of eight Claims endpoints still records sorryAx, consistent with the 13 source placeholders in that baseline chain. The completed generic theorem is consequently distinguished from the unfinished universal circuit-count claim. Fresh proof reproduction does not itself establish mathematical novelty or authorize the provisional award.

The generic theorem belongs in the original partial-results review, not among established nonnovel mathematics. Its endpoint does not contain a circuit or CNOT-count inequality. The Synthesis module assumes the global improved recursion. Non-generic nodes, circuit extraction, Gray-code realization and the sixth-CNOT cancellation bridge remain outside the completed chain; Claims.synth returns the older 22/48 statement through an admitted baseline theorem. The universal 21/48 bound remains held. Formal additions are dated 4October, after the author's pre-cutoff acknowledgment of proof gaps. Intake is accepted for review under the user's late-intake instruction, preserving the proof-completion date. Source: \url{https://github.com/bowrango/lean-unitary/tree/7014b63ee2cfd01c641a7f539c73b664a2aa3d66} .

\subsection{8.2. Geometric and automata packets}
Raj Harshit Srirangam's Kobon reductions lack the verified complete closure needed for original-result acceptance. Frederik Smits van Oyen's Heilbronn covering has an incomplete global cover. Qichao Wang's Černý partial-result manifests arrived, but the actual source packet was absent in the material available to this audit. These limitations should be stated accurately. A metadata manifest cannot substitute for a construction; a local geometric cover cannot establish an unverified global cover; a conditional reduction cannot be called a solved parent problem.

The Cerny correspondence expressly identifies and signs the entrant name Qichao Wang, matching the event registration, while its account display name differs. This review credits the submitted name; the identity mapping, preferred publication name and affiliation require author confirmation. An account display name alone does not establish another entrant or author.

Each has a concrete next scientific step: deliver the exact frozen sources and endpoint statements, isolate the missing assumption or region, and reproduce the relevant checker before determining novelty. Their absence from the novel-results paper is an evidence decision, not an assertion that the authors' proposed ideas are mathematically uninteresting.

\subsection{8.3. SYSTASYS finite inquiry}
Hichem Benali's SYSTASYS inquiry described a finite cyclic problem in Z/25Z. The independent audit confirms maximum size 16 for a subset avoiding nonconstant five-term arithmetic progressions, and exactly 2,000 such size 16 sets containing 0. However, the preserved inquiry lacks a submitted proof packet, novelty comparison and admitted competition target. It is retained as a potentially useful finite benchmark, not an accepted new result and not a solution of Erdős 3. Its coordinatewise-extremality assertion remains separate from the two finite facts independently recounted.

\section{9. Reproducibility And A Practical Library Release}
A useful public formalization release should identify immutable sources, exact Lean and dependency pins, the selected endpoint theorem types and their complete import closures. The build ledger should list what was freshly run, what used external prebuilt dependencies, and what still relies on supplied logs. Expected axioms such as propositional extensionality, classical choice and quotient soundness are distinct from admitted proof holes or a custom unproved conjecture. Each final endpoint should have its own explicit dependency audit.

Libraries should preserve source licenses and authorship headers. Reusable integration does not require redistributing upstream files without permission; an authorized fetch route and pinned source reference can be preferable where licensing is unresolved. A mathematics paper should attribute the original theorem and the formal development separately, while AI/model and compute disclosures explain how the work was produced. Private email addresses and mailbox evidence should not be published merely because they helped resolve internal attribution.

The best collective release structure is thematic: arithmetic and valuations; finite combinatorics and graphs; progression and analytic interfaces; exact search certificates; and a clearly separated unresolved research appendix. It should deduplicate shared statements and retain per-module contributor provenance. Small example certificates can be supplied as test fixtures beside theorems, while substantial complete formalizations deserve explanatory proof architecture. This arrangement lets future researchers reuse the strongest material without inheriting misleading parent-problem claims.

\section{10. Contributor And Affiliation Note}
The companion author inventory records source-backed affiliations and their evidence level. HTPeo's pinned TEAM.md supplies author-declared affiliations for Woohyuk Kang (HTPeo and Kyung Hee University, CS\&E), Hyunjin Lee (University of Cambridge), Sanghyeon Lee (Korea University, Security) and Youchan Oh (Seoul National University, TI). The department abbreviations are preserved rather than expanded speculatively. The Sakana team affiliation is established by its company-owned packet and author correspondence; optional academic dual affiliations require the authors' preferred publication choices.

Chandragupt's supplied public profile identifies the Cluster Innovation Centre, University of Delhi. Rohith's current personal website and Northeastern University's official profile identify the Autonomy and Intelligence Lab at Northeastern. Leon's public profile declares a fall visiting role at Harvard SEAS and study at TUM; Cameron's current public profile lists State Street, with publication affiliation and identity matching still to confirm. Matt's earlier arXiv manuscript v2, dated 9 January 2026, records Purdue University, Department of Mechanical Engineering; his current preferred publication affiliation remains to be confirmed. His personal website supplies corroborating academic context, while separately declared Purdue study and MathWorks employment are not automatically paper dual-affiliations or competition resource classes. Wilson and Chaewon have no verified current institutional affiliation in the supplied materials. Jamie explicitly asserts no organization. Missing affiliations are printed as unverified, not silently invented as 'independent researcher'.

For the container section, Luca is the formalization lead and Mateus/Tiago are statement reviewers; their preferred full names require confirmation. Augusto and Felipe can retain event participation acknowledgment without an invented submitted-work contribution. Prior theorem authors, library authors and external formalization contributors must retain their own attribution even when they were not event entrants. This draft does not assume that every person in a team roster has approved coauthorship of the collective paper. The final author order, contribution statement and affiliations require author review before archive submission or external circulation.

\section{Source Index}
Central preserved source and review cards: \url{https://github.com/alejandrozu/openmath-2026-judging/tree/d0ed487f487fa7ec814ff705ab55249983fe8707/OpenMath-Judging/review} .
Sakana original team repository: \url{https://github.com/SakanaAI/autolab-100} ; immutable delivered ZIP SHA 256 f22e456eecf1a5c071aafa0666fa517dc9ebbd102f483a963f694ee64748f769.
HTPeo: \url{https://github.com/lavaskiller/openmath-2026-htpeo/tree/76c63b8e425e551a930060a6871b84fe3b4780db} .
Chandragupt: \url{https://github.com/ChandraguptSharma07/matrix-multiplication-tensor-3x3/tree/479c2416b6261ee841052eef4881df0e40054f3f} .
Leanification: \url{https://github.com/lazyluca/hyprgraph_containers/tree/adcaf45c06fa70407f1abb2df80e1bd801d970b3} .
Leon/Cameron: \url{https://github.com/Koerbser/erdos-3-autolab/tree/551e388b0a7a933e68cadbab099342af326c1dc2} .
Wilson: \url{https://github.com/wilsonwu-ai/sundai-erdos-3/tree/492fd601f9c76e44dcb87f683a5ea6ccb7aaefc9} .
Jamie: \url{https://github.com/grandchallenge/MATHSOLVE/tree/2fd9fd643bf3e9e98d121dc2e7af1cf1be411ae2} .
Matt: \url{https://github.com/bowrango/lean-unitary/tree/57d6c96992a58f43d48d61c2a44353d902813cda/submission} .
Campos-Samotij mathematical source: \url{https://doi.org/10.1007/s00493-026-00214-1} ; \url{https://arxiv.org/abs/2408.06617} .
Ramsey prior announcement: \url{https://math.technion.ac.il/events/noam-feinstein/} .
Kobon pre-event 470: \url{https://github.com/alejandrozu/kobon-proof/commit/99fdc8ec1ef8b1fb22c3da32b011b7361762e958} .

APPENDIX A. EXACT MACHINE TABLES AND DEFINITIONS

Each deterministic six-state, two-symbol Turing-machine entry wDN writes bitw, moves D=L/R, and enters N; H halts. Execution starts in stateA, head 0, on the all-zero bi-infinite tape. Step count includes the transition enteringH. Visited span and final 1 count are separate quantities.

State | Sakana read 0 | Sakana read 1 | HTPeo read 0 | HTPeo read 1
A | 1RB | 0LC | 0LE | 1LA
B | 1RD | 1RB | 1LC | 1RH
C | 0LH | 1LE | 1RA | 1RF
D | 1LA | 1LD | 0RF | 0RD
E | 0LA | 1LF | 1RD | 1LB
F | 0LE | 0RD | 1LC | 1RD

Reviewer-owned blank-tape simulation reproduces 1,235,211 steps,1,519 ones and span 2,028 for Sakana, and 249,881 steps,554 ones and span 735 for HTPeo. These are finite trajectories, not maximality theorems. The Sakana literal machine SHA-256 is8cd4b6a0ea5ac0333c12cbd2b3f4b1f04c7dbb970b360b4e2086efbb5cf1d726. HTPeo source: \url{https://github.com/lavaskiller/openmath-2026-htpeo/blob/76c63b8e425e551a930060a6871b84fe3b4780db/entries/hills/busy-beaver-6-n0rang2/solution.json} . Sakana's submitted FOCUS-BB6 packet and independent simulator are preserved in the central archive.

IsDComplete(A) is formally eventual existence, for each natural n, of finite s contained in A, IsAntichain(divides,s), and sum(s)=n. E292.A is\{n:exists finite S, S subset[1,n], n in S, sum\_\{m in S\}(1:Q)/m=1\}. nonnovel\_math\_appendix.json preserves the exact Lean predicates, source URLs, hashes and both machine tables. Definition sources: \url{https://github.com/lavaskiller/openmath-2026-htpeo/blob/76c63b8e425e551a930060a6871b84fe3b4780db/entries/erdos-m2-formalizations/artifact/bundle/Erdos123.lean} ; \url{https://github.com/lavaskiller/openmath-2026-htpeo/blob/76c63b8e425e551a930060a6871b84fe3b4780db/entries/erdos-m2-formalizations/artifact/bundle/Erdos292.lean} .

Dated Matt affiliation sources: \url{https://mattbowring.com/} and \url{https://arxiv.org/html/2512.23720v2} . These sources concern author attribution; they do not certify the new generic balanced polar-pivot theorem or the global improved CNOT bound.

\clearpage
\section{Author and affiliation register}
Affiliations are source-backed declarations, not institutional endorsement. Confirm preferred name, publication affiliation and authorship order before external release. Do not infer paper authorship from event roster or organizer role.

\subsection{Rachel Teo (E02)}
Rachel S. Y. Teo on personal academic website; event uses Rachel Teo

Recorded role: Joint Sakana packet; do not invent per-result split

Sakana AI. Evidence: Event team affiliation; author correspondence and Sakana-owned team repository.

{\small \url{https://github.com/SakanaAI/autolab-100}\par}

National University of Singapore. Evidence: Author-declared on personal academic website fetched October 4; year-of-study text may be stale.

{\small \url{https://rachtsy.github.io/}\par}

\subsection{Wenyi Wang (E02)}
Recorded role: Joint Sakana packet; do not invent per-result split

Sakana AI. Evidence: Author public profile and team packet.

{\small \url{https://sa.linkedin.com/in/wenyiw}\par}

{\small \url{https://github.com/SakanaAI/autolab-100}\par}

King Abdullah University of Science and Technology (KAUST). Evidence: Author LinkedIn education lists 2022-present; preferred paper dual affiliation must be confirmed.

{\small \url{https://sa.linkedin.com/in/wenyiw}\par}

\subsection{Hamdi Abderrahmene (E02)}
Public LinkedIn uses Abderrahmene Hamdi; confirm desired order

Recorded role: Joint Sakana packet; do not invent per-result split

Sakana AI. Evidence: Author LinkedIn Research Intern August 2026-present, corroborated team repository.

{\small \url{https://jp.linkedin.com/in/abderrahmenehamdi}\par}

{\small \url{https://github.com/SakanaAI/autolab-100}\par}

Nagoya University. Evidence: Author LinkedIn master's programme 2026-2028; preferred paper dual affiliation must be confirmed.

{\small \url{https://jp.linkedin.com/in/abderrahmenehamdi}\par}

\subsection{Woohyuk Kang (E08)}
Recorded role: DMS, selected Erdős formalizations, leading Ramsey search and certificate, Kobon 39 packet; team assembly

HTPeo; Kyung Hee University, CS\&E. Evidence: Personally entered author declaration in pinned TEAM.md and manuscript; preserve department abbreviation.

{\small \url{https://github.com/lavaskiller/openmath-2026-htpeo/blob/76c63b8e425e551a930060a6871b84fe3b4780db/TEAM.md}\par}

{\small \url{https://github.com/lavaskiller/openmath-2026-htpeo/blob/76c63b8e425e551a930060a6871b84fe3b4780db/paper/htpeo-openmath-2026.tex}\par}

\subsection{Hyunjin Lee (E08)}
Recorded role: Separate Ramsey final/validation entry; Erdős 1038 with Sanghyeon, not claimed; public-seed matrix search

University of Cambridge. Evidence: Personally entered author declaration in pinned TEAM.md.

{\small \url{https://github.com/lavaskiller/openmath-2026-htpeo/blob/76c63b8e425e551a930060a6871b84fe3b4780db/TEAM.md}\par}

\subsection{Sanghyeon Lee (E08)}
Recorded role: Busy Beaver 249881 certificate; Ramsey leaderboard record; Erdős 1038 with Hyunjin, not claimed

Korea University, Security. Evidence: Author declaration recorded in pinned TEAM.md; preserve department wording.

{\small \url{https://github.com/lavaskiller/openmath-2026-htpeo/blob/76c63b8e425e551a930060a6871b84fe3b4780db/TEAM.md}\par}

\subsection{Youchan Oh (E08)}
Recorded role: Known Kobon 18=93 reconstruction

Seoul National University, TI. Evidence: Author declaration recorded in pinned TEAM.md; do not guess TI expansion.

{\small \url{https://github.com/lavaskiller/openmath-2026-htpeo/blob/76c63b8e425e551a930060a6871b84fe3b4780db/TEAM.md}\par}

\subsection{Chandragupt Sharma (E04)}
Recorded role: Sole named matrix support 138 packet author

Cluster Innovation Centre, University of Delhi. Evidence: Author public LinkedIn profile; identity linked by supplied registration profile; confirm preferred paper affiliation.

{\small \url{https://in.linkedin.com/in/chandraguptsharma}\par}

LinkedIn also lists Mathematical Sciences Foundation experience; role/date insufficient to assert that as paper institution.

\subsection{Rohith Poola (E57)}
Recorded role: Sole named Kobon catalogue author

Autonomy and Intelligence Lab, Northeastern University. Evidence: Official university March 9 2026 profile and current author website.

{\small \url{https://coe.northeastern.edu/news/driven-by-curiosity-rohith-poolas-journey-in-autonomous-robotics/}\par}

{\small \url{https://Rohith18p.github.io/}\par}

\subsection{Leon Koerbs (E11)}
GitHub and LinkedIn use Leon Körbs; use author-confirmed spelling

Recorded role: Joint Erdős 3 Token Burner packet with Cameron Fen; per-result roles not divided

Harvard John A. Paulson School of Engineering and Applied Sciences, Intelligent Interactive Systems Group. Evidence: Author public LinkedIn says Visiting Scholar for fall term; not independent institutional appointment verification.

{\small \url{https://de.linkedin.com/in/leon-koerbs}\par}

Technical University of Munich, School of Management. Evidence: Current author LinkedIn study 2025-2027; earlier official exchange report corroborates identity.

{\small \url{https://de.linkedin.com/in/leon-koerbs}\par}

{\small \url{https://cms.mgt.tum.de/fileadmin/mgt.tum.de/downloads/international_exchange_programs/experience_reports/Erfahrungsbericht_Leon_Koerbs_SMU_SS_23_DE_TUMSOMex_min.pdf}\par}

\subsection{Cameron Fen (E11)}
Recorded role: Joint Erdős 3 Token Burner packet with Leon

State Street. Evidence: Current author public LinkedIn identifies State Street; identity-to-event roster and permitted publication affiliation require confirmation.

{\small \url{https://www.linkedin.com/in/cameron-fen-phd-945237225}\par}

{\small \url{https://cameronfen.github.io/}\par}

University of Michigan is documented education, not asserted current affiliation.

\subsection{Wilson Wu (E65)}
Recorded role: Sole named supporting-formalizations packet author

Affiliation not verified in the supplied materials. Evidence: GitHub and personal page identify author but no institution/employer.

{\small \url{https://github.com/wilsonwu-ai}\par}

{\small \url{https://wilsonwu-ai.github.io/}\par}

\subsection{Jamie Steeg (E10)}
Recorded role: Sole named H1/H4/H7 entrant; GCL/GCT project provenance not paper affiliation

No organizational affiliation asserted. Evidence: Explicit human authority declaration in frozen packet; individual and self-resourced.

{\small \url{https://github.com/grandchallenge/MATHSOLVE/blob/2fd9fd643bf3e9e98d121dc2e7af1cf1be411ae2/work_packages/OPENMATH_2026/COMPETITION_PACKETS/HUMAN_AUTHORITY_DECLARATION.json}\par}

\subsection{Chaewon Yoon (E05)}
Recorded role: Collatz 512-rule packet

Affiliation not verified in the supplied materials. Evidence: Named private submission email and artifact establish author, not institution.

\subsection{Matt Bowring (E03)}
Recorded role: Generic matrix eigenphase-pivot partial theorem received 4 October; actual exact-pin replay completed 5 October with all 26 authored source modules compiling and all six selected generic endpoint axiom outputs standard-kernel only. Universal improved CNOT synthesis claim remains unproved; separate Claims diagnostics retain sorryAx.

Purdue University, Department of Mechanical Engineering (documented January 2026; current preferred publication affiliation to confirm). Evidence: Author-declared institutional affiliation in arXiv manuscript v2 dated 9 January 2026 and corroborated by personal website; this dated paper affiliation does not establish a current appointment or competition resource class.

{\small \url{https://mattbowring.com/}\par}

{\small \url{https://arxiv.org/html/2512.23720v2}\par}

{\small \url{https://www.linkedin.com/in/matt-bowring-463610149}\par}

\subsection{Luke Van Seters (E01)}
Recorded role: Joint Ramsey packet with Madhan Jothimani; preserved manuscript provides author names, no affiliations

Klaviyo. Evidence: Authenticated public LinkedIn profile header reviewed 4 October 2026 at approximately 22:58 UTC states Director of Engineering at Klaviyo. This is a dated profile declaration; preferred publication affiliation awaits author confirmation.

{\small \url{https://www.linkedin.com/in/lukevanseters/}\par}

{\small \url{https://github.com/lukevs/k4-ramsey/blob/2cf216ba885fbb1b2d1b2d4920efeaa6ba438842/README.md}\par}

Earlier Agency AI attribution remains unconfirmed historical information and is superseded as a current-profile description. Northeastern University appears as profile education, not a confirmed current university affiliation. Declared competition resources determine the competition class separately from employment or education.

\subsection{Madhan Jothimani (E01)}
Recorded role: Joint Ramsey packet with Luke Van Seters

Northeastern University. Evidence: Author public LinkedIn lists education from 2024 to the present; exact programme and preferred paper affiliation require author confirmation.

{\small \url{https://www.linkedin.com/in/madhanjothimani}\par}

{\small \url{https://github.com/lukevs/k4-ramsey/blob/2cf216ba885fbb1b2d1b2d4920efeaa6ba438842/README.md}\par}

Profile lists Infosys experience without sufficiently clear current dates; not silently asserted as publication employer.

\subsection{Luca Seiki Pereira Fujii (E06)}
Recorded role: Container formalization lead, scope/design and statement review; mention-only administrative decision

IMPA Tech. Evidence: Author explicitly identifies five team members as IMPA Tech-account registrants in focused email; do not rely solely on email domain.

\subsection{Mateus (E06)}
Recorded role: Statement reviewer; do not invent full author name

IMPA Tech. Evidence: Author team declaration; full preferred name unresolved; Luma says Mateus Mundstock, source CREDITS says Matheus.

{\small \url{https://github.com/lazyluca/hyprgraph_containers/blob/adcaf45c06fa70407f1abb2df80e1bd801d970b3/CREDITS.md}\par}

\subsection{Tiago (E06)}
Recorded role: Statement reviewer

IMPA Tech. Evidence: Author team declaration; registration links Tiago Sánchez Ribeiro but exact packet-name expansion awaits confirmation.

{\small \url{https://github.com/lazyluca/hyprgraph_containers/blob/adcaf45c06fa70407f1abb2df80e1bd801d970b3/CREDITS.md}\par}

\subsection{Augusto (E06)}
Recorded role: Event roster member; Luca explicitly says did not contribute to submitted work; acknowledge participation, not automatic paper authorship

IMPA Tech. Evidence: Author team declaration; possible full name Augusto Schaefer Huff from registration, not silently merged.

\subsection{Felipe Leria (E06)}
Recorded role: Event roster member; Luca explicitly says did not contribute to submitted work; acknowledge participation, not automatic paper authorship

IMPA Tech. Evidence: Author team declaration; preferred complete name not expanded from email address.

\subsection{Frederik Smits van Oyen (E12)}
Recorded role: Incomplete Heilbronn covering and a separate Ramsey project without an author-designated final certificate

Affiliation not verified in the supplied materials. Evidence: Named author in frozen TEAM record; neither name similarity nor a GitHub account establishes an institution.

{\small \url{https://github.com/alejandrozu/openmath-2026-judging/blob/d0ed487f487fa7ec814ff705ab55249983fe8707/OpenMath-Judging/contestants/E12/TEAM.json}\par}

{\small \url{https://github.com/FSvOAI/Heilbronn/tree/b8da072acb21d2b0acc8daa038600dddfc8c67e9}\par}

This is the E12 author; do not substitute Frederik Baetens.

\subsection{Raj Harshit Srirangam (E09)}
Recorded role: Kobon upper-bound reductions and exact certificates with unresolved fourfold-point and global closure conditions

Affiliation not verified in the supplied materials. Evidence: Named author in pinned contestant TEAM record and repository; no confirmed present institution.

{\small \url{https://github.com/alejandrozu/openmath-2026-judging/blob/d0ed487f487fa7ec814ff705ab55249983fe8707/OpenMath-Judging/contestants/E09/TEAM.json}\par}

{\small \url{https://github.com/srirangam-r/kobon_triangles/tree/eed14659a9b2f4ca12777da5d557f2b620b966f6}\par}

\subsection{Qichao Wang (E17)}
Recorded role: Cerny partial-formalization manifests; actual source ZIP remains unavailable

Hebei University of Technology. Evidence: Author-declared undergraduate Automation study in event registration; publication affiliation and exact programme styling require author confirmation.

{\small \url{https://github.com/alejandrozu/openmath-2026-judging/blob/d0ed487f487fa7ec814ff705ab55249983fe8707/OpenMath-Judging/contestants/EVENT/12-correspondence/3937a54b1ac3-OpenMath-Luma-roster.json}\par}

{\small \url{https://www.linkedin.com/in/%E5%90%AF%E8%B6%85-%E7%8E%8B-a5a69043b/}\par}

Credit uses the registration and expressly submitted name Qichao Wang: the final correspondence identifies the entrant and signs that name. The account display name differs (Cui Lijun); its identity mapping and the preferred publication name are not independently verified and require author confirmation. No second entrant or coauthor is inferred.

\subsection{Hichem Benali (Unassigned correspondence)}
Recorded role: Finite cyclic five-AP inquiry in Z/25Z; author source proof packet and official target/admission remain unresolved

Affiliation not verified in the supplied materials. Evidence: SYSTASYS is a correspondence/project label, not independently verified as an employer or institution.

Private correspondence identifies the author. No outbound correspondence or institutional inference was made.

\clearpage
\section{Exhaustive source-status appendix}
This appendix preserves 119 lower-level formalization, baseline, unfinished and administrative records. These are evidence states, not 119 accepted theorems. Entries related to the original-results volume are cross-references rather than duplicate scientific claims. The source snapshot is October 3; the fresh verification ledger and October 4 decisions supersede older status language. Luke/Madhan's newly recovered packet is treated in the original-results review.

\begin{samepage}
\subsection{E03-SYNTHESIS  -  Improved synthesis bound}
{\small Contributors: Matt Bowring | Family: UNITARY\par}

\textbf{Recorded scope:} Claimed (21/48)4\textasciicircum{}n \ensuremath{-} (3/2)2\textasciicircum{}n + 2 CNOTs; n=3:18 CX; n=4:90 CX.

\end{samepage}
\textbf{Source classification:} Original generic balanced polar-pivot partial result; universal 21/48 circuit-count claim remains unproved and belongs in the unresolved appendix.

\textbf{Prior comparison:} Potential new coefficient if proved. Compare to Krol-Al-Ars 2024, not only the older 23/48 baseline.

\textbf{Formal/source status:} New exact source 7014b63ee2cfd01c641a7f539c73b664a2aa3d66, Lean 4.30.0-rc2: actual independent replay completed 2026-10-05T01:38:41Z; all 26 source modules PASS and 6/6 selected generic-pivot axiom outputs standard-kernel only. Separate 8/8 Claims audit outputs include sorryAx, consistent with 13 baseline source placeholders. The generic endpoint takes invertible X/Y, dimensions 4(q+1), separated crossings, unrestricted real alpha and beta in[0,pi]; no global constructive circuit/CNOT bridge is proved.

\textbf{Research or reuse implication:} Accepted competitive score remains 0 pending chair/novelty/resource-class review. Recommended conditional P2=.05 yields 33.9488 from D=824 under approved scoring; this is for the completed narrow generic partial result. Universal 21/48 claim scores 0; hypotheticalP 3=.15=101.8464 is not adopted.

{\small \url{https://github.com/alejandrozu/openmath-2026-judging/blob/d0ed487f487fa7ec814ff705ab55249983fe8707/OpenMath-Judging/review/entries/E03-SYNTHESIS.txt}\par}

\begin{samepage}
\subsection{E04-GAUGE  -  Gauge and inequivalence structure}
{\small Contributors: Chandragupt Sharma | Family: MATRIX-3\par}

\textbf{Recorded scope:} Gauged scheme:126 entries \ensuremath{\pm}1,12 entries \ensuremath{\pm}2 or \ensuremath{\pm}1/2; 83 additions plus 12 shifts, CSE 77 versus cited Smirnov 70; one-parameter family and rank-multiset inequivalence.

\end{samepage}
\textbf{Source classification:} Baseline/known-result candidate, or incomplete claim with prior overlap; see exact evidence

\textbf{Prior comparison:} Potentially substantive structural analysis of the new construction. The family identity is now independently checked algebraically, but its general parameter theorem is not included in the supplied Lean certificate; no separate formal-only award.

\textbf{Formal/source status:} The dated October 3 archive said no fresh kernel build; that historical state is superseded for the exact finite certificates. The current SUPPORT 138 receipt at source commit 479c2416b6261ee841052eef4881df0e40054f3f freshly compiles all seven archived Lean modules under Lean 4.34.1 and prints all twelve finite validity, identity and support endpoints with no selected admission or custom axiom. This includes the support-138 scheme, its support-138 gauged variant, the known support-139 auxiliary scheme and the separate support-143 scheme. These endpoints use ordinary kernel computation, rather than native\_decide. The exact independent checker separately confirms all 729 tensor identities and rational/source correspondence. The universal rational parameter-family identity outside s=0,-2 is independently checked by ordinary exact algebra, but no universal parameter-family Lean theorem occurs among these seven sources or twelve selected endpoints. README scale-factor entries remain historically stale; actual source/data correspondence passes. See the current SUPPORT 138 chapter and Verification/chandragupt-pinned-build.json for exact fresh scope.

\textbf{Research or reuse implication:} Supporting gauge and parameter structure belongs to the same SUPPORT 138 construction family. Supporting gauge score remains 0 and there are no additional family points. Exact finite gauged validity is included in the completed certificate; universal parameter, global inequivalence and optimality claims retain their separate ordinary/formal qualifications.

{\small \url{https://github.com/alejandrozu/openmath-2026-judging/blob/d0ed487f487fa7ec814ff705ab55249983fe8707/OpenMath-Judging/review/entries/E04-GAUGE.txt}\par}

\begin{samepage}
\subsection{E04-CORE 143  -  Separate support-143 construction}
{\small Contributors: Chandragupt Sharma | Family: MATRIX-3\par}

\textbf{Recorded scope:} Core-1809 rank 23/support 143; empirical unsuccessful search for 137.

\end{samepage}
\textbf{Source classification:} Baseline/known-result candidate, or incomplete claim with prior overlap; see exact evidence

\textbf{Prior comparison:} No eligible original mathematical delta established in the available packet. This is not a finding that all underlying work lacks novelty.

\textbf{Formal/source status:} The current SUPPORT 138 replay at frozen commit 479c2416b6261ee841052eef4881df0e40054f3f also verifies this separate support-143 scheme: its finite identity and support endpoints are included in the seven-module, twelve-output ordinary-kernel PASS detailed in E04-GAUGE above. The dated October 3 no-build wording is retained as historical evidence and is superseded for these finite certificates. The independent exact checker confirms rational/source correspondence; README scale entries remain historically stale. Universal rational parameter-family identities remain ordinary exact algebra, with no universal parameter Lean theorem in this source set. See Verification/chandragupt-pinned-build.json. This supporting scheme earns no additional family credit.

\textbf{Research or reuse implication:} Propose this exact scope and family once, subject to the stated source, formal, novelty and deadline checks.

{\small \url{https://github.com/alejandrozu/openmath-2026-judging/blob/d0ed487f487fa7ec814ff705ab55249983fe8707/OpenMath-Judging/review/entries/E04-CORE143.txt}\par}

\begin{samepage}
\subsection{E06-THEOREM-B  -  Known container lemma formalization}
{\small Contributors: Luca Seiki Pereira Fujii, Mateus, Tiago | Family: M2-CONTAINERS\par}

\textbf{Recorded scope:} TheoremB\_unconditional plus Proposition 2.2 and supporting Lemmas 4.1-4.3.

\end{samepage}
\textbf{Source classification:} Known-mathematics formalization / formal-library novelty and usefulness require verification

\textbf{Prior comparison:} Known Campos-Samotij mathematics; candidate new useful formalization only.

\textbf{Formal/source status:} 13 local Lean sources inspected. Submitted logs list only standard axioms. Nonempty-family qualification avoids undefined conditional expectation/log 0; LinearOrder on finite V can be chosen existentially. No fresh kernel build or two independent referee reviews.

\textbf{Research or reuse implication:} Mention the container formalization and its reusable library value. The user's explicit October 4 ruling excludes Leanification from competitive M2 placement and points; the earlier .8 adjustment is superseded.

{\small \url{https://github.com/alejandrozu/openmath-2026-judging/blob/d0ed487f487fa7ec814ff705ab55249983fe8707/OpenMath-Judging/review/entries/E06-THEOREM-B.txt}\par}

\begin{samepage}
\subsection{E07-RAMTASTIC  -  Pre-event Ramsey exploration}
{\small Contributors:  | Family: RAMSEY-K4\par}

\textbf{Recorded scope:} Public weighted search repository and pre-event template work; no designated event proof/result.

\end{samepage}
\textbf{Source classification:} Unresolved or missing-artifact record; nonacceptance does not prove nonnovelty

\textbf{Prior comparison:} No eligible original mathematical delta established in the available packet. This is not a finding that all underlying work lacks novelty.

\textbf{Formal/source status:} Repository is research evidence, not a final certified claim.

\textbf{Research or reuse implication:} Hold missing payload; current proposed points 0. Board metadata cannot establish an eligible exact theorem.

{\small \url{https://github.com/alejandrozu/openmath-2026-judging/blob/d0ed487f487fa7ec814ff705ab55249983fe8707/OpenMath-Judging/review/entries/E07-RAMTASTIC.txt}\par}

\begin{samepage}
\subsection{E08-G-PM  -  Petersen / Schönberger perfect matchings}
{\small Contributors: Woohyuk Kang, Hyunjin Lee, Sanghyeon Lee, Youchan Oh | Family: M2-G-PM\par}

\textbf{Recorded scope:} def Star6.Bridgeless (G : SimpleGraph V) : Prop := \ensuremath{\forall} e \ensuremath{\in} G.edgeSet, \ensuremath{\neg} G.IsBridge e
theorem Star6.simple\_schoenberger (hconn : G.Connected) (hreg : G.IsRegularOfDegree 3) (hbr : Bridgeless G)
    \{e : Sym2 V\} (he : e \ensuremath{\in} G.edgeSet) :
    (\ensuremath{\exists} M : G.Subgraph, M.IsPerfectMatching \ensuremath{\wedge} e \ensuremath{\in} M.edgeSet) \ensuremath{\wedge} (\ensuremath{\exists} M : G.Subgraph, M.IsPerfectMatching \ensuremath{\wedge} e \ensuremath{\notin} M.edgeSet)
theorem Star6.simple\_petersen\_connected (hconn : G.Connected) (hreg : G.IsRegularOfDegree 3) (hbr : Bridgeless G) :
    \ensuremath{\exists} M : G.Subgraph, M.IsPerfectMatching


\end{samepage}
\textbf{Source classification:} Known-mathematics formalization / formal-library novelty and usefulness require verification

\textbf{Prior comparison:} Known mathematics; only a genuinely new, faithful, useful, reusable event-window formalization may count. Exact prior-library overlap is not yet independently cleared.

\textbf{Formal/source status:} Packet reports pinned builds and standard axioms; source correspondence and full novelty remain review tasks. Some same-file declarations are explicitly unclaimed duplicates.

\textbf{Research or reuse implication:} Propose this exact scope and family once, subject to the stated source, formal, novelty and deadline checks.

{\small \url{https://github.com/alejandrozu/openmath-2026-judging/blob/d0ed487f487fa7ec814ff705ab55249983fe8707/OpenMath-Judging/review/entries/E08-G-PM.txt}\par}

\begin{samepage}
\subsection{E08-E942  -  Unbounded powerful-number counts between consecutive squares}
{\small Contributors: Woohyuk Kang, Hyunjin Lee, Sanghyeon Lee, Youchan Oh | Family: M2-E942\par}

\textbf{Recorded scope:} @[category textbook, AMS 11] theorem erdos\_942.variants.limsup : atTop.limsup (((fun (n : \ensuremath{\mathbb{N}}) \ensuremath{\mapsto} (n : \ensuremath{\mathbb{N}}\ensuremath{\infty})) \ensuremath{\circ} erdos\_942.h)) = \ensuremath{\top}


\end{samepage}
\textbf{Source classification:} Known-mathematics formalization / formal-library novelty and usefulness require verification

\textbf{Prior comparison:} Known mathematics; only a genuinely new, faithful, useful, reusable event-window formalization may count. Exact prior-library overlap is not yet independently cleared.

\textbf{Formal/source status:} Packet reports pinned builds and standard axioms; source correspondence and full novelty remain review tasks. Some same-file declarations are explicitly unclaimed duplicates.

\textbf{Research or reuse implication:} Propose this exact scope and family once, subject to the stated source, formal, novelty and deadline checks.

{\small \url{https://github.com/alejandrozu/openmath-2026-judging/blob/d0ed487f487fa7ec814ff705ab55249983fe8707/OpenMath-Judging/review/entries/E08-E942.txt}\par}

\begin{samepage}
\subsection{E08-E44  -  Greedy Sidon lower bound N \ensuremath{\leq} |A|\ensuremath{^3}}
{\small Contributors: Woohyuk Kang, Hyunjin Lee, Sanghyeon Lee, Youchan Oh | Family: M2-E44\par}

\textbf{Recorded scope:} @[category textbook, AMS 5 11] theorem greedy\_sidon\_construction (N : \ensuremath{\mathbb{N}}) (hN : 1 \ensuremath{\leq} N) : \ensuremath{\exists}\ensuremath{^{e}} (A \ensuremath{\subseteq} Finset.Icc 1 N), IsSidon (A : Set \ensuremath{\mathbb{N}}) \ensuremath{\wedge} N \ensuremath{\leq} A.card \textasciicircum{} 3


\end{samepage}
\textbf{Source classification:} Known-mathematics formalization / formal-library novelty and usefulness require verification

\textbf{Prior comparison:} Known mathematics; only a genuinely new, faithful, useful, reusable event-window formalization may count. Exact prior-library overlap is not yet independently cleared.

\textbf{Formal/source status:} Packet reports pinned builds and standard axioms; source correspondence and full novelty remain review tasks. Some same-file declarations are explicitly unclaimed duplicates.

\textbf{Research or reuse implication:} Propose this exact scope and family once, subject to the stated source, formal, novelty and deadline checks.

{\small \url{https://github.com/alejandrozu/openmath-2026-judging/blob/d0ed487f487fa7ec814ff705ab55249983fe8707/OpenMath-Judging/review/entries/E08-E44.txt}\par}

\begin{samepage}
\subsection{E08-E123  -  d-completeness of powers of 2 and 3}
{\small Contributors: Woohyuk Kang, Hyunjin Lee, Sanghyeon Lee, Youchan Oh | Family: M2-E123\par}

\textbf{Recorded scope:} @[category research solved, AMS 11] theorem erdos\_123.variants.powers\_2\_3 : IsDComplete (↑(powers 2) * ↑(powers 3))


\end{samepage}
\textbf{Source classification:} Known-mathematics formalization / formal-library novelty and usefulness require verification

\textbf{Prior comparison:} Known mathematics; only a genuinely new, faithful, useful, reusable event-window formalization may count. Exact prior-library overlap is not yet independently cleared.

\textbf{Formal/source status:} Packet reports pinned builds and standard axioms; source correspondence and full novelty remain review tasks. Some same-file declarations are explicitly unclaimed duplicates.

\textbf{Research or reuse implication:} Propose this exact scope and family once, subject to the stated source, formal, novelty and deadline checks.

{\small \url{https://github.com/alejandrozu/openmath-2026-judging/blob/d0ed487f487fa7ec814ff705ab55249983fe8707/OpenMath-Judging/review/entries/E08-E123.txt}\par}

\begin{samepage}
\subsection{E08-E918  -  Impossible exact countable chromatic restrictions}
{\small Contributors: Woohyuk Kang, Hyunjin Lee, Sanghyeon Lee, Youchan Oh | Family: M2-E918\par}

\textbf{Recorded scope:} @[category textbook, AMS 5] theorem erdos\_918.variants.eq\_aleph\_0.parts.i : \ensuremath{\neg}\ensuremath{\exists} (V : Type u) (G : SimpleGraph V), \#V = \ensuremath{\aleph}\_ 2 \ensuremath{\wedge} G.chromaticCardinal = \ensuremath{\aleph}\_ 2 \ensuremath{\wedge} \ensuremath{\forall} (W : Set V) (\_ : \#W = \ensuremath{\aleph}\ensuremath{_1}), (G.induce W).chromaticCardinal = \ensuremath{\aleph}\ensuremath{_0}
@[category textbook, AMS 5] theorem erdos\_918.variants.eq\_aleph\_0\_all\_subgraphs.parts.i : \ensuremath{\neg}\ensuremath{\exists} (V : Type u) (G : SimpleGraph V), \#V = \ensuremath{\aleph}\_ 2 \ensuremath{\wedge} G.chromaticCardinal = \ensuremath{\aleph}\_ 2 \ensuremath{\wedge} \ensuremath{\forall} (H : G.Subgraph) (\_ : \#H.verts = \ensuremath{\aleph}\ensuremath{_1}), H.coe.chromaticCardinal = \ensuremath{\aleph}\ensuremath{_0}
@[category textbook, AMS 5] theorem erdos\_918.variants.eq\_aleph\_0\_all\_subgraphs.parts.ii : \ensuremath{\neg}\ensuremath{\exists} (V : Type u) (G : SimpleGraph V), \#V = \ensuremath{\aleph}\_ (\ensuremath{\omega} + 1) \ensuremath{\wedge} G.chromaticCardinal = \ensuremath{\aleph}\ensuremath{_1} \ensuremath{\wedge} \ensuremath{\forall} (H : G.Subgraph) (\_ : \#H.verts = \ensuremath{\aleph}\_ \ensuremath{\omega}), H.coe.chromaticCardinal = \ensuremath{\aleph}\ensuremath{_0}


\end{samepage}
\textbf{Source classification:} Known-mathematics formalization / formal-library novelty and usefulness require verification

\textbf{Prior comparison:} Known mathematics; only a genuinely new, faithful, useful, reusable event-window formalization may count. Exact prior-library overlap is not yet independently cleared.

\textbf{Formal/source status:} Packet reports pinned builds and standard axioms; source correspondence and full novelty remain review tasks. Some same-file declarations are explicitly unclaimed duplicates.

\textbf{Research or reuse implication:} Propose this exact scope and family once, subject to the stated source, formal, novelty and deadline checks.

{\small \url{https://github.com/alejandrozu/openmath-2026-judging/blob/d0ed487f487fa7ec814ff705ab55249983fe8707/OpenMath-Judging/review/entries/E08-E918.txt}\par}

\begin{samepage}
\subsection{E08-E292  -  Multiplicative closure / doubling / no prime powers}
{\small Contributors: Woohyuk Kang, Hyunjin Lee, Sanghyeon Lee, Youchan Oh | Family: M2-E292\par}

\textbf{Recorded scope:} @[category research solved, AMS 11] theorem erdos\_292.variants.mul : \ensuremath{\forall} m \ensuremath{\in} A, \ensuremath{\forall} n \ensuremath{\in} A, m * n \ensuremath{\in} A
@[category research solved, AMS 11] theorem erdos\_292.variants.two\_mul : \ensuremath{\forall} n \ensuremath{\in} A, 1 < n \ensuremath{\to} 2 * n \ensuremath{\in} A
@[category research solved, AMS 11] theorem erdos\_292.variants.prime\_pow : \ensuremath{\forall} n \ensuremath{\in} A, \ensuremath{\neg} IsPrimePow n


\end{samepage}
\textbf{Source classification:} Known-mathematics formalization / formal-library novelty and usefulness require verification

\textbf{Prior comparison:} Known mathematics; only a genuinely new, faithful, useful, reusable event-window formalization may count. Exact prior-library overlap is not yet independently cleared.

\textbf{Formal/source status:} Packet reports pinned builds and standard axioms; source correspondence and full novelty remain review tasks. Some same-file declarations are explicitly unclaimed duplicates.

\textbf{Research or reuse implication:} Propose this exact scope and family once, subject to the stated source, formal, novelty and deadline checks.

{\small \url{https://github.com/alejandrozu/openmath-2026-judging/blob/d0ed487f487fa7ec814ff705ab55249983fe8707/OpenMath-Judging/review/entries/E08-E292.txt}\par}

\begin{samepage}
\subsection{E08-E395  -  Radius-one reverse Littlewood-Offord counterexample}
{\small Contributors: Woohyuk Kang, Hyunjin Lee, Sanghyeon Lee, Youchan Oh | Family: M2-E395\par}

\textbf{Recorded scope:} @[category research solved, AMS 5 60] theorem erdos\_395.variants.one : answer(False) \ensuremath{\leftrightarrow} \ensuremath{\exists} c : \ensuremath{\mathbb{R}}, 0 < c \ensuremath{\wedge} \ensuremath{\forall} n : \ensuremath{\mathbb{N}}, 0 < n \ensuremath{\to} \ensuremath{\forall} z : Fin n \ensuremath{\to} \ensuremath{\mathbb{C}}, (\ensuremath{\forall} i, ‖z i‖ = 1) \ensuremath{\to} c / n \ensuremath{\leq} (signedSumCount z 1 : \ensuremath{\mathbb{R}}) / 2 \textasciicircum{} n


\end{samepage}
\textbf{Source classification:} Known-mathematics formalization / formal-library novelty and usefulness require verification

\textbf{Prior comparison:} Known mathematics; only a genuinely new, faithful, useful, reusable event-window formalization may count. Exact prior-library overlap is not yet independently cleared.

\textbf{Formal/source status:} Packet reports pinned builds and standard axioms; source correspondence and full novelty remain review tasks. Some same-file declarations are explicitly unclaimed duplicates.

\textbf{Research or reuse implication:} Propose this exact scope and family once, subject to the stated source, formal, novelty and deadline checks.

{\small \url{https://github.com/alejandrozu/openmath-2026-judging/blob/d0ed487f487fa7ec814ff705ab55249983fe8707/OpenMath-Judging/review/entries/E08-E395.txt}\par}

\begin{samepage}
\subsection{E08-E698  -  Erdős-Szekeres binomial gcd inequality}
{\small Contributors: Woohyuk Kang, Hyunjin Lee, Sanghyeon Lee, Youchan Oh | Family: M2-E698\par}

\textbf{Recorded scope:} @[category research solved, AMS 5 11] theorem erdos\_698.variants.erdos\_szekeres (n i j : \ensuremath{\mathbb{N}}) (hi : 1 \ensuremath{\leq} i) (hij : i < j) (hj : j \ensuremath{\leq} n / 2) : (n.choose i : \ensuremath{\mathbb{R}}) / (j.choose i : \ensuremath{\mathbb{R}}) \ensuremath{\leq} (Nat.gcd (n.choose i) (n.choose j) : \ensuremath{\mathbb{R}}) \ensuremath{\wedge} (2 : \ensuremath{\mathbb{R}}) \textasciicircum{} i \ensuremath{\leq} (n.choose i : \ensuremath{\mathbb{R}}) / (j.choose i : \ensuremath{\mathbb{R}})


\end{samepage}
\textbf{Source classification:} Known-mathematics formalization / formal-library novelty and usefulness require verification

\textbf{Prior comparison:} Known mathematics; only a genuinely new, faithful, useful, reusable event-window formalization may count. Exact prior-library overlap is not yet independently cleared.

\textbf{Formal/source status:} Packet reports pinned builds and standard axioms; source correspondence and full novelty remain review tasks. Some same-file declarations are explicitly unclaimed duplicates.

\textbf{Research or reuse implication:} Propose this exact scope and family once, subject to the stated source, formal, novelty and deadline checks.

{\small \url{https://github.com/alejandrozu/openmath-2026-judging/blob/d0ed487f487fa7ec814ff705ab55249983fe8707/OpenMath-Judging/review/entries/E08-E698.txt}\par}

\begin{samepage}
\subsection{E08-E939  -  Known r=7,8 powerful-number sum examples}
{\small Contributors: Woohyuk Kang, Hyunjin Lee, Sanghyeon Lee, Youchan Oh | Family: M2-E939\par}

\textbf{Recorded scope:} @[category research solved, AMS 11] theorem erdos\_939.variants.seven : (Erdos939Sums 7).Nonempty
@[category research solved, AMS 11] theorem erdos\_939.variants.eight : (Erdos939Sums 8).Nonempty


\end{samepage}
\textbf{Source classification:} Known-mathematics formalization / formal-library novelty and usefulness require verification

\textbf{Prior comparison:} Known mathematics; only a genuinely new, faithful, useful, reusable event-window formalization may count. Exact prior-library overlap is not yet independently cleared.

\textbf{Formal/source status:} Packet reports pinned builds and standard axioms; source correspondence and full novelty remain review tasks. Some same-file declarations are explicitly unclaimed duplicates.

\textbf{Research or reuse implication:} Propose this exact scope and family once, subject to the stated source, formal, novelty and deadline checks.

{\small \url{https://github.com/alejandrozu/openmath-2026-judging/blob/d0ed487f487fa7ec814ff705ab55249983fe8707/OpenMath-Judging/review/entries/E08-E939.txt}\par}

\begin{samepage}
\subsection{E08-E477  -  Known polynomial-range additive nonunique representation}
{\small Contributors: Woohyuk Kang, Hyunjin Lee, Sanghyeon Lee, Youchan Oh | Family: M2-E477\par}

\textbf{Recorded scope:} @[category research solved, AMS 12] theorem erdos\_477.variants.S\_sq :
    letI f := X \textasciicircum{} 2
    \ensuremath{\forall} A : Set \ensuremath{\mathbb{Z}}, \ensuremath{\exists} z, \ensuremath{\neg} \ensuremath{\exists}! a \ensuremath{\in} A \ensuremath{\times}\ensuremath{^{s}} (Set.range f.eval), z = a.1 + a.2
@[category research solved, AMS 12] theorem erdos\_477.variants.degree\_two\_dvd\_condition\_b\_ne\_zero \{a b c : \ensuremath{\mathbb{Z}}\} (ha : a \ensuremath{\neq} 0) (hb : b \ensuremath{\neq} 0) (hab : a \ensuremath{\mid} b) :
    let f := a - X \textasciicircum{} 2 + b - X + C c
    \ensuremath{\forall} A : Set \ensuremath{\mathbb{Z}}, \ensuremath{\exists} z, \ensuremath{\neg} \ensuremath{\exists}! a \ensuremath{\in} A \ensuremath{\times}\ensuremath{^{s}} (Set.range f.eval), z = a.1 + a.2


\end{samepage}
\textbf{Source classification:} Known-mathematics formalization / formal-library novelty and usefulness require verification

\textbf{Prior comparison:} Known mathematics; only a genuinely new, faithful, useful, reusable event-window formalization may count. Exact prior-library overlap is not yet independently cleared.

\textbf{Formal/source status:} Packet reports pinned builds and standard axioms; source correspondence and full novelty remain review tasks. Some same-file declarations are explicitly unclaimed duplicates.

\textbf{Research or reuse implication:} Propose this exact scope and family once, subject to the stated source, formal, novelty and deadline checks.

{\small \url{https://github.com/alejandrozu/openmath-2026-judging/blob/d0ed487f487fa7ec814ff705ab55249983fe8707/OpenMath-Judging/review/entries/E08-E477.txt}\par}

\begin{samepage}
\subsection{E08-E295  -  Egyptian fraction existence helper}
{\small Contributors: Woohyuk Kang, Hyunjin Lee, Sanghyeon Lee, Youchan Oh | Family: M2-E295\par}

\textbf{Recorded scope:} @[category textbook, AMS 5 11] lemma exists\_k (N : \ensuremath{\mathbb{N}}) : \ensuremath{\exists} (k : \ensuremath{\mathbb{N}}) (n : Fin k \ensuremath{\to} \ensuremath{\mathbb{N}}), (\ensuremath{\forall} i, N \ensuremath{\leq} n i) \ensuremath{\wedge} StrictMono n \ensuremath{\wedge} \ensuremath{\sum} i, (1 / n i : \ensuremath{\mathbb{R}}) = 1


\end{samepage}
\textbf{Source classification:} Known-mathematics formalization / formal-library novelty and usefulness require verification

\textbf{Prior comparison:} Known mathematics; only a genuinely new, faithful, useful, reusable event-window formalization may count. Exact prior-library overlap is not yet independently cleared.

\textbf{Formal/source status:} Packet reports pinned builds and standard axioms; source correspondence and full novelty remain review tasks. Some same-file declarations are explicitly unclaimed duplicates.

\textbf{Research or reuse implication:} Propose this exact scope and family once, subject to the stated source, formal, novelty and deadline checks.

{\small \url{https://github.com/alejandrozu/openmath-2026-judging/blob/d0ed487f487fa7ec814ff705ab55249983fe8707/OpenMath-Judging/review/entries/E08-E295.txt}\par}

\begin{samepage}
\subsection{E08-E703  -  Known intersecting-family t=0 value}
{\small Contributors: Woohyuk Kang, Hyunjin Lee, Sanghyeon Lee, Youchan Oh | Family: M2-E703\par}

\textbf{Recorded scope:} @[category textbook, AMS 5] theorem erdos\_703.variants.zero (n : \ensuremath{\mathbb{N}}) (hn : 1 \ensuremath{\leq} n) : T n 0 = 2 \textasciicircum{} (n - 1)


\end{samepage}
\textbf{Source classification:} Known-mathematics formalization / formal-library novelty and usefulness require verification

\textbf{Prior comparison:} Known mathematics; only a genuinely new, faithful, useful, reusable event-window formalization may count. Exact prior-library overlap is not yet independently cleared.

\textbf{Formal/source status:} Packet reports pinned builds and standard axioms; source correspondence and full novelty remain review tasks. Some same-file declarations are explicitly unclaimed duplicates.

\textbf{Research or reuse implication:} Propose this exact scope and family once, subject to the stated source, formal, novelty and deadline checks.

{\small \url{https://github.com/alejandrozu/openmath-2026-judging/blob/d0ed487f487fa7ec814ff705ab55249983fe8707/OpenMath-Judging/review/entries/E08-E703.txt}\par}

\begin{samepage}
\subsection{E08-E748  -  Known sum-free-set lower bound}
{\small Contributors: Woohyuk Kang, Hyunjin Lee, Sanghyeon Lee, Youchan Oh | Family: M2-E748\par}

\textbf{Recorded scope:} @[category textbook, AMS 5 11] theorem erdos\_748.variants.lower\_bound (n : \ensuremath{\mathbb{N}}) : (2 : \ensuremath{\mathbb{R}}) \textasciicircum{} ((n : \ensuremath{\mathbb{R}}) / 2) \ensuremath{\leq} (f n : \ensuremath{\mathbb{R}})


\end{samepage}
\textbf{Source classification:} Known-mathematics formalization / formal-library novelty and usefulness require verification

\textbf{Prior comparison:} Known mathematics; only a genuinely new, faithful, useful, reusable event-window formalization may count. Exact prior-library overlap is not yet independently cleared.

\textbf{Formal/source status:} Packet reports pinned builds and standard axioms; source correspondence and full novelty remain review tasks. Some same-file declarations are explicitly unclaimed duplicates.

\textbf{Research or reuse implication:} Propose this exact scope and family once, subject to the stated source, formal, novelty and deadline checks.

{\small \url{https://github.com/alejandrozu/openmath-2026-judging/blob/d0ed487f487fa7ec814ff705ab55249983fe8707/OpenMath-Judging/review/entries/E08-E748.txt}\par}

\begin{samepage}
\subsection{E08-E1136  -  Multiples-of-three density baseline}
{\small Contributors: Woohyuk Kang, Hyunjin Lee, Sanghyeon Lee, Youchan Oh | Family: M2-E1136\par}

\textbf{Recorded scope:} @[category research solved, AMS 11] theorem erdos\_1136.variants.multiples\_of\_three : AvoidsPowersOfTwo \{n : \ensuremath{\mathbb{N}} | 3 \ensuremath{\mid} n\} \ensuremath{\wedge} Set.HasDensity \{n : \ensuremath{\mathbb{N}} | 3 \ensuremath{\mid} n\} (1 / 3)


\end{samepage}
\textbf{Source classification:} Known-mathematics formalization / formal-library novelty and usefulness require verification

\textbf{Prior comparison:} Known mathematics; only a genuinely new, faithful, useful, reusable event-window formalization may count. Exact prior-library overlap is not yet independently cleared.

\textbf{Formal/source status:} Packet reports pinned builds and standard axioms; source correspondence and full novelty remain review tasks. Some same-file declarations are explicitly unclaimed duplicates.

\textbf{Research or reuse implication:} Propose this exact scope and family once, subject to the stated source, formal, novelty and deadline checks.

{\small \url{https://github.com/alejandrozu/openmath-2026-judging/blob/d0ed487f487fa7ec814ff705ab55249983fe8707/OpenMath-Judging/review/entries/E08-E1136.txt}\par}

\begin{samepage}
\subsection{E08-E358  -  Sylvester consecutive sums / odd divisors}
{\small Contributors: Woohyuk Kang, Hyunjin Lee, Sanghyeon Lee, Youchan Oh | Family: M2-E358\par}

\textbf{Recorded scope:} @[category textbook, AMS 5 11] theorem f\_id : f id = fun n \ensuremath{\mapsto} \#\{d \ensuremath{\in} n.divisors | Odd d\}


\end{samepage}
\textbf{Source classification:} Known-mathematics formalization / formal-library novelty and usefulness require verification

\textbf{Prior comparison:} Known mathematics; only a genuinely new, faithful, useful, reusable event-window formalization may count. Exact prior-library overlap is not yet independently cleared.

\textbf{Formal/source status:} Packet reports pinned builds and standard axioms; source correspondence and full novelty remain review tasks. Some same-file declarations are explicitly unclaimed duplicates.

\textbf{Research or reuse implication:} Propose this exact scope and family once, subject to the stated source, formal, novelty and deadline checks.

{\small \url{https://github.com/alejandrozu/openmath-2026-judging/blob/d0ed487f487fa7ec814ff705ab55249983fe8707/OpenMath-Judging/review/entries/E08-E358.txt}\par}

\begin{samepage}
\subsection{E08-E619  -  Maximal triangle-free supergraph diameter bound}
{\small Contributors: Woohyuk Kang, Hyunjin Lee, Sanghyeon Lee, Youchan Oh | Family: M2-E619\par}

\textbf{Recorded scope:} @[category research solved, AMS 5] theorem erdos\_619.variants.add\_edges\_diam\_three \{V : Type*\} [Fintype V] (G : SimpleGraph V) (hG : G.Connected) (hG' : G.CliqueFree 3) :
    \ensuremath{\exists} H : SimpleGraph V, G \ensuremath{\leq} H \ensuremath{\wedge} H.CliqueFree 3 \ensuremath{\wedge} H.ediam \ensuremath{\leq} 3


\end{samepage}
\textbf{Source classification:} Known-mathematics formalization / formal-library novelty and usefulness require verification

\textbf{Prior comparison:} Known mathematics; only a genuinely new, faithful, useful, reusable event-window formalization may count. Exact prior-library overlap is not yet independently cleared.

\begin{samepage}
\textbf{Formal/source status:} Packet reports pinned builds and standard axioms; source correspondence and full novelty remain review tasks. Some same-file declarations are explicitly unclaimed duplicates.

\end{samepage}
\textbf{Research or reuse implication:} Propose this exact scope and family once, subject to the stated source, formal, novelty and deadline checks.

{\small \url{https://github.com/alejandrozu/openmath-2026-judging/blob/d0ed487f487fa7ec814ff705ab55249983fe8707/OpenMath-Judging/review/entries/E08-E619.txt}\par}

\begin{samepage}
\subsection{E08-E1148  -  Known difference-of-squares weaker bound}
{\small Contributors: Woohyuk Kang, Hyunjin Lee, Sanghyeon Lee, Youchan Oh | Family: M2-E1148\par}

\textbf{Recorded scope:} @[category research solved, AMS 11] theorem erdos\_1148.variants.weaker : \ensuremath{\forall} n, erdos\_1148\_weaker\_prop n


\end{samepage}
\textbf{Source classification:} Known-mathematics formalization / formal-library novelty and usefulness require verification

\textbf{Prior comparison:} Known mathematics; only a genuinely new, faithful, useful, reusable event-window formalization may count. Exact prior-library overlap is not yet independently cleared.

\textbf{Formal/source status:} Packet reports pinned builds and standard axioms; source correspondence and full novelty remain review tasks. Some same-file declarations are explicitly unclaimed duplicates.

\textbf{Research or reuse implication:} Propose this exact scope and family once, subject to the stated source, formal, novelty and deadline checks.

{\small \url{https://github.com/alejandrozu/openmath-2026-judging/blob/d0ed487f487fa7ec814ff705ab55249983fe8707/OpenMath-Judging/review/entries/E08-E1148.txt}\par}

\begin{samepage}
\subsection{E08-BB6  -  Finite six-state baseline witness}
{\small Contributors: Woohyuk Kang, Hyunjin Lee, Sanghyeon Lee, Youchan Oh | Family: BB6\par}

\textbf{Recorded scope:} 249881steps554ones735span, no global record

\end{samepage}
\textbf{Source classification:} Baseline/known-result candidate, or incomplete claim with prior overlap; see exact evidence

\textbf{Prior comparison:} No eligible original mathematical delta established in the available packet. This is not a finding that all underlying work lacks novelty.

\textbf{Formal/source status:} Exact arithmetic certificates/source available for review; independent Collatz check 234/234. No fresh full kernel acceptance.

\textbf{Research or reuse implication:} Propose this exact scope and family once, subject to the stated source, formal, novelty and deadline checks.

{\small \url{https://github.com/alejandrozu/openmath-2026-judging/blob/d0ed487f487fa7ec814ff705ab55249983fe8707/OpenMath-Judging/review/entries/E08-BB6.txt}\par}

\begin{samepage}
\subsection{E08-CHSH  -  Known CHSH witness}
{\small Contributors: Woohyuk Kang, Hyunjin Lee, Sanghyeon Lee, Youchan Oh | Family: GROTHENDIECK\par}

\textbf{Recorded scope:} \ensuremath{\surd}2-scale rational certificate; not new general bound

\end{samepage}
\textbf{Source classification:} Baseline/known-result candidate, or incomplete claim with prior overlap; see exact evidence

\textbf{Prior comparison:} No eligible original mathematical delta established in the available packet. This is not a finding that all underlying work lacks novelty.

\textbf{Formal/source status:} Exact arithmetic certificates/source available for review; independent Collatz check 234/234. No fresh full kernel acceptance.

\textbf{Research or reuse implication:} Propose this exact scope and family once, subject to the stated source, formal, novelty and deadline checks.

{\small \url{https://github.com/alejandrozu/openmath-2026-judging/blob/d0ed487f487fa7ec814ff705ab55249983fe8707/OpenMath-Judging/review/entries/E08-CHSH.txt}\par}

\begin{samepage}
\subsection{E08-MATRIX 139  -  Known support 139 tensor}
{\small Contributors: Woohyuk Kang, Hyunjin Lee, Sanghyeon Lee, Youchan Oh | Family: MATRIX-3\par}

\textbf{Recorded scope:} Rank 23/support 139 baseline

\end{samepage}
\textbf{Source classification:} Baseline/known-result candidate, or incomplete claim with prior overlap; see exact evidence

\textbf{Prior comparison:} No eligible original mathematical delta established in the available packet. This is not a finding that all underlying work lacks novelty.

\textbf{Formal/source status:} Exact arithmetic certificates/source available for review; independent Collatz check 234/234. No fresh full kernel acceptance.

\textbf{Research or reuse implication:} Propose this exact scope and family once, subject to the stated source, formal, novelty and deadline checks.

{\small \url{https://github.com/alejandrozu/openmath-2026-judging/blob/d0ed487f487fa7ec814ff705ab55249983fe8707/OpenMath-Judging/review/entries/E08-MATRIX139.txt}\par}

\begin{samepage}
\subsection{E08-COLLATZ 234  -  234-rule baseline-equivalent coverage}
{\small Contributors: Woohyuk Kang, Hyunjin Lee, Sanghyeon Lee, Youchan Oh | Family: COLLATZ\par}

\textbf{Recorded scope:} 234rules, independently1765odd classes mod 4096: exactly Chaewon union

\end{samepage}
\textbf{Source classification:} Baseline/known-result candidate, or incomplete claim with prior overlap; see exact evidence

\textbf{Prior comparison:} No eligible original mathematical delta established in the available packet. This is not a finding that all underlying work lacks novelty.

\textbf{Formal/source status:} Exact arithmetic certificates/source available for review; independent Collatz check 234/234. No fresh full kernel acceptance.

\textbf{Research or reuse implication:} Propose this exact scope and family once, subject to the stated source, formal, novelty and deadline checks.

{\small \url{https://github.com/alejandrozu/openmath-2026-judging/blob/d0ed487f487fa7ec814ff705ab55249983fe8707/OpenMath-Judging/review/entries/E08-COLLATZ234.txt}\par}

\begin{samepage}
\subsection{E08-KOBON 18  -  Known n18=93 reconstruction}
{\small Contributors: Woohyuk Kang, Hyunjin Lee, Sanghyeon Lee, Youchan Oh | Family: KOBON\par}

\textbf{Recorded scope:} n18=93, known lower bound

\end{samepage}
\textbf{Source classification:} Baseline/known-result candidate, or incomplete claim with prior overlap; see exact evidence

\textbf{Prior comparison:} No eligible original mathematical delta established in the available packet. This is not a finding that all underlying work lacks novelty.

\textbf{Formal/source status:} Exact arithmetic certificates/source available for review; independent Collatz check 234/234. No fresh full kernel acceptance.

\textbf{Research or reuse implication:} Propose this exact scope and family once, subject to the stated source, formal, novelty and deadline checks.

{\small \url{https://github.com/alejandrozu/openmath-2026-judging/blob/d0ed487f487fa7ec814ff705ab55249983fe8707/OpenMath-Judging/review/entries/E08-KOBON18.txt}\par}

\begin{samepage}
\subsection{E09-UPPER 94  -  Arbitrary K(18)\ensuremath{\leq}94}
{\small Contributors: Raj Harshit Srirangam | Family: KOBON\par}

\textbf{Recorded scope:} All straight-line/pseudoline arrangements including arbitrary multiple points have\ensuremath{\leq}94 triangular faces.

\end{samepage}
\textbf{Source classification:} Unresolved or missing-artifact record; nonacceptance does not prove nonnovelty

\textbf{Prior comparison:} Candidate new extension beyond simple-arrangement bounds; specific multiplicity scope is the key difference.

\textbf{Formal/source status:} Rational LP/DP/SAT/DRAT reports are retained, but complete geometry\ensuremath{\to}constraint\ensuremath{\to}certificate soundness chain has not received human or Lean review. The open 93 stage is distinct from claimed\ensuremath{\leq}94 and must not erase that claim.

\textbf{Research or reuse implication:} Current formal-only proposal 0. The README explicitly states no formal proof and unproved guarded geometric facts. Potential parent union 17.8695(p=.05) only if the entire approved geometry-to-certificate chain is accepted; not one award per finite order.

{\small \url{https://github.com/alejandrozu/openmath-2026-judging/blob/d0ed487f487fa7ec814ff705ab55249983fe8707/OpenMath-Judging/review/entries/E09-UPPER94.txt}\par}

\begin{samepage}
\subsection{E09-TRIPLES 93  -  Multiplicity\ensuremath{\leq}3 upper 93}
{\small Contributors: Raj Harshit Srirangam | Family: KOBON\par}

\textbf{Recorded scope:} T\ensuremath{\leq}93 at 18lines without multiplicity\ensuremath{\geq}4; sharp using known 93 lower bound.

\end{samepage}
\textbf{Source classification:} Baseline/known-result candidate, or incomplete claim with prior overlap; see exact evidence

\textbf{Prior comparison:} Candidate scoped mathematical delta; exact baseline and independently sealed provenance are required. A missing gallery value is not evidence of a world record.

\textbf{Formal/source status:} Rational LP/DP/SAT/DRAT reports are retained, but complete geometry\ensuremath{\to}constraint\ensuremath{\to}certificate soundness chain has not received human or Lean review. The open 93 stage is distinct from claimed\ensuremath{\leq}94 and must not erase that claim.

\textbf{Research or reuse implication:} Current formal-only proposal 0. The README explicitly states no formal proof and unproved guarded geometric facts. Potential parent union 17.8695(p=.05) only if the entire approved geometry-to-certificate chain is accepted; not one award per finite order.

{\small \url{https://github.com/alejandrozu/openmath-2026-judging/blob/d0ed487f487fa7ec814ff705ab55249983fe8707/OpenMath-Judging/review/entries/E09-TRIPLES93.txt}\par}

\begin{samepage}
\subsection{E09-K14  -  Claimed complete finite n14 case}
{\small Contributors: Raj Harshit Srirangam | Family: KOBON\par}

\textbf{Recorded scope:} K(14)=54 arbitrary multiplicities; finite child case of global Kobon.

\end{samepage}
\textbf{Source classification:} Baseline/known-result candidate, or incomplete claim with prior overlap; see exact evidence

\textbf{Prior comparison:} Candidate new arbitrary-multiplicity closure, subject to actual historical scope and certificate coverage.

\textbf{Formal/source status:} Rational LP/DP/SAT/DRAT reports are retained, but complete geometry\ensuremath{\to}constraint\ensuremath{\to}certificate soundness chain has not received human or Lean review. The open 93 stage is distinct from claimed\ensuremath{\leq}94 and must not erase that claim.

\textbf{Research or reuse implication:} Current formal-only proposal 0. The README explicitly states no formal proof and unproved guarded geometric facts. Potential parent union 17.8695(p=.05) only if the entire approved geometry-to-certificate chain is accepted; not one award per finite order.

{\small \url{https://github.com/alejandrozu/openmath-2026-judging/blob/d0ed487f487fa7ec814ff705ab55249983fe8707/OpenMath-Judging/review/entries/E09-K14.txt}\par}

\begin{samepage}
\subsection{E09-EVEN  -  Triple-point bounds across even n}
{\small Contributors: Raj Harshit Srirangam | Family: KOBON\par}

\textbf{Recorded scope:} Restricted BBL bound for even n6..40; local budget/automaton/DP lemmas.

\end{samepage}
\textbf{Source classification:} Baseline/known-result candidate, or incomplete claim with prior overlap; see exact evidence

\textbf{Prior comparison:} Candidate scoped mathematical delta; exact baseline and independently sealed provenance are required. A missing gallery value is not evidence of a world record.

\textbf{Formal/source status:} Rational LP/DP/SAT/DRAT reports are retained, but complete geometry\ensuremath{\to}constraint\ensuremath{\to}certificate soundness chain has not received human or Lean review. The open 93 stage is distinct from claimed\ensuremath{\leq}94 and must not erase that claim.

\textbf{Research or reuse implication:} Current formal-only proposal 0. The README explicitly states no formal proof and unproved guarded geometric facts. Potential parent union 17.8695(p=.05) only if the entire approved geometry-to-certificate chain is accepted; not one award per finite order.

{\small \url{https://github.com/alejandrozu/openmath-2026-judging/blob/d0ed487f487fa7ec814ff705ab55249983fe8707/OpenMath-Judging/review/entries/E09-EVEN.txt}\par}

\begin{samepage}
\subsection{E09-ALL 8  -  Residual K18=93 obstruction}
{\small Contributors: Raj Harshit Srirangam | Family: KOBON\par}

\textbf{Recorded scope:} Stage 1 capacity inequality unresolved;281/416UNSAT cubes,135undecided; K18=93 not proved.

\end{samepage}
\textbf{Source classification:} Baseline/known-result candidate, or incomplete claim with prior overlap; see exact evidence

\textbf{Prior comparison:} Candidate scoped mathematical delta; exact baseline and independently sealed provenance are required. A missing gallery value is not evidence of a world record.

\textbf{Formal/source status:} Rational LP/DP/SAT/DRAT reports are retained, but complete geometry\ensuremath{\to}constraint\ensuremath{\to}certificate soundness chain has not received human or Lean review. The open 93 stage is distinct from claimed\ensuremath{\leq}94 and must not erase that claim.

\textbf{Research or reuse implication:} Current formal-only proposal 0. The README explicitly states no formal proof and unproved guarded geometric facts. Potential parent union 17.8695(p=.05) only if the entire approved geometry-to-certificate chain is accepted; not one award per finite order.

{\small \url{https://github.com/alejandrozu/openmath-2026-judging/blob/d0ed487f487fa7ec814ff705ab55249983fe8707/OpenMath-Judging/review/entries/E09-ALL8.txt}\par}

\begin{samepage}
\subsection{E09-K10  -  Known K10 / n18 construction}
{\small Contributors: Raj Harshit Srirangam | Family: KOBON\par}

\textbf{Recorded scope:} K10=25 and n18=93 record reconstructions.

\end{samepage}
\textbf{Source classification:} Baseline/known-result candidate, or incomplete claim with prior overlap; see exact evidence

\textbf{Prior comparison:} No eligible original mathematical delta established in the available packet. This is not a finding that all underlying work lacks novelty.

\textbf{Formal/source status:} Certificates may be reusable research evidence.

\textbf{Research or reuse implication:} Current formal-only proposal 0. The README explicitly states no formal proof and unproved guarded geometric facts. Potential parent union 17.8695(p=.05) only if the entire approved geometry-to-certificate chain is accepted; not one award per finite order.

{\small \url{https://github.com/alejandrozu/openmath-2026-judging/blob/d0ed487f487fa7ec814ff705ab55249983fe8707/OpenMath-Judging/review/entries/E09-K10.txt}\par}

\begin{samepage}
\subsection{E10-H7  -  AP conclusion representation bridge}
{\small Contributors: Jamie Steeg | Family: M2-AP-CONCLUSION\par}

\textbf{Recorded scope:} ap\_subprogression\_of\_le; ap\_frequently\_iff\_every\_length.

\end{samepage}
\textbf{Source classification:} Known-mathematics formalization / formal-library novelty and usefulness require verification

\textbf{Prior comparison:} Known mathematics; only a genuinely new, faithful, useful, reusable event-window formalization may count. Exact prior-library overlap is not yet independently cleared.

\textbf{Formal/source status:} Source and pinned compile/standard-axiom receipts present; final hill/workflow fields null. No fresh build.

\textbf{Research or reuse implication:} Propose this exact scope and family once, subject to the stated source, formal, novelty and deadline checks.

{\small \url{https://github.com/alejandrozu/openmath-2026-judging/blob/d0ed487f487fa7ec814ff705ab55249983fe8707/OpenMath-Judging/review/entries/E10-H7.txt}\par}

\begin{samepage}
\subsection{E10-H1PART  -  Conditional hypothetical 95 exclusion}
{\small Contributors: Jamie Steeg | Family: KOBON\par}

\textbf{Recorded scope:} n18 hypothetical 95/q\ensuremath{\leq}5 exclusion under predecessor local-fan/clean-line lemmas; q\ensuremath{\geq}6open.

\end{samepage}
\textbf{Source classification:} Unresolved or missing-artifact record; nonacceptance does not prove nonnovelty

\textbf{Prior comparison:} Conditional organization of the obstruction may be useful, but no completed eligible original advance is established.

\textbf{Formal/source status:} Predecessor facts/return evidence and replay are not all independently discharged.

\textbf{Research or reuse implication:} Score 0 for the original family until independently accepted formal geometric subclaims are identified; retain the potential P1=.01 scenario only as a hold note.

{\small \url{https://github.com/alejandrozu/openmath-2026-judging/blob/d0ed487f487fa7ec814ff705ab55249983fe8707/OpenMath-Judging/review/entries/E10-H1PART.txt}\par}

\begin{samepage}
\subsection{E10-H1BASE  -  Known Bader 93}
{\small Contributors: Jamie Steeg | Family: KOBON\par}

\textbf{Recorded scope:} Known18-line93reconstruction

\end{samepage}
\textbf{Source classification:} Baseline/known-result candidate, or incomplete claim with prior overlap; see exact evidence

\textbf{Prior comparison:} No eligible original mathematical delta established in the available packet. This is not a finding that all underlying work lacks novelty.

\textbf{Formal/source status:} Queued native reports and prepared packet are not official final receipts.

\textbf{Research or reuse implication:} Propose this exact scope and family once, subject to the stated source, formal, novelty and deadline checks.

{\small \url{https://github.com/alejandrozu/openmath-2026-judging/blob/d0ed487f487fa7ec814ff705ab55249983fe8707/OpenMath-Judging/review/entries/E10-H1BASE.txt}\par}

\begin{samepage}
\subsection{E10-H2  -  Finite baseline witness}
{\small Contributors: Jamie Steeg | Family: BB6\par}

\textbf{Recorded scope:} 89911steps185ones541span, not review-ready

\end{samepage}
\textbf{Source classification:} Baseline/known-result candidate, or incomplete claim with prior overlap; see exact evidence

\textbf{Prior comparison:} No eligible original mathematical delta established in the available packet. This is not a finding that all underlying work lacks novelty.

\textbf{Formal/source status:} Queued native reports and prepared packet are not official final receipts.

\textbf{Research or reuse implication:} Propose this exact scope and family once, subject to the stated source, formal, novelty and deadline checks.

{\small \url{https://github.com/alejandrozu/openmath-2026-judging/blob/d0ed487f487fa7ec814ff705ab55249983fe8707/OpenMath-Judging/review/entries/E10-H2.txt}\par}

\begin{samepage}
\subsection{E10-H5  -  CHSH baseline}
{\small Contributors: Jamie Steeg | Family: GROTHENDIECK\par}

\textbf{Recorded scope:} 275807/195025 ratio, known \ensuremath{\surd}2 scale

\end{samepage}
\textbf{Source classification:} Baseline/known-result candidate, or incomplete claim with prior overlap; see exact evidence

\textbf{Prior comparison:} No eligible original mathematical delta established in the available packet. This is not a finding that all underlying work lacks novelty.

\textbf{Formal/source status:} Queued native reports and prepared packet are not official final receipts.

\textbf{Research or reuse implication:} Propose this exact scope and family once, subject to the stated source, formal, novelty and deadline checks.

{\small \url{https://github.com/alejandrozu/openmath-2026-judging/blob/d0ed487f487fa7ec814ff705ab55249983fe8707/OpenMath-Judging/review/entries/E10-H5.txt}\par}

\begin{samepage}
\subsection{E10-H6  -  Laderman baseline}
{\small Contributors: Jamie Steeg | Family: MATRIX-3\par}

\textbf{Recorded scope:} Rank 23support 153

\end{samepage}
\textbf{Source classification:} Baseline/known-result candidate, or incomplete claim with prior overlap; see exact evidence

\textbf{Prior comparison:} No eligible original mathematical delta established in the available packet. This is not a finding that all underlying work lacks novelty.

\textbf{Formal/source status:} Queued native reports and prepared packet are not official final receipts.

\textbf{Research or reuse implication:} Propose this exact scope and family once, subject to the stated source, formal, novelty and deadline checks.

{\small \url{https://github.com/alejandrozu/openmath-2026-judging/blob/d0ed487f487fa7ec814ff705ab55249983fe8707/OpenMath-Judging/review/entries/E10-H6.txt}\par}

\begin{samepage}
\subsection{E11-THREE  -  Known length-three case}
{\small Contributors: Leon Koerbs, Cameron Fen | Family: M2-E3-THREE\par}

\textbf{Recorded scope:} Formal divergent-reciprocal implication for 3-term APs using external Kelley-Meka formalization.

\end{samepage}
\textbf{Source classification:} Known-mathematics formalization / formal-library novelty and usefulness require verification

\textbf{Prior comparison:} Mathematical result known; formal integration novelty requires baseline and imported-source comparison.

\textbf{Formal/source status:} Source includes vendored dependencies and fetch\_plby route; external 44files not redistributed. Author reports standard-axiom check.

\textbf{Research or reuse implication:} M2 candidate only, deduplicated from supporting dyadic/AP infrastructure.

{\small \url{https://github.com/alejandrozu/openmath-2026-judging/blob/d0ed487f487fa7ec814ff705ab55249983fe8707/OpenMath-Judging/review/entries/E11-THREE.txt}\par}

\begin{samepage}
\subsection{E11-DYADIC  -  Per-length extremal equivalence}
{\small Contributors: Leon Koerbs, Cameron Fen | Family: M2-E3-DYADIC\par}

\textbf{Recorded scope:} For everyK\ensuremath{\geq}3, dyadic rK(2\textasciicircum{}j)/2\textasciicircum{}j summability iff every divergent set containsK-term AP.

\end{samepage}
\textbf{Source classification:} Known-mathematics formalization / formal-library novelty and usefulness require verification

\textbf{Prior comparison:} Candidate reusable formalization/reformulation; sealed independence or distinct extension controls separate entrant credit.

\textbf{Formal/source status:} Converse.lean provides named closed maxCardDyadicSummable\_iff. Earlier lean/equivalence/solution.lean ends in a real sorry on the parent goal; do not advertise that as full proof or reject the separately closed Converse endpoint.

\textbf{Research or reuse implication:} One possible M2 family, with exact closed endpoint and no all-length solution credit.

{\small \url{https://github.com/alejandrozu/openmath-2026-judging/blob/d0ed487f487fa7ec814ff705ab55249983fe8707/OpenMath-Judging/review/entries/E11-DYADIC.txt}\par}

\begin{samepage}
\subsection{E11-GOWERS  -  Four-term analytic groundwork}
{\small Contributors: Leon Koerbs, Cameron Fen | Family: M2-GOWERS 4\par}

\textbf{Recorded scope:} norm\_apCount4\_pow8\_le\_gowers3 and associated closed counterexample-audit lemmas.

\end{samepage}
\textbf{Source classification:} Known-mathematics formalization / formal-library novelty and usefulness require verification

\textbf{Prior comparison:} Known mathematics; only a genuinely new, faithful, useful, reusable event-window formalization may count. Exact prior-library overlap is not yet independently cleared.

\textbf{Formal/source status:} Submitted module described as compiling without holes; exact axiom output for every endpoint not yet independently rerun.

\textbf{Research or reuse implication:} Propose this exact scope and family once, subject to the stated source, formal, novelty and deadline checks.

{\small \url{https://github.com/alejandrozu/openmath-2026-judging/blob/d0ed487f487fa7ec814ff705ab55249983fe8707/OpenMath-Judging/review/entries/E11-GOWERS.txt}\par}

\begin{samepage}
\subsection{E11-HARMONIC  -  Finite four-AP-free harmonic construction}
{\small Contributors: Leon Koerbs, Cameron Fen | Family: ERDOS 169\par}

\textbf{Recorded scope:} Certified finite reciprocal sum 4.314 atN\ensuremath{\leq}10\textasciicircum{}6.

\end{samepage}
\textbf{Source classification:} Unresolved or missing-artifact record; nonacceptance does not prove nonnovelty

\textbf{Prior comparison:} No eligible original mathematical delta established in the available packet. This is not a finding that all underlying work lacks novelty.

\textbf{Formal/source status:} Hill integer/bitmap checks reported; finite harmonic sum not a divergent series.

\textbf{Research or reuse implication:} Propose this exact scope and family once, subject to the stated source, formal, novelty and deadline checks.

{\small \url{https://github.com/alejandrozu/openmath-2026-judging/blob/d0ed487f487fa7ec814ff705ab55249983fe8707/OpenMath-Judging/review/entries/E11-HARMONIC.txt}\par}

\begin{samepage}
\subsection{E12-HEILBRONN  -  Incomplete nine-point cover}
{\small Contributors: Frederik Smits van Oyen | Family: HEILBRONN 9\par}

\textbf{Recorded scope:} Target0.027426211734693878; boundary uncovered;2581unresolved529missing leaves.

\end{samepage}
\textbf{Source classification:} Unresolved or missing-artifact record; nonacceptance does not prove nonnovelty

\textbf{Prior comparison:} Potential partial covering/certificate engineering; no accepted global theorem or complete new bound established.

\textbf{Formal/source status:} Author floating-bound audit uses tolerance/protocol; no complete formal global certificate.

\textbf{Research or reuse implication:} Open score 0. A separately specified exact certified local subproblem can be reconsidered; do not assign the attempted global target p>0 merely from coverage percentages.

{\small \url{https://github.com/alejandrozu/openmath-2026-judging/blob/d0ed487f487fa7ec814ff705ab55249983fe8707/OpenMath-Judging/review/entries/E12-HEILBRONN.txt}\par}

\begin{samepage}
\subsection{E15-RAMSEY  -  Scored Ramsey lead}
{\small Contributors:  | Family: RAMSEY-K4\par}

\textbf{Recorded scope:} density\_ppt30142250482; public repo only.gitkeep pre-event.

\end{samepage}
\textbf{Source classification:} Baseline/known-result candidate, or incomplete claim with prior overlap; see exact evidence

\textbf{Prior comparison:} Candidate scoped mathematical delta; exact baseline and independently sealed provenance are required. A missing gallery value is not evidence of a world record.

\textbf{Formal/source status:} No current exact template or proof.

\textbf{Research or reuse implication:} Hold missing payload; current proposed points 0. Board metadata cannot establish an eligible exact theorem.

{\small \url{https://github.com/alejandrozu/openmath-2026-judging/blob/d0ed487f487fa7ec814ff705ab55249983fe8707/OpenMath-Judging/review/entries/E15-RAMSEY.txt}\par}

\begin{samepage}
\subsection{E17-CORE 1  -  \texttt{Orion.potential\_along\_word}, \texttt{Orion.lower\_bound}, \texttt{Orion.cardinality\_limit}}
{\small Contributors: Qichao Wang | Family: CERNY-PART-1\par}

\textbf{Recorded scope:} For arbitrary step systems, local potential inequalities telescope along every list; a zero-potential goal forces word length \ensuremath{\geq} initial potential. The cardinality-slope lemma is conditional on its stated slope assumptions.

\end{samepage}
\textbf{Source classification:} Known-mathematics formalization / formal-library novelty and usefulness require verification

\textbf{Prior comparison:} Known mathematics; only a genuinely new, faithful, useful, reusable event-window formalization may count. Exact prior-library overlap is not yet independently cleared.

\textbf{Formal/source status:} Only four manifests, not the 27MBsourceZIP. Compiler/axiom logs are author descriptions. No local Lean source or fresh fidelity check.

\textbf{Research or reuse implication:} Propose this exact scope and family once, subject to the stated source, formal, novelty and deadline checks.

{\small \url{https://github.com/alejandrozu/openmath-2026-judging/blob/d0ed487f487fa7ec814ff705ab55249983fe8707/OpenMath-Judging/review/entries/E17-CORE1.txt}\par}

\begin{samepage}
\subsection{E17-CORE 2  -  \texttt{OrionResearch.potential\_lower\_bound}, \texttt{OrionResearch.quotient\_local}, \texttt{OrionResearch.monotone\_simulation\_transfer}, `OrionResearc}
{\small Contributors: Qichao Wang | Family: CERNY-PART-2\par}

\textbf{Recorded scope:} Arbitrary relational walks and existential quotient-edge local transfer; inclusion-type monotone simulation; finite dyadic timing lower bounds, optimal-tail rigidity and buffer arithmetic; eventual constancy of decreasing natural sequences. The genealogy timing assumptions are explicit.

\end{samepage}
\textbf{Source classification:} Known-mathematics formalization / formal-library novelty and usefulness require verification

\textbf{Prior comparison:} Known mathematics; only a genuinely new, faithful, useful, reusable event-window formalization may count. Exact prior-library overlap is not yet independently cleared.

\textbf{Formal/source status:} Only four manifests, not the 27MBsourceZIP. Compiler/axiom logs are author descriptions. No local Lean source or fresh fidelity check.

\textbf{Research or reuse implication:} Propose this exact scope and family once, subject to the stated source, formal, novelty and deadline checks.

{\small \url{https://github.com/alejandrozu/openmath-2026-judging/blob/d0ed487f487fa7ec814ff705ab55249983fe8707/OpenMath-Judging/review/entries/E17-CORE2.txt}\par}

\begin{samepage}
\subsection{E17-CORE 3  -  \texttt{OrionAnchors.p2\_rotation}, \texttt{OrionAnchors.p2\_merge}, \texttt{OrionAnchors.p3\_rotation}, \texttt{OrionAnchors.p3\_merge}, `OrionAnchors.p3\_full\_va}
{\small Contributors: Qichao Wang | Family: CERNY-PART-3\par}

\textbf{Recorded scope:} Symbolic two-/three-anchor rotation and merge inequalities for all valid cardinalities and bits, and the full three-anchor potential value for n\ensuremath{\geq}5.

\end{samepage}
\textbf{Source classification:} Known-mathematics formalization / formal-library novelty and usefulness require verification

\textbf{Prior comparison:} Known mathematics; only a genuinely new, faithful, useful, reusable event-window formalization may count. Exact prior-library overlap is not yet independently cleared.

\textbf{Formal/source status:} Only four manifests, not the 27MBsourceZIP. Compiler/axiom logs are author descriptions. No local Lean source or fresh fidelity check.

\textbf{Research or reuse implication:} Propose this exact scope and family once, subject to the stated source, formal, novelty and deadline checks.

{\small \url{https://github.com/alejandrozu/openmath-2026-judging/blob/d0ed487f487fa7ec814ff705ab55249983fe8707/OpenMath-Judging/review/entries/E17-CORE3.txt}\par}

\begin{samepage}
\subsection{E17-CORE 4  -  \texttt{OrionThreeGraph.one\_step}, \texttt{OrionThreeGraph.full\_lower\_bound}}
{\small Contributors: Qichao Wang | Family: CERNY-PART-4\par}

\textbf{Recorded scope:} In the explicitly defined count/window graph, any walk from the full feature to a cardinality-one goal has length \ensuremath{\geq}3n\ensuremath{-}4, n\ensuremath{\geq}5. This is a lower bound in that graph, not the general exact quotient formula.

\end{samepage}
\textbf{Source classification:} Known-mathematics formalization / formal-library novelty and usefulness require verification

\textbf{Prior comparison:} Known mathematics; only a genuinely new, faithful, useful, reusable event-window formalization may count. Exact prior-library overlap is not yet independently cleared.

\textbf{Formal/source status:} Only four manifests, not the 27MBsourceZIP. Compiler/axiom logs are author descriptions. No local Lean source or fresh fidelity check.

\textbf{Research or reuse implication:} Propose this exact scope and family once, subject to the stated source, formal, novelty and deadline checks.

{\small \url{https://github.com/alejandrozu/openmath-2026-judging/blob/d0ed487f487fa7ec814ff705ab55249983fe8707/OpenMath-Judging/review/entries/E17-CORE4.txt}\par}

\begin{samepage}
\subsection{E17-CORE 5  -  \texttt{OrionExample.one\_step\_fidelity}, \texttt{OrionExample.optimal\_subset\_path} (separate namespaces in separate compilation units)}
{\small Contributors: Qichao Wang | Family: CERNY-PART-5\par}

\textbf{Recorded scope:} In the separately defined finite C3 and C4 powerset tables, the listed witnesses attain the singleton goal and every goal-reaching list is at least as long: lengths 4 and 9 respectively. Every table edge has a checked one-step concrete subset-image correspondence.

\end{samepage}
\textbf{Source classification:} Known-mathematics formalization / formal-library novelty and usefulness require verification

\textbf{Prior comparison:} Known mathematics; only a genuinely new, faithful, useful, reusable event-window formalization may count. Exact prior-library overlap is not yet independently cleared.

\textbf{Formal/source status:} Only four manifests, not the 27MBsourceZIP. Compiler/axiom logs are author descriptions. No local Lean source or fresh fidelity check.

\textbf{Research or reuse implication:} Propose this exact scope and family once, subject to the stated source, formal, novelty and deadline checks.

{\small \url{https://github.com/alejandrozu/openmath-2026-judging/blob/d0ed487f487fa7ec814ff705ab55249983fe8707/OpenMath-Judging/review/entries/E17-CORE5.txt}\par}

\begin{samepage}
\subsection{E17-CORE 6  -  \texttt{Orion.DirectAttack.layered\_invariants}, \texttt{Orion.DirectAttack.first\_collapse\_cost}}
{\small Contributors: Qichao Wang | Family: CERNY-PART-6\par}

\textbf{Recorded scope:} For any \texttt{LayeredCertificate} satisfying all local index, internal/entrance-potential and K/U entrance hypotheses, the five invariants hold along every path. A path ending in K or U followed by its collapse letter costs at least the requested bound.

\end{samepage}
\textbf{Source classification:} Known-mathematics formalization / formal-library novelty and usefulness require verification

\textbf{Prior comparison:} Known mathematics; only a genuinely new, faithful, useful, reusable event-window formalization may count. Exact prior-library overlap is not yet independently cleared.

\textbf{Formal/source status:} Only four manifests, not the 27MBsourceZIP. Compiler/axiom logs are author descriptions. No local Lean source or fresh fidelity check.

\textbf{Research or reuse implication:} Propose this exact scope and family once, subject to the stated source, formal, novelty and deadline checks.

{\small \url{https://github.com/alejandrozu/openmath-2026-judging/blob/d0ed487f487fa7ec814ff705ab55249983fe8707/OpenMath-Judging/review/entries/E17-CORE6.txt}\par}

\begin{samepage}
\subsection{E17-CORE 7  -  \texttt{KariInverseShrink.literal\_preimage\_correct}, \texttt{KariInverseShrink.no\_monotone\_inverse\_reset}, `KariInverseShrink.reset\_needs\_shrink}
{\small Contributors: Qichao Wang | Family: CERNY-PART-7\par}

\textbf{Recorded scope:} For the literal six-state preimage-mask system, no cardinality-nondecreasing inverse word takes any singleton to the full mask. Every inverse word reaching full has \ensuremath{\geq}1 decreasing step; a specific length-25 inverse word from singleton 4 reaches full with exactly one decrease.

\end{samepage}
\textbf{Source classification:} Known-mathematics formalization / formal-library novelty and usefulness require verification

\textbf{Prior comparison:} Known mathematics; only a genuinely new, faithful, useful, reusable event-window formalization may count. Exact prior-library overlap is not yet independently cleared.

\textbf{Formal/source status:} Only four manifests, not the 27MBsourceZIP. Compiler/axiom logs are author descriptions. No local Lean source or fresh fidelity check.

\textbf{Research or reuse implication:} Propose this exact scope and family once, subject to the stated source, formal, novelty and deadline checks.

{\small \url{https://github.com/alejandrozu/openmath-2026-judging/blob/d0ed487f487fa7ec814ff705ab55249983fe8707/OpenMath-Judging/review/entries/E17-CORE7.txt}\par}

\begin{samepage}
\subsection{E17-CORE 8  -  \texttt{KariProductBase.reset\_length\_lower}, \texttt{KariProductBase.to45\_length\_lower}, \texttt{KariProductBase.mono44\_length\_lower}, `KariProductBase}
{\small Contributors: Qichao Wang | Family: CERNY-PART-8\par}

\textbf{Recorded scope:} In that same literal base system, every singleton-to-full inverse word has length \ensuremath{\geq}25. The last decreasing edge is 45\ensuremath{\to}44 or 47\ensuremath{\to}60; a sole decrease is 47\ensuremath{\to}60. The stated prefix/monotone-suffix hypotheses at 45/44 force cost \ensuremath{\geq}27; phase witnesses attain costs 25/27.

\end{samepage}
\textbf{Source classification:} Known-mathematics formalization / formal-library novelty and usefulness require verification

\textbf{Prior comparison:} Known mathematics; only a genuinely new, faithful, useful, reusable event-window formalization may count. Exact prior-library overlap is not yet independently cleared.

\textbf{Formal/source status:} Only four manifests, not the 27MBsourceZIP. Compiler/axiom logs are author descriptions. No local Lean source or fresh fidelity check.

\textbf{Research or reuse implication:} Propose this exact scope and family once, subject to the stated source, formal, novelty and deadline checks.

{\small \url{https://github.com/alejandrozu/openmath-2026-judging/blob/d0ed487f487fa7ec814ff705ab55249983fe8707/OpenMath-Judging/review/entries/E17-CORE8.txt}\par}

\begin{samepage}
\subsection{E17-CORE 9  -  \texttt{ReflectionCoreCut.cut\_closed}, \texttt{ReflectionCoreCut.run\_in\_layers}, \texttt{ReflectionCoreCut.no\_avoid\_first\_eight}, `ReflectionCoreCut.wi}
{\small Contributors: Qichao Wang | Family: CERNY-PART-9\par}

\textbf{Recorded scope:} In the abstract integer-hole recurrence, checked layers close through eight steps; every word of length \ensuremath{\leq}8 has \texttt{SafeRun}. The specified eight-letter prefix ends with holes [0,1,3], which are unsafe.

\end{samepage}
\textbf{Source classification:} Known-mathematics formalization / formal-library novelty and usefulness require verification

\textbf{Prior comparison:} Known mathematics; only a genuinely new, faithful, useful, reusable event-window formalization may count. Exact prior-library overlap is not yet independently cleared.

\textbf{Formal/source status:} Only four manifests, not the 27MBsourceZIP. Compiler/axiom logs are author descriptions. No local Lean source or fresh fidelity check.

\textbf{Research or reuse implication:} Propose this exact scope and family once, subject to the stated source, formal, novelty and deadline checks.

{\small \url{https://github.com/alejandrozu/openmath-2026-judging/blob/d0ed487f487fa7ec814ff705ab55249983fe8707/OpenMath-Judging/review/entries/E17-CORE9.txt}\par}

\begin{samepage}
\subsection{E17-PAPER  -  Unformalized all-parameter families}
{\small Contributors: Qichao Wang | Family: CERNY\par}

\textbf{Recorded scope:} Anchor B\_r(n), layered Figure 3, Kari Cartesian products and reflection reset 5n\ensuremath{-}11; general concrete bridges/all-n arguments not fully formalized.

\end{samepage}
\textbf{Source classification:} Unresolved or missing-artifact record; nonacceptance does not prove nonnovelty

\textbf{Prior comparison:} Potential method-obstruction and lower-bound research; no improved general reset upper bound is claimed.

\textbf{Formal/source status:} Original paper/proofs/source absent locally; exact abstract components do not certify concrete parameterized automata.

\textbf{Research or reuse implication:} Open score 0. Keep nine scoped-core M2 candidates on hold, with explicit proposed deduplication, pending the actual source.

{\small \url{https://github.com/alejandrozu/openmath-2026-judging/blob/d0ed487f487fa7ec814ff705ab55249983fe8707/OpenMath-Judging/review/entries/E17-PAPER.txt}\par}

\begin{samepage}
\subsection{E57-K13  -  13 lines/47triangles}
{\small Contributors: Rohith Poola | Family: KOBON\par}

\textbf{Recorded scope:} K(13)\ensuremath{\geq}47; packet classification: ties published; reported published 47.

\end{samepage}
\textbf{Source classification:} Baseline/known-result candidate, or incomplete claim with prior overlap; see exact evidence

\textbf{Prior comparison:} No eligible original mathematical delta established in the available packet. This is not a finding that all underlying work lacks novelty.

\textbf{Formal/source status:} Independent integer homogeneous-intersection count matches submitted count; no duplicate lines. No Lean proof supplied; approved certificate trust boundary still needs chair decision.

\textbf{Research or reuse implication:} Propose this exact scope and family once, subject to the stated source, formal, novelty and deadline checks.

{\small \url{https://github.com/alejandrozu/openmath-2026-judging/blob/d0ed487f487fa7ec814ff705ab55249983fe8707/OpenMath-Judging/review/entries/E57-K13.txt}\par}

\begin{samepage}
\subsection{E57-K18  -  18 lines/93triangles}
{\small Contributors: Rohith Poola | Family: KOBON\par}

\textbf{Recorded scope:} K(18)\ensuremath{\geq}93; packet classification: ties published; reported published 93.

\end{samepage}
\textbf{Source classification:} Baseline/known-result candidate, or incomplete claim with prior overlap; see exact evidence

\textbf{Prior comparison:} No eligible original mathematical delta established in the available packet. This is not a finding that all underlying work lacks novelty.

\textbf{Formal/source status:} Independent integer homogeneous-intersection count matches submitted count; no duplicate lines. No Lean proof supplied; approved certificate trust boundary still needs chair decision.

\textbf{Research or reuse implication:} Propose this exact scope and family once, subject to the stated source, formal, novelty and deadline checks.

{\small \url{https://github.com/alejandrozu/openmath-2026-judging/blob/d0ed487f487fa7ec814ff705ab55249983fe8707/OpenMath-Judging/review/entries/E57-K18.txt}\par}

\begin{samepage}
\subsection{E57-K19  -  19 lines/107triangles}
{\small Contributors: Rohith Poola | Family: KOBON\par}

\textbf{Recorded scope:} K(19)\ensuremath{\geq}107; packet classification: ties published; reported published 107.

\end{samepage}
\textbf{Source classification:} Baseline/known-result candidate, or incomplete claim with prior overlap; see exact evidence

\textbf{Prior comparison:} No eligible original mathematical delta established in the available packet. This is not a finding that all underlying work lacks novelty.

\textbf{Formal/source status:} Independent integer homogeneous-intersection count matches submitted count; no duplicate lines. No Lean proof supplied; approved certificate trust boundary still needs chair decision.

\textbf{Research or reuse implication:} Propose this exact scope and family once, subject to the stated source, formal, novelty and deadline checks.

{\small \url{https://github.com/alejandrozu/openmath-2026-judging/blob/d0ed487f487fa7ec814ff705ab55249983fe8707/OpenMath-Judging/review/entries/E57-K19.txt}\par}

\begin{samepage}
\subsection{E57-K31  -  31 lines/290triangles}
{\small Contributors: Rohith Poola | Family: KOBON\par}

\textbf{Recorded scope:} K(31)\ensuremath{\geq}290; packet classification: below published (299); reported published 299.

\end{samepage}
\textbf{Source classification:} Baseline/known-result candidate, or incomplete claim with prior overlap; see exact evidence

\textbf{Prior comparison:} No eligible original mathematical delta established in the available packet. This is not a finding that all underlying work lacks novelty.

\textbf{Formal/source status:} Independent integer homogeneous-intersection count matches submitted count; no duplicate lines. No Lean proof supplied; approved certificate trust boundary still needs chair decision.

\textbf{Research or reuse implication:} Propose this exact scope and family once, subject to the stated source, formal, novelty and deadline checks.

{\small \url{https://github.com/alejandrozu/openmath-2026-judging/blob/d0ed487f487fa7ec814ff705ab55249983fe8707/OpenMath-Judging/review/entries/E57-K31.txt}\par}

\begin{samepage}
\subsection{E57-K35  -  35 lines/338triangles}
{\small Contributors: Rohith Poola | Family: KOBON\par}

\textbf{Recorded scope:} K(35)\ensuremath{\geq}338; packet classification: below published (385); reported published 385.

\end{samepage}
\textbf{Source classification:} Baseline/known-result candidate, or incomplete claim with prior overlap; see exact evidence

\textbf{Prior comparison:} No eligible original mathematical delta established in the available packet. This is not a finding that all underlying work lacks novelty.

\textbf{Formal/source status:} Independent integer homogeneous-intersection count matches submitted count; no duplicate lines. No Lean proof supplied; approved certificate trust boundary still needs chair decision.

\textbf{Research or reuse implication:} Propose this exact scope and family once, subject to the stated source, formal, novelty and deadline checks.

{\small \url{https://github.com/alejandrozu/openmath-2026-judging/blob/d0ed487f487fa7ec814ff705ab55249983fe8707/OpenMath-Judging/review/entries/E57-K35.txt}\par}

\begin{samepage}
\subsection{E57-K37  -  37 lines/431triangles}
{\small Contributors: Rohith Poola | Family: KOBON\par}

\textbf{Recorded scope:} K(37)\ensuremath{\geq}431; packet classification: ties published; reported published 431.

\end{samepage}
\textbf{Source classification:} Baseline/known-result candidate, or incomplete claim with prior overlap; see exact evidence

\textbf{Prior comparison:} No eligible original mathematical delta established in the available packet. This is not a finding that all underlying work lacks novelty.

\textbf{Formal/source status:} Independent integer homogeneous-intersection count matches submitted count; no duplicate lines. No Lean proof supplied; approved certificate trust boundary still needs chair decision.

\textbf{Research or reuse implication:} Propose this exact scope and family once, subject to the stated source, formal, novelty and deadline checks.

{\small \url{https://github.com/alejandrozu/openmath-2026-judging/blob/d0ed487f487fa7ec814ff705ab55249983fe8707/OpenMath-Judging/review/entries/E57-K37.txt}\par}

\begin{samepage}
\subsection{E57-K38  -  38 lines/450triangles}
{\small Contributors: Rohith Poola | Family: KOBON\par}

\textbf{Recorded scope:} K(38)\ensuremath{\geq}450; packet classification: ties published; reported published 450.

\end{samepage}
\textbf{Source classification:} Baseline/known-result candidate, or incomplete claim with prior overlap; see exact evidence

\textbf{Prior comparison:} No eligible original mathematical delta established in the available packet. This is not a finding that all underlying work lacks novelty.

\textbf{Formal/source status:} Independent integer homogeneous-intersection count matches submitted count; no duplicate lines. No Lean proof supplied; approved certificate trust boundary still needs chair decision.

\textbf{Research or reuse implication:} Propose this exact scope and family once, subject to the stated source, formal, novelty and deadline checks.

{\small \url{https://github.com/alejandrozu/openmath-2026-judging/blob/d0ed487f487fa7ec814ff705ab55249983fe8707/OpenMath-Judging/review/entries/E57-K38.txt}\par}

\begin{samepage}
\subsection{E57-K41  -  41 lines/533triangles}
{\small Contributors: Rohith Poola | Family: KOBON\par}

\textbf{Recorded scope:} K(41)\ensuremath{\geq}533; packet classification: ties published; reported published 533.

\end{samepage}
\textbf{Source classification:} Baseline/known-result candidate, or incomplete claim with prior overlap; see exact evidence

\textbf{Prior comparison:} No eligible original mathematical delta established in the available packet. This is not a finding that all underlying work lacks novelty.

\textbf{Formal/source status:} Independent integer homogeneous-intersection count matches submitted count; no duplicate lines. No Lean proof supplied; approved certificate trust boundary still needs chair decision.

\textbf{Research or reuse implication:} Propose this exact scope and family once, subject to the stated source, formal, novelty and deadline checks.

{\small \url{https://github.com/alejandrozu/openmath-2026-judging/blob/d0ed487f487fa7ec814ff705ab55249983fe8707/OpenMath-Judging/review/entries/E57-K41.txt}\par}

\begin{samepage}
\subsection{E57-K43  -  43 lines/587triangles}
{\small Contributors: Rohith Poola | Family: KOBON\par}

\textbf{Recorded scope:} K(43)\ensuremath{\geq}587; packet classification: ties published; reported published 587.

\end{samepage}
\textbf{Source classification:} Baseline/known-result candidate, or incomplete claim with prior overlap; see exact evidence

\textbf{Prior comparison:} No eligible original mathematical delta established in the available packet. This is not a finding that all underlying work lacks novelty.

\textbf{Formal/source status:} Independent integer homogeneous-intersection count matches submitted count; no duplicate lines. No Lean proof supplied; approved certificate trust boundary still needs chair decision.

\textbf{Research or reuse implication:} Propose this exact scope and family once, subject to the stated source, formal, novelty and deadline checks.

{\small \url{https://github.com/alejandrozu/openmath-2026-judging/blob/d0ed487f487fa7ec814ff705ab55249983fe8707/OpenMath-Judging/review/entries/E57-K43.txt}\par}

\begin{samepage}
\subsection{E57-K45  -  45 lines/645triangles}
{\small Contributors: Rohith Poola | Family: KOBON\par}

\textbf{Recorded scope:} K(45)\ensuremath{\geq}645; packet classification: ties published; reported published 645.

\end{samepage}
\textbf{Source classification:} Baseline/known-result candidate, or incomplete claim with prior overlap; see exact evidence

\textbf{Prior comparison:} No eligible original mathematical delta established in the available packet. This is not a finding that all underlying work lacks novelty.

\textbf{Formal/source status:} Independent integer homogeneous-intersection count matches submitted count; no duplicate lines. No Lean proof supplied; approved certificate trust boundary still needs chair decision.

\textbf{Research or reuse implication:} Propose this exact scope and family once, subject to the stated source, formal, novelty and deadline checks.

{\small \url{https://github.com/alejandrozu/openmath-2026-judging/blob/d0ed487f487fa7ec814ff705ab55249983fe8707/OpenMath-Judging/review/entries/E57-K45.txt}\par}

\begin{samepage}
\subsection{E57-K46  -  46 lines/667triangles}
{\small Contributors: Rohith Poola | Family: KOBON\par}

\textbf{Recorded scope:} K(46)\ensuremath{\geq}667; packet classification: ties published; reported published 667.

\end{samepage}
\textbf{Source classification:} Baseline/known-result candidate, or incomplete claim with prior overlap; see exact evidence

\textbf{Prior comparison:} No eligible original mathematical delta established in the available packet. This is not a finding that all underlying work lacks novelty.

\textbf{Formal/source status:} Independent integer homogeneous-intersection count matches submitted count; no duplicate lines. No Lean proof supplied; approved certificate trust boundary still needs chair decision.

\textbf{Research or reuse implication:} Propose this exact scope and family once, subject to the stated source, formal, novelty and deadline checks.

{\small \url{https://github.com/alejandrozu/openmath-2026-judging/blob/d0ed487f487fa7ec814ff705ab55249983fe8707/OpenMath-Judging/review/entries/E57-K46.txt}\par}

\begin{samepage}
\subsection{E57-K49  -  49 lines/767triangles}
{\small Contributors: Rohith Poola | Family: KOBON\par}

\textbf{Recorded scope:} K(49)\ensuremath{\geq}767; packet classification: ties published; reported published 767.

\end{samepage}
\textbf{Source classification:} Baseline/known-result candidate, or incomplete claim with prior overlap; see exact evidence

\textbf{Prior comparison:} No eligible original mathematical delta established in the available packet. This is not a finding that all underlying work lacks novelty.

\textbf{Formal/source status:} Independent integer homogeneous-intersection count matches submitted count; no duplicate lines. No Lean proof supplied; approved certificate trust boundary still needs chair decision.

\textbf{Research or reuse implication:} Propose this exact scope and family once, subject to the stated source, formal, novelty and deadline checks.

{\small \url{https://github.com/alejandrozu/openmath-2026-judging/blob/d0ed487f487fa7ec814ff705ab55249983fe8707/OpenMath-Judging/review/entries/E57-K49.txt}\par}

\begin{samepage}
\subsection{E57-K50  -  50 lines/792triangles}
{\small Contributors: Rohith Poola | Family: KOBON\par}

\textbf{Recorded scope:} K(50)\ensuremath{\geq}792; packet classification: ties published; reported published 792.

\end{samepage}
\textbf{Source classification:} Baseline/known-result candidate, or incomplete claim with prior overlap; see exact evidence

\textbf{Prior comparison:} No eligible original mathematical delta established in the available packet. This is not a finding that all underlying work lacks novelty.

\textbf{Formal/source status:} Independent integer homogeneous-intersection count matches submitted count; no duplicate lines. No Lean proof supplied; approved certificate trust boundary still needs chair decision.

\textbf{Research or reuse implication:} Propose this exact scope and family once, subject to the stated source, formal, novelty and deadline checks.

{\small \url{https://github.com/alejandrozu/openmath-2026-judging/blob/d0ed487f487fa7ec814ff705ab55249983fe8707/OpenMath-Judging/review/entries/E57-K50.txt}\par}

\begin{samepage}
\subsection{E57-K53  -  53 lines/901triangles}
{\small Contributors: Rohith Poola | Family: KOBON\par}

\textbf{Recorded scope:} K(53)\ensuremath{\geq}901; packet classification: ties published; reported published 901.

\end{samepage}
\textbf{Source classification:} Baseline/known-result candidate, or incomplete claim with prior overlap; see exact evidence

\textbf{Prior comparison:} No eligible original mathematical delta established in the available packet. This is not a finding that all underlying work lacks novelty.

\textbf{Formal/source status:} Independent integer homogeneous-intersection count matches submitted count; no duplicate lines. No Lean proof supplied; approved certificate trust boundary still needs chair decision.

\textbf{Research or reuse implication:} Propose this exact scope and family once, subject to the stated source, formal, novelty and deadline checks.

{\small \url{https://github.com/alejandrozu/openmath-2026-judging/blob/d0ed487f487fa7ec814ff705ab55249983fe8707/OpenMath-Judging/review/entries/E57-K53.txt}\par}

\begin{samepage}
\subsection{E57-K54  -  54 lines/927triangles}
{\small Contributors: Rohith Poola | Family: KOBON\par}

\textbf{Recorded scope:} K(54)\ensuremath{\geq}927; packet classification: ties published; reported published 927.

\end{samepage}
\textbf{Source classification:} Baseline/known-result candidate, or incomplete claim with prior overlap; see exact evidence

\textbf{Prior comparison:} No eligible original mathematical delta established in the available packet. This is not a finding that all underlying work lacks novelty.

\textbf{Formal/source status:} Independent integer homogeneous-intersection count matches submitted count; no duplicate lines. No Lean proof supplied; approved certificate trust boundary still needs chair decision.

\textbf{Research or reuse implication:} Propose this exact scope and family once, subject to the stated source, formal, novelty and deadline checks.

{\small \url{https://github.com/alejandrozu/openmath-2026-judging/blob/d0ed487f487fa7ec814ff705ab55249983fe8707/OpenMath-Judging/review/entries/E57-K54.txt}\par}

\begin{samepage}
\subsection{E57-K57  -  57 lines/1045triangles}
{\small Contributors: Rohith Poola | Family: KOBON\par}

\textbf{Recorded scope:} K(57)\ensuremath{\geq}1045; packet classification: ties published; reported published 1045.

\end{samepage}
\textbf{Source classification:} Baseline/known-result candidate, or incomplete claim with prior overlap; see exact evidence

\textbf{Prior comparison:} No eligible original mathematical delta established in the available packet. This is not a finding that all underlying work lacks novelty.

\textbf{Formal/source status:} Independent integer homogeneous-intersection count matches submitted count; no duplicate lines. No Lean proof supplied; approved certificate trust boundary still needs chair decision.

\textbf{Research or reuse implication:} Propose this exact scope and family once, subject to the stated source, formal, novelty and deadline checks.

{\small \url{https://github.com/alejandrozu/openmath-2026-judging/blob/d0ed487f487fa7ec814ff705ab55249983fe8707/OpenMath-Judging/review/entries/E57-K57.txt}\par}

\begin{samepage}
\subsection{E57-K65  -  65 lines/1365triangles}
{\small Contributors: Rohith Poola | Family: KOBON\par}

\textbf{Recorded scope:} K(65)\ensuremath{\geq}1365; packet classification: ties published; reported published 1365.

\end{samepage}
\textbf{Source classification:} Baseline/known-result candidate, or incomplete claim with prior overlap; see exact evidence

\textbf{Prior comparison:} No eligible original mathematical delta established in the available packet. This is not a finding that all underlying work lacks novelty.

\textbf{Formal/source status:} Independent integer homogeneous-intersection count matches submitted count; no duplicate lines. No Lean proof supplied; approved certificate trust boundary still needs chair decision.

\textbf{Research or reuse implication:} Propose this exact scope and family once, subject to the stated source, formal, novelty and deadline checks.

{\small \url{https://github.com/alejandrozu/openmath-2026-judging/blob/d0ed487f487fa7ec814ff705ab55249983fe8707/OpenMath-Judging/review/entries/E57-K65.txt}\par}

\begin{samepage}
\subsection{E57-K7  -  7 lines/11triangles}
{\small Contributors: Rohith Poola | Family: KOBON\par}

\textbf{Recorded scope:} K(7)\ensuremath{\geq}11; packet classification: ties published; reported published 11.

\end{samepage}
\textbf{Source classification:} Baseline/known-result candidate, or incomplete claim with prior overlap; see exact evidence

\textbf{Prior comparison:} No eligible original mathematical delta established in the available packet. This is not a finding that all underlying work lacks novelty.

\textbf{Formal/source status:} Independent integer homogeneous-intersection count matches submitted count; no duplicate lines. No Lean proof supplied; approved certificate trust boundary still needs chair decision.

\textbf{Research or reuse implication:} Propose this exact scope and family once, subject to the stated source, formal, novelty and deadline checks.

{\small \url{https://github.com/alejandrozu/openmath-2026-judging/blob/d0ed487f487fa7ec814ff705ab55249983fe8707/OpenMath-Judging/review/entries/E57-K7.txt}\par}

\begin{samepage}
\subsection{E57-K73  -  73 lines/1727triangles}
{\small Contributors: Rohith Poola | Family: KOBON\par}

\textbf{Recorded scope:} K(73)\ensuremath{\geq}1727; packet classification: ties published; reported published 1727.

\end{samepage}
\textbf{Source classification:} Baseline/known-result candidate, or incomplete claim with prior overlap; see exact evidence

\textbf{Prior comparison:} No eligible original mathematical delta established in the available packet. This is not a finding that all underlying work lacks novelty.

\textbf{Formal/source status:} Independent integer homogeneous-intersection count matches submitted count; no duplicate lines. No Lean proof supplied; approved certificate trust boundary still needs chair decision.

\textbf{Research or reuse implication:} Propose this exact scope and family once, subject to the stated source, formal, novelty and deadline checks.

{\small \url{https://github.com/alejandrozu/openmath-2026-judging/blob/d0ed487f487fa7ec814ff705ab55249983fe8707/OpenMath-Judging/review/entries/E57-K73.txt}\par}

\begin{samepage}
\subsection{E57-K97  -  97 lines/3071triangles}
{\small Contributors: Rohith Poola | Family: KOBON\par}

\textbf{Recorded scope:} K(97)\ensuremath{\geq}3071; packet classification: ties published; reported published 3071.

\end{samepage}
\textbf{Source classification:} Baseline/known-result candidate, or incomplete claim with prior overlap; see exact evidence

\textbf{Prior comparison:} No eligible original mathematical delta established in the available packet. This is not a finding that all underlying work lacks novelty.

\textbf{Formal/source status:} Independent integer homogeneous-intersection count matches submitted count; no duplicate lines. No Lean proof supplied; approved certificate trust boundary still needs chair decision.

\textbf{Research or reuse implication:} Propose this exact scope and family once, subject to the stated source, formal, novelty and deadline checks.

{\small \url{https://github.com/alejandrozu/openmath-2026-judging/blob/d0ed487f487fa7ec814ff705ab55249983fe8707/OpenMath-Judging/review/entries/E57-K97.txt}\par}

\begin{samepage}
\subsection{E57-R4  -  Generalized Blanc conjecture / pseudolines}
{\small Contributors: Rohith Poola | Family: KOBON\par}

\textbf{Recorded scope:} Generalized bound conjecture unproved;SATpseudolines481/475-480 not realized as straight lines.

\end{samepage}
\textbf{Source classification:} Unresolved or missing-artifact record; nonacceptance does not prove nonnovelty

\textbf{Prior comparison:} No eligible original mathematical delta established in the available packet. This is not a finding that all underlying work lacks novelty.

\textbf{Formal/source status:} No formal theorem.

\textbf{Research or reuse implication:} Propose this exact scope and family once, subject to the stated source, formal, novelty and deadline checks.

{\small \url{https://github.com/alejandrozu/openmath-2026-judging/blob/d0ed487f487fa7ec814ff705ab55249983fe8707/OpenMath-Judging/review/entries/E57-R4.txt}\par}

\begin{samepage}
\subsection{E65-DYADIC  -  Exact hill/extremal equivalence}
{\small Contributors: Wilson Wu | Family: M2-E3-DYADIC\par}

\textbf{Recorded scope:} Original hill iff \ensuremath{\Sigma}j r\_k(2\textasciicircum{}j)/2\textasciicircum{}j converges for everyk\ensuremath{\geq}4; neither side proved.

\end{samepage}
\textbf{Source classification:} Known-mathematics formalization / formal-library novelty and usefulness require verification

\textbf{Prior comparison:} Candidate useful formal bridge, overlapping Leon/Cameron's family. Supplied pinned-image checks and the separate equivalence hill are evidence for this theorem, not the original hill.

\textbf{Formal/source status:} The completed exact recovered-source Windows checkpoint at 07:05:52 UTC on 5 October 2026 checked 187 dependency-library modules: 183 granular and four whole-resource. Six remaining FC dependency modules and the six FC-dependent entrant modules, including the dyadic endpoint, are still pending; no current FC-dependent organizer audit or selected output has run. The separate earlier exact Std one-source/14-print witness is accepted for possible later evidence composition. The final 200-body/56-output composite is not qualified. Historical author logs and the answer(sorry) statement-wrapper caveat are retained. No full-project PASS or originality score is conferred.

\textbf{Research or reuse implication:} One proposed M2 equivalence family after independent provenance and scoped formal-novelty review; no original Erdős 3 points.

{\small \url{https://github.com/alejandrozu/openmath-2026-judging/blob/d0ed487f487fa7ec814ff705ab55249983fe8707/OpenMath-Judging/review/entries/E65-DYADIC.txt}\par}

\begin{samepage}
\subsection{E65-THREE  -  Conditional three-term analytic bridges}
{\small Contributors: Wilson Wu | Family: M2-E3-THREE\par}

\textbf{Recorded scope:} Kelley-Meka/Bloom-Sisask density hypotheses imply known three-term target; analytic bounds are not proved here.

\end{samepage}
\textbf{Source classification:} Known-mathematics formalization / formal-library novelty and usefulness require verification

\textbf{Prior comparison:} Known analytic mathematics with a potentially new useful bridge; count only if distinct from the dyadic/AP library family.

\textbf{Formal/source status:} Parameters are explicit assumptions; clean axiom output does not discharge them. The independently reviewed completed checkpoint at 07:05:52 UTC on 5 October 2026 contains 187 recovered FC dependency source successes (183 granular plus four whole), with the six FC-dependent entrant modules and six current audits/42 selected outputs still pending. It does not certify those unrun conditional endpoints or discharge their analytic premises; the prior exact Std one-source/14-print witness remains separate.

\textbf{Research or reuse implication:} Supporting/conditional row; conservative M2 union includes it with the AP interface, not a standalone full 3-term proof.

{\small \url{https://github.com/alejandrozu/openmath-2026-judging/blob/d0ed487f487fa7ec814ff705ab55249983fe8707/OpenMath-Judging/review/entries/E65-THREE.txt}\par}

\begin{samepage}
\subsection{E65-ELEMENTARY  -  Elementary AP vocabulary / special cases}
{\small Contributors: Wilson Wu | Family: M2-AP-CONCLUSION\par}

\textbf{Recorded scope:} Cofinite sets, affine tails, unboundedness\ensuremath{\to}two-term, AP truncation and library bridges.

\end{samepage}
\textbf{Source classification:} Known-mathematics formalization / formal-library novelty and usefulness require verification

\textbf{Prior comparison:} Known mathematics; only a genuinely new, faithful, useful, reusable event-window formalization may count. Exact prior-library overlap is not yet independently cleared.

\begin{samepage}
\textbf{Formal/source status:} Actual independent exact-source Std-only replay on 5 October 2026: Erdos3SpecialCases source 1/1 and its selected audit 14/14 PASS; five axiom-free, nine standard-kernel-only, no selected sorry/native/unrecognized axiom. ContainsAP has a natural start and positive natural step; ArbitrarilyLongAPs quantifies all finite natural lengths. The recovered-source Windows replay separately reached a completed checkpoint at 07:05:52 UTC on 5 October 2026 with 187 FC dependency-library modules checked (183 granular plus four whole). Six further FC modules and all six FC-dependent entrant modules, including the organizer Set/cardinality/Filter bridge, remain pending; zero of their six current audits/42 selected outputs have run. The exact prior Std one-source/14-print witness is separately accepted, but a 200-body/56-output final composition is not qualified. This is partial known-mathematics verification, with no full-project PASS, original-container execution or new reciprocal-divergence result.

\end{samepage}
\textbf{Research or reuse implication:} Propose this exact scope and family once, subject to the stated source, formal, novelty and deadline checks.

{\small \url{https://github.com/alejandrozu/openmath-2026-judging/blob/d0ed487f487fa7ec814ff705ab55249983fe8707/OpenMath-Judging/review/entries/E65-ELEMENTARY.txt}\par}

\begin{samepage}
\subsection{E34-HILL 1  -  alejandrozu/matrix-multiplication-tensor-3x3}
{\small Contributors:  | Family: MATRIX-3\par}

\textbf{Recorded scope:} Rank 23, support 139 on September 21: matches the old support baseline; pre-event provenance needs confirmation. No final packet identified.

\end{samepage}
\textbf{Source classification:} Baseline/known-result candidate, or incomplete claim with prior overlap; see exact evidence

\textbf{Prior comparison:} No eligible original mathematical delta established in the available packet. This is not a finding that all underlying work lacks novelty.

\textbf{Formal/source status:} Only board metadata/access leads; no original mathematical proof or exact final packet. Observed run predates event; later eligible delta not established.

\textbf{Research or reuse implication:} Hold missing payload; current proposed points 0. Board metadata cannot establish an eligible exact theorem.

{\small \url{https://github.com/alejandrozu/openmath-2026-judging/blob/d0ed487f487fa7ec814ff705ab55249983fe8707/OpenMath-Judging/review/entries/E34-HILL1.txt}\par}

\begin{samepage}
\subsection{E35-HILL 1  -  alejandrozu/grothendieck-constant-witnesses}
{\small Contributors:  | Family: GROTHENDIECK\par}

\textbf{Recorded scope:} Grothendieck gap 1.414213 with area 4 and 80-bit certificate: compatible with the known two-by-two CHSH witness, not evidence of improving the best general lower bound.

\end{samepage}
\textbf{Source classification:} Baseline/known-result candidate, or incomplete claim with prior overlap; see exact evidence

\textbf{Prior comparison:} No eligible original mathematical delta established in the available packet. This is not a finding that all underlying work lacks novelty.

\textbf{Formal/source status:} Only board metadata/access leads; no original mathematical proof or exact final packet. 

\textbf{Research or reuse implication:} Hold missing payload; current proposed points 0. Board metadata cannot establish an eligible exact theorem.

{\small \url{https://github.com/alejandrozu/openmath-2026-judging/blob/d0ed487f487fa7ec814ff705ab55249983fe8707/OpenMath-Judging/review/entries/E35-HILL1.txt}\par}

\begin{samepage}
\subsection{E36-HILL 1  -  alejandrozu/matrix-multiplication-tensor-3x3}
{\small Contributors:  | Family: MATRIX-3\par}

\textbf{Recorded scope:} Rank 23, support 139: baseline-matching scored artifact. GitHub name Advaith Appajodu and Gopal registration associate this person with Chandragupt's group; the final entrant roster remains unconfirmed.

\end{samepage}
\textbf{Source classification:} Baseline/known-result candidate, or incomplete claim with prior overlap; see exact evidence

\textbf{Prior comparison:} No eligible original mathematical delta established in the available packet. This is not a finding that all underlying work lacks novelty.

\textbf{Formal/source status:} Only board metadata/access leads; no original mathematical proof or exact final packet. 

\textbf{Research or reuse implication:} Hold missing payload; current proposed points 0. Board metadata cannot establish an eligible exact theorem.

{\small \url{https://github.com/alejandrozu/openmath-2026-judging/blob/d0ed487f487fa7ec814ff705ab55249983fe8707/OpenMath-Judging/review/entries/E36-HILL1.txt}\par}

\begin{samepage}
\subsection{E37-HILL 1  -  alejandrozu/kobon-triangles}
{\small Contributors:  | Family: KOBON\par}

\textbf{Recorded scope:} Kobon n=18: 90 triangles, below the known 93 construction; scored September 21 before the event.

\end{samepage}
\textbf{Source classification:} Baseline/known-result candidate, or incomplete claim with prior overlap; see exact evidence

\textbf{Prior comparison:} No eligible original mathematical delta established in the available packet. This is not a finding that all underlying work lacks novelty.

\textbf{Formal/source status:} Only board metadata/access leads; no original mathematical proof or exact final packet. Observed run predates event; later eligible delta not established.

\textbf{Research or reuse implication:} Hold missing payload; current proposed points 0. Board metadata cannot establish an eligible exact theorem.

{\small \url{https://github.com/alejandrozu/openmath-2026-judging/blob/d0ed487f487fa7ec814ff705ab55249983fe8707/OpenMath-Judging/review/entries/E37-HILL1.txt}\par}

\begin{samepage}
\subsection{E38-HILL 1  -  alejandrozu/kobon-triangles}
{\small Contributors:  | Family: KOBON\par}

\textbf{Recorded scope:} Kobon n=18: 93, tying the known lower bound; September 21 pre-event score. No inspectable final packet identified.

\end{samepage}
\textbf{Source classification:} Baseline/known-result candidate, or incomplete claim with prior overlap; see exact evidence

\textbf{Prior comparison:} No eligible original mathematical delta established in the available packet. This is not a finding that all underlying work lacks novelty.

\textbf{Formal/source status:} Only board metadata/access leads; no original mathematical proof or exact final packet. Observed run predates event; later eligible delta not established.

\textbf{Research or reuse implication:} Hold missing payload; current proposed points 0. Board metadata cannot establish an eligible exact theorem.

{\small \url{https://github.com/alejandrozu/openmath-2026-judging/blob/d0ed487f487fa7ec814ff705ab55249983fe8707/OpenMath-Judging/review/entries/E38-HILL1.txt}\par}

\begin{samepage}
\subsection{E39-HILL 1  -  alejandrozu/kobon-triangles}
{\small Contributors:  | Family: KOBON\par}

\textbf{Recorded scope:} Kobon n=18: 93, a valid benchmark tie rather than the missing 94 arrangement or an upper-bound proof.

\end{samepage}
\textbf{Source classification:} Baseline/known-result candidate, or incomplete claim with prior overlap; see exact evidence

\textbf{Prior comparison:} No eligible original mathematical delta established in the available packet. This is not a finding that all underlying work lacks novelty.

\textbf{Formal/source status:} Only board metadata/access leads; no original mathematical proof or exact final packet. 

\textbf{Research or reuse implication:} Hold missing payload; current proposed points 0. Board metadata cannot establish an eligible exact theorem.

{\small \url{https://github.com/alejandrozu/openmath-2026-judging/blob/d0ed487f487fa7ec814ff705ab55249983fe8707/OpenMath-Judging/review/entries/E39-HILL1.txt}\par}

\begin{samepage}
\subsection{E40-HILL 1  -  alejandrozu/kobon-triangles}
{\small Contributors:  | Family: KOBON\par}

\textbf{Recorded scope:} Kobon n=18: 93, baseline tie on September 20. No current final packet identified.

\end{samepage}
\textbf{Source classification:} Baseline/known-result candidate, or incomplete claim with prior overlap; see exact evidence

\textbf{Prior comparison:} No eligible original mathematical delta established in the available packet. This is not a finding that all underlying work lacks novelty.

\textbf{Formal/source status:} Only board metadata/access leads; no original mathematical proof or exact final packet. Observed run predates event; later eligible delta not established.

\textbf{Research or reuse implication:} Hold missing payload; current proposed points 0. Board metadata cannot establish an eligible exact theorem.

{\small \url{https://github.com/alejandrozu/openmath-2026-judging/blob/d0ed487f487fa7ec814ff705ab55249983fe8707/OpenMath-Judging/review/entries/E40-HILL1.txt}\par}

\begin{samepage}
\subsection{E41-HILL 1  -  alejandrozu/busy-beaver-6-certificates}
{\small Contributors:  | Family: BB6\par}

\textbf{Recorded scope:} BB6 validation: 178,727 steps, 180 ones, span 539. Scored candidate, substantially below the validation leader. No machine/trace packet identified.

\end{samepage}
\textbf{Source classification:} Baseline/known-result candidate, or incomplete claim with prior overlap; see exact evidence

\textbf{Prior comparison:} No eligible original mathematical delta established in the available packet. This is not a finding that all underlying work lacks novelty.

\textbf{Formal/source status:} Only board metadata/access leads; no original mathematical proof or exact final packet. 

\textbf{Research or reuse implication:} Hold missing payload; current proposed points 0. Board metadata cannot establish an eligible exact theorem.

{\small \url{https://github.com/alejandrozu/openmath-2026-judging/blob/d0ed487f487fa7ec814ff705ab55249983fe8707/OpenMath-Judging/review/entries/E41-HILL1.txt}\par}

\begin{samepage}
\subsection{E42-HILL 1  -  alejandrozu/matrix-multiplication-tensor-3x3}
{\small Contributors:  | Family: MATRIX-3\par}

\textbf{Recorded scope:} Matrix rank 23, support 153 on September 20: worse than support 139 and the current 138 candidate.

\end{samepage}
\textbf{Source classification:} Unresolved or missing-artifact record; nonacceptance does not prove nonnovelty

\textbf{Prior comparison:} No eligible original mathematical delta established in the available packet. This is not a finding that all underlying work lacks novelty.

\textbf{Formal/source status:} Only board metadata/access leads; no original mathematical proof or exact final packet. Observed run predates event; later eligible delta not established.

\textbf{Research or reuse implication:} Hold missing payload; current proposed points 0. Board metadata cannot establish an eligible exact theorem.

{\small \url{https://github.com/alejandrozu/openmath-2026-judging/blob/d0ed487f487fa7ec814ff705ab55249983fe8707/OpenMath-Judging/review/entries/E42-HILL1.txt}\par}

\begin{samepage}
\subsection{E43-HILL 1  -  alejandrozu/clique-cluster-ramsey-multiplicity}
{\small Contributors:  | Family: RAMSEY-K4\par}

\textbf{Recorded scope:} Ramsey validation density 0.030142036534 beats the frozen 0.03014227343194 reference. Exact graph, rational weights and blow-up proof still absent. No Ramsey repository among the public repositories checked.

\end{samepage}
\textbf{Source classification:} Baseline/known-result candidate, or incomplete claim with prior overlap; see exact evidence

\textbf{Prior comparison:} Candidate scoped mathematical delta; exact baseline and independently sealed provenance are required. A missing gallery value is not evidence of a world record.

\textbf{Formal/source status:} Only board metadata/access leads; no original mathematical proof or exact final packet. 

\textbf{Research or reuse implication:} Hold missing payload; current proposed points 0. Board metadata cannot establish an eligible exact theorem.

{\small \url{https://github.com/alejandrozu/openmath-2026-judging/blob/d0ed487f487fa7ec814ff705ab55249983fe8707/OpenMath-Judging/review/entries/E43-HILL1.txt}\par}

\begin{samepage}
\subsection{E44-HILL 1  -  alejandrozu/grothendieck-constant-witnesses}
{\small Contributors:  | Family: GROTHENDIECK\par}

\textbf{Recorded scope:} Grothendieck gap 1.414213, area 4, 80 bits: known small witness scale; no new general lower-bound advance established.

\end{samepage}
\textbf{Source classification:} Baseline/known-result candidate, or incomplete claim with prior overlap; see exact evidence

\textbf{Prior comparison:} No eligible original mathematical delta established in the available packet. This is not a finding that all underlying work lacks novelty.

\textbf{Formal/source status:} Only board metadata/access leads; no original mathematical proof or exact final packet. 

\textbf{Research or reuse implication:} Hold missing payload; current proposed points 0. Board metadata cannot establish an eligible exact theorem.

{\small \url{https://github.com/alejandrozu/openmath-2026-judging/blob/d0ed487f487fa7ec814ff705ab55249983fe8707/OpenMath-Judging/review/entries/E44-HILL1.txt}\par}

\begin{samepage}
\subsection{E45-HILL 1  -  alejandrozu/busy-beaver-6-certificates}
{\small Contributors:  | Family: BB6\par}

\textbf{Recorded scope:} BB6 validation: 232,746 steps, 554 ones, span 684. Finite certified-run benchmark candidate; no final mathematical packet identified.

\end{samepage}
\textbf{Source classification:} Unresolved or missing-artifact record; nonacceptance does not prove nonnovelty

\textbf{Prior comparison:} No eligible original mathematical delta established in the available packet. This is not a finding that all underlying work lacks novelty.

\textbf{Formal/source status:} Only board metadata/access leads; no original mathematical proof or exact final packet. 

\textbf{Research or reuse implication:} Hold missing payload; current proposed points 0. Board metadata cannot establish an eligible exact theorem.

{\small \url{https://github.com/alejandrozu/openmath-2026-judging/blob/d0ed487f487fa7ec814ff705ab55249983fe8707/OpenMath-Judging/review/entries/E45-HILL1.txt}\par}

\begin{samepage}
\subsection{E46-HILL 1  -  alejandrozu/busy-beaver-6-certificates}
{\small Contributors:  | Family: BB6\par}

\textbf{Recorded scope:} BB6 validation: 249,881 steps, 554 ones, span 735 on September 21; ties Hamdi/n0rang 2 numerically but precedes event. Same metrics alone do not prove the same machine.

\end{samepage}
\textbf{Source classification:} Baseline/known-result candidate, or incomplete claim with prior overlap; see exact evidence

\textbf{Prior comparison:} No eligible original mathematical delta established in the available packet. This is not a finding that all underlying work lacks novelty.

\textbf{Formal/source status:} Only board metadata/access leads; no original mathematical proof or exact final packet. Observed run predates event; later eligible delta not established.

\textbf{Research or reuse implication:} Hold missing payload; current proposed points 0. Board metadata cannot establish an eligible exact theorem.

{\small \url{https://github.com/alejandrozu/openmath-2026-judging/blob/d0ed487f487fa7ec814ff705ab55249983fe8707/OpenMath-Judging/review/entries/E46-HILL1.txt}\par}

\begin{samepage}
\subsection{E48-HILL 1  -  alejandrozu/kobon-triangles}
{\small Contributors:  | Family: KOBON\par}

\textbf{Recorded scope:} Kobon n=18: 93 on September 21, known construction tie. No final packet identified.

\end{samepage}
\textbf{Source classification:} Baseline/known-result candidate, or incomplete claim with prior overlap; see exact evidence

\textbf{Prior comparison:} No eligible original mathematical delta established in the available packet. This is not a finding that all underlying work lacks novelty.

\textbf{Formal/source status:} Only board metadata/access leads; no original mathematical proof or exact final packet. Observed run predates event; later eligible delta not established.

\textbf{Research or reuse implication:} Hold missing payload; current proposed points 0. Board metadata cannot establish an eligible exact theorem.

{\small \url{https://github.com/alejandrozu/openmath-2026-judging/blob/d0ed487f487fa7ec814ff705ab55249983fe8707/OpenMath-Judging/review/entries/E48-HILL1.txt}\par}

\begin{samepage}
\subsection{E49-HILL 1  -  alejandrozu/kobon-triangles}
{\small Contributors:  | Family: KOBON\par}

\textbf{Recorded scope:} Kobon n=18: 93 on September 28, ties the known construction; exact final arrangement/paper not identified.

\end{samepage}
\textbf{Source classification:} Baseline/known-result candidate, or incomplete claim with prior overlap; see exact evidence

\textbf{Prior comparison:} No eligible original mathematical delta established in the available packet. This is not a finding that all underlying work lacks novelty.

\textbf{Formal/source status:} Only board metadata/access leads; no original mathematical proof or exact final packet. 

\textbf{Research or reuse implication:} Hold missing payload; current proposed points 0. Board metadata cannot establish an eligible exact theorem.

{\small \url{https://github.com/alejandrozu/openmath-2026-judging/blob/d0ed487f487fa7ec814ff705ab55249983fe8707/OpenMath-Judging/review/entries/E49-HILL1.txt}\par}

\begin{samepage}
\subsection{E50-HILL 1  -  alejandrozu/busy-beaver-6-certificates}
{\small Contributors:  | Family: BB6\par}

\textbf{Recorded scope:} BB6 validation: 249,614 steps, 543 ones, span 973. Near the bounded validation leader, but no global BB6 advance or complete packet established.

\end{samepage}
\textbf{Source classification:} Unresolved or missing-artifact record; nonacceptance does not prove nonnovelty

\textbf{Prior comparison:} No eligible original mathematical delta established in the available packet. This is not a finding that all underlying work lacks novelty.

\textbf{Formal/source status:} Only board metadata/access leads; no original mathematical proof or exact final packet. 

\textbf{Research or reuse implication:} Hold missing payload; current proposed points 0. Board metadata cannot establish an eligible exact theorem.

{\small \url{https://github.com/alejandrozu/openmath-2026-judging/blob/d0ed487f487fa7ec814ff705ab55249983fe8707/OpenMath-Judging/review/entries/E50-HILL1.txt}\par}

\begin{samepage}
\subsection{E52-HILL 1  -  alejandrozu/kobon-triangles}
{\small Contributors:  | Family: KOBON\par}

\textbf{Recorded scope:} Kobon n=18: 93, known lower-bound tie. AutoLab links GitHub kevinsrun (display name Icer). Likely Kevin Srun registration overlap, not conclusively confirmed.

\end{samepage}
\textbf{Source classification:} Baseline/known-result candidate, or incomplete claim with prior overlap; see exact evidence

\textbf{Prior comparison:} No eligible original mathematical delta established in the available packet. This is not a finding that all underlying work lacks novelty.

\textbf{Formal/source status:} Only board metadata/access leads; no original mathematical proof or exact final packet. 

\textbf{Research or reuse implication:} Hold missing payload; current proposed points 0. Board metadata cannot establish an eligible exact theorem.

{\small \url{https://github.com/alejandrozu/openmath-2026-judging/blob/d0ed487f487fa7ec814ff705ab55249983fe8707/OpenMath-Judging/review/entries/E52-HILL1.txt}\par}

\begin{samepage}
\subsection{E53-HILL 1  -  alejandrozu/kobon-triangles}
{\small Contributors:  | Family: KOBON\par}

\textbf{Recorded scope:} Kobon n=18: 91 on September 20, below 93. GitHub names Luiz Felipe Barbosa; no final event packet identified.

\end{samepage}
\textbf{Source classification:} Baseline/known-result candidate, or incomplete claim with prior overlap; see exact evidence

\textbf{Prior comparison:} No eligible original mathematical delta established in the available packet. This is not a finding that all underlying work lacks novelty.

\textbf{Formal/source status:} Only board metadata/access leads; no original mathematical proof or exact final packet. Observed run predates event; later eligible delta not established.

\textbf{Research or reuse implication:} Hold missing payload; current proposed points 0. Board metadata cannot establish an eligible exact theorem.

{\small \url{https://github.com/alejandrozu/openmath-2026-judging/blob/d0ed487f487fa7ec814ff705ab55249983fe8707/OpenMath-Judging/review/entries/E53-HILL1.txt}\par}

\begin{samepage}
\subsection{E55-HILL 1  -  alejandrozu/kobon-triangles}
{\small Contributors:  | Family: KOBON\par}

\textbf{Recorded scope:} Kobon n=39: 468, behind current 471 and Rohith 470. A larger-n exact construction can be useful; no proof of optimality or final packet identified.

\end{samepage}
\textbf{Source classification:} Baseline/known-result candidate, or incomplete claim with prior overlap; see exact evidence

\textbf{Prior comparison:} Candidate scoped mathematical delta; exact baseline and independently sealed provenance are required. A missing gallery value is not evidence of a world record.

\textbf{Formal/source status:} Only board metadata/access leads; no original mathematical proof or exact final packet. 

\textbf{Research or reuse implication:} Hold missing payload; current proposed points 0. Board metadata cannot establish an eligible exact theorem.

{\small \url{https://github.com/alejandrozu/openmath-2026-judging/blob/d0ed487f487fa7ec814ff705ab55249983fe8707/OpenMath-Judging/review/entries/E55-HILL1.txt}\par}

\begin{samepage}
\subsection{E58-HILL 1  -  alejandrozu/busy-beaver-6-certificates}
{\small Contributors:  | Family: BB6\par}

\textbf{Recorded scope:} BB6 validation 228641 steps, 647 ones, span 864.

\end{samepage}
\textbf{Source classification:} Unresolved or missing-artifact record; nonacceptance does not prove nonnovelty

\textbf{Prior comparison:} No eligible original mathematical delta established in the available packet. This is not a finding that all underlying work lacks novelty.

\textbf{Formal/source status:} Only board metadata/access leads; no original mathematical proof or exact final packet. 

\textbf{Research or reuse implication:} Hold missing payload; current proposed points 0. Board metadata cannot establish an eligible exact theorem.

{\small \url{https://github.com/alejandrozu/openmath-2026-judging/blob/d0ed487f487fa7ec814ff705ab55249983fe8707/OpenMath-Judging/review/entries/E58-HILL1.txt}\par}

\begin{samepage}
\subsection{E61-HILL 1  -  alejandrozu/kobon-triangles}
{\small Contributors:  | Family: KOBON\par}

\textbf{Recorded scope:} Kobon n=18: 93 tie; Grothendieck gap 1.414213, area 4, 80 bits. Both are baseline-level benchmark results.

\end{samepage}
\textbf{Source classification:} Baseline/known-result candidate, or incomplete claim with prior overlap; see exact evidence

\textbf{Prior comparison:} No eligible original mathematical delta established in the available packet. This is not a finding that all underlying work lacks novelty.

\textbf{Formal/source status:} Only board metadata/access leads; no original mathematical proof or exact final packet. 

\textbf{Research or reuse implication:} Hold missing payload; current proposed points 0. Board metadata cannot establish an eligible exact theorem.

{\small \url{https://github.com/alejandrozu/openmath-2026-judging/blob/d0ed487f487fa7ec814ff705ab55249983fe8707/OpenMath-Judging/review/entries/E61-HILL1.txt}\par}

\begin{samepage}
\subsection{E61-HILL 2  -  alejandrozu/grothendieck-constant-witnesses}
{\small Contributors:  | Family: GROTHENDIECK\par}

\textbf{Recorded scope:} Kobon n=18: 93 tie; Grothendieck gap 1.414213, area 4, 80 bits. Both are baseline-level benchmark results.

\end{samepage}
\textbf{Source classification:} Baseline/known-result candidate, or incomplete claim with prior overlap; see exact evidence

\textbf{Prior comparison:} No eligible original mathematical delta established in the available packet. This is not a finding that all underlying work lacks novelty.

\textbf{Formal/source status:} Only board metadata/access leads; no original mathematical proof or exact final packet. 

\textbf{Research or reuse implication:} Hold missing payload; current proposed points 0. Board metadata cannot establish an eligible exact theorem.

{\small \url{https://github.com/alejandrozu/openmath-2026-judging/blob/d0ed487f487fa7ec814ff705ab55249983fe8707/OpenMath-Judging/review/entries/E61-HILL2.txt}\par}

\begin{samepage}
\subsection{E63-HILL 1  -  alejandrozu/clique-cluster-ramsey-multiplicity}
{\small Contributors:  | Family: RAMSEY-K4\par}

\textbf{Recorded scope:} BB6 held-out 572171 steps/1035 ones/span 1380; validation is a different 233892-step score. Earlier AImerican mathematical workflow repo (September 7) is not a final event packet.

\end{samepage}
\textbf{Source classification:} Unresolved or missing-artifact record; nonacceptance does not prove nonnovelty

\textbf{Prior comparison:} No eligible original mathematical delta established in the available packet. This is not a finding that all underlying work lacks novelty.

\textbf{Formal/source status:} Only board metadata/access leads; no original mathematical proof or exact final packet. 

\textbf{Research or reuse implication:} Hold missing payload; current proposed points 0. Board metadata cannot establish an eligible exact theorem.

{\small \url{https://github.com/alejandrozu/openmath-2026-judging/blob/d0ed487f487fa7ec814ff705ab55249983fe8707/OpenMath-Judging/review/entries/E63-HILL1.txt}\par}

\begin{samepage}
\subsection{E63-HILL 2  -  alejandrozu/busy-beaver-6-certificates}
{\small Contributors:  | Family: BB6\par}

\textbf{Recorded scope:} BB6 held-out 572171 steps/1035 ones/span 1380; validation is a different 233892-step score. Earlier AImerican mathematical workflow repo (September 7) is not a final event packet.

\end{samepage}
\textbf{Source classification:} Unresolved or missing-artifact record; nonacceptance does not prove nonnovelty

\textbf{Prior comparison:} No eligible original mathematical delta established in the available packet. This is not a finding that all underlying work lacks novelty.

\textbf{Formal/source status:} Only board metadata/access leads; no original mathematical proof or exact final packet. 

\textbf{Research or reuse implication:} Hold missing payload; current proposed points 0. Board metadata cannot establish an eligible exact theorem.

{\small \url{https://github.com/alejandrozu/openmath-2026-judging/blob/d0ed487f487fa7ec814ff705ab55249983fe8707/OpenMath-Judging/review/entries/E63-HILL2.txt}\par}

\begin{samepage}
\subsection{E64-HILL 1  -  alejandrozu/kobon-triangles}
{\small Contributors:  | Family: KOBON\par}

\textbf{Recorded scope:} Kobon n=18: 93. Possibly Will's registration (email willow.pine.2011), but blank GitHub identity/no recent public repository does not confirm it.

\end{samepage}
\textbf{Source classification:} Unresolved or missing-artifact record; nonacceptance does not prove nonnovelty

\textbf{Prior comparison:} No eligible original mathematical delta established in the available packet. This is not a finding that all underlying work lacks novelty.

\textbf{Formal/source status:} Only board metadata/access leads; no original mathematical proof or exact final packet. 

\textbf{Research or reuse implication:} Hold missing payload; current proposed points 0. Board metadata cannot establish an eligible exact theorem.

{\small \url{https://github.com/alejandrozu/openmath-2026-judging/blob/d0ed487f487fa7ec814ff705ab55249983fe8707/OpenMath-Judging/review/entries/E64-HILL1.txt}\par}

\begin{samepage}
\subsection{E66-HILL 1  -  alejandrozu/clique-cluster-ramsey-multiplicity}
{\small Contributors:  | Family: RAMSEY-K4\par}

\textbf{Recorded scope:} Ramsey 0.030141721098 beats frozen reference, but dated September 21. Collatz 100\% weighted target coverage with three rules and BB6 246,872 steps are also pre-event. Identity is blank; do not equate yavol with Yev Barkalov.

\end{samepage}
\textbf{Source classification:} Baseline/known-result candidate, or incomplete claim with prior overlap; see exact evidence

\textbf{Prior comparison:} Candidate scoped mathematical delta; exact baseline and independently sealed provenance are required. A missing gallery value is not evidence of a world record.

\textbf{Formal/source status:} Only board metadata/access leads; no original mathematical proof or exact final packet. Observed run predates event; later eligible delta not established.

\textbf{Research or reuse implication:} Hold missing payload; current proposed points 0. Board metadata cannot establish an eligible exact theorem.

{\small \url{https://github.com/alejandrozu/openmath-2026-judging/blob/d0ed487f487fa7ec814ff705ab55249983fe8707/OpenMath-Judging/review/entries/E66-HILL1.txt}\par}

\begin{samepage}
\subsection{E66-HILL 2  -  alejandrozu/collatz-modular-descent}
{\small Contributors:  | Family: COLLATZ\par}

\textbf{Recorded scope:} Ramsey 0.030141721098 beats frozen reference, but dated September 21. Collatz 100\% weighted target coverage with three rules and BB6 246,872 steps are also pre-event. Identity is blank; do not equate yavol with Yev Barkalov.

\end{samepage}
\textbf{Source classification:} Unresolved or missing-artifact record; nonacceptance does not prove nonnovelty

\textbf{Prior comparison:} No eligible original mathematical delta established in the available packet. This is not a finding that all underlying work lacks novelty.

\textbf{Formal/source status:} Only board metadata/access leads; no original mathematical proof or exact final packet. Observed run predates event; later eligible delta not established.

\textbf{Research or reuse implication:} Hold missing payload; current proposed points 0. Board metadata cannot establish an eligible exact theorem.

{\small \url{https://github.com/alejandrozu/openmath-2026-judging/blob/d0ed487f487fa7ec814ff705ab55249983fe8707/OpenMath-Judging/review/entries/E66-HILL2.txt}\par}

\begin{samepage}
\subsection{E66-HILL 3  -  alejandrozu/busy-beaver-6-certificates}
{\small Contributors:  | Family: BB6\par}

\textbf{Recorded scope:} Ramsey 0.030141721098 beats frozen reference, but dated September 21. Collatz 100\% weighted target coverage with three rules and BB6 246,872 steps are also pre-event. Identity is blank; do not equate yavol with Yev Barkalov.

\end{samepage}
\textbf{Source classification:} Unresolved or missing-artifact record; nonacceptance does not prove nonnovelty

\textbf{Prior comparison:} No eligible original mathematical delta established in the available packet. This is not a finding that all underlying work lacks novelty.

\textbf{Formal/source status:} Only board metadata/access leads; no original mathematical proof or exact final packet. Observed run predates event; later eligible delta not established.

\textbf{Research or reuse implication:} Hold missing payload; current proposed points 0. Board metadata cannot establish an eligible exact theorem.

{\small \url{https://github.com/alejandrozu/openmath-2026-judging/blob/d0ed487f487fa7ec814ff705ab55249983fe8707/OpenMath-Judging/review/entries/E66-HILL3.txt}\par}

\begin{samepage}
\subsection{E02-a 046969-i  -  A046969 Conjecture I: exact Bernoulli denominator}
{\small Contributors: Rachel Teo, Wenyi Wang, Hamdi Abderrahmene | Family: M2-BERNOULLI\par}

\textbf{Recorded scope:} For a(n)=den(B\_(2n)/(2n(2n\ensuremath{-}1))) when n>0 (a(0)=0), if p>3 and both p and 2p\ensuremath{-}1 are prime, then a(p)=12(2p\ensuremath{-}1); consequently a(p)/12 is prime.

\end{samepage}
\textbf{Source classification:} Known-mathematics formalization / formal-library novelty and usefulness require verification

\textbf{Prior comparison:} Strong classical mathematical overlap; exact prior published ConjectureI theorem still unresolved Targeted Mathlib Bernoulli source inspection did not find this exact prime-pair formula, but found preexisting basic formal theory

\textbf{Formal/source status:} Packet evidence: Independent 5-module replay completed2026-10-01T13:09:42.764394Z, standard final formula and target axioms; exact source correspondence recorded.. Endpoint signatures and declared axioms are recorded below; author/internal replay is not this reviewer's fresh rebuild. Known theorem route: only genuinely new faithful useful reusable formal delta can qualify.

\textbf{Research or reuse implication:} Propose this exact scope and family once, subject to the stated source, formal, novelty and deadline checks.

{\small \url{https://github.com/alejandrozu/openmath-2026-judging/blob/d0ed487f487fa7ec814ff705ab55249983fe8707/OpenMath-Judging/review/entries/E02-a046969-i.txt}\par}

\begin{samepage}
\subsection{E02-opdp67-quaternion  -  OPDP 67 explicit two-witness quaternion scalar definition}
{\small Contributors: Rachel Teo, Wenyi Wang, Hamdi Abderrahmene | Family: M2-QUATERNION\par}

\textbf{Recorded scope:} For every linearly ordered commutative ring R satisfying CommRing, LinearOrder and IsStrictOrderedRing, a Hamilton quaternion q over R is scalar iff there exist unrestricted quaternions u,v with ((uv\ensuremath{-}vu)\textasciicircum{}2+4)\textasciicircum{}2\ensuremath{-}(qu\ensuremath{-}uq)\textasciicircum{}2\ensuremath{-}(qv\ensuremath{-}vq)\textasciicircum{}2=0. Specializations define scalar Z in the Lipschitz order and scalar Q in rational Hamilton quaternions.

\end{samepage}
\textbf{Source classification:} Known-mathematics formalization / formal-library novelty and usefulness require verification

\textbf{Prior comparison:} Known mathematical outcome; no original-open claim Useful formalization is plausible, not guaranteed by successful local compilation or absence of an exact name in a narrow search The latest 8-file pinned Mathlib scan distinguishes preexisting algebra/norm/commutation APIs from the explicit existential equation. It is not an exhaustive semantic search of every proof assistant or repository.

\textbf{Formal/source status:} Packet evidence: Recorded two-witness module build and explicit final generic/rational/integer axiom audit; only standard axioms. No independent clean second replay claimed. Endpoint signatures and declared axioms are recorded below; author/internal replay is not this reviewer's fresh rebuild. Known theorem route: only genuinely new faithful useful reusable formal delta can qualify.

\textbf{Research or reuse implication:} Propose this exact scope and family once, subject to the stated source, formal, novelty and deadline checks.

{\small \url{https://github.com/alejandrozu/openmath-2026-judging/blob/d0ed487f487fa7ec814ff705ab55249983fe8707/OpenMath-Judging/review/entries/E02-opdp67-quaternion.txt}\par}

\begin{samepage}
\subsection{E02-a 135508  -  A135508: prime-index gcd criterion and residue classes}
{\small Contributors: Rachel Teo, Wenyi Wang, Hamdi Abderrahmene | Family: M2-LCM-RECURRENCE\par}

\textbf{Recorded scope:} Set x(0)=0, x(1)=1 and x(n+1)=2x(n)+lcm(x(n),n+1) for n>=1; for n>0 let a(n)=x(n+1)/x(n)\ensuremath{-}2. For each prime p>3, a(p\ensuremath{-}1)=p iff gcd(x(p\ensuremath{-}3),p\ensuremath{-}2)>1. If p\ensuremath{\geq}7 is prime and p mod 3=2 or p mod 5=2, then a(p\ensuremath{-}1)=p. More generally an acquired divisor d>1 dividing x(m) and p\ensuremath{-}2, with 1\ensuremath{\leq}m\ensuremath{\leq}p\ensuremath{-}3, suffices.

\end{samepage}
\textbf{Source classification:} Known-mathematics formalization / formal-library novelty and usefulness require verification

\textbf{Prior comparison:} Prior conditional full result and local sufficient unconditional claims must be compared at hypothesis level Classical/elementary specialization does not alone settle modality

\textbf{Formal/source status:} Packet evidence: Strict 4.33.1 replay passed 2026-10-02T01:22:25Z. Four in-source axiom reports standard; later expanded-axioms.log also explicitly audits the acquired-divisor theorem with standard axioms.. Endpoint signatures and declared axioms are recorded below; author/internal replay is not this reviewer's fresh rebuild. Known theorem route: only genuinely new faithful useful reusable formal delta can qualify.

\textbf{Research or reuse implication:} Propose this exact scope and family once, subject to the stated source, formal, novelty and deadline checks.

{\small \url{https://github.com/alejandrozu/openmath-2026-judging/blob/d0ed487f487fa7ec814ff705ab55249983fe8707/OpenMath-Judging/review/entries/E02-a135508.txt}\par}

\begin{samepage}
\subsection{E02-apery-odd-formalization  -  A005258 Apéry prime-power irreducibility}
{\small Contributors: Rachel Teo, Wenyi Wang, Hamdi Abderrahmene | Family: M2-APERY\par}

\textbf{Recorded scope:} For B\_n(X)=\ensuremath{\Sigma}\_(k=0)\textasciicircum{}n binom(n,k)\textasciicircum{}2 binom(n+k,k)X\textasciicircum{}k in Q[X], every odd prime p and r\ensuremath{\geq}1 satisfying p\ensuremath{^2} \ensuremath{\not\mid} 2\textasciicircum{}(p\ensuremath{-}1)\ensuremath{-}1 has Irreducible(B\_(p\textasciicircum{}r)). This card requests review only of any genuinely new useful formalization beyond the prior exposed mathematical sketch; it asserts no original-open novelty.

\end{samepage}
\textbf{Source classification:} Known-mathematics formalization / formal-library novelty and usefulness require verification

\textbf{Prior comparison:} Completed prior binary chain is not invalidated by the unresolved full theorem's sorry Odd-family overlap is an informal sketch, not an earlier verified odd-family Lean endpoint The prior higher-valuation formula is incorrect/unjustified and is not adopted; valid non-Wieferich sufficient implication survives Earlier no-exact-prior reports are superseded by the 2Oct positive overlap

\textbf{Formal/source status:} Packet evidence: Seven modules checked successfully on both pins; retained successful first-five logs and patched final-two logs are stated honestly. Source hashes match both receipts; final target standard axioms.. Endpoint signatures and declared axioms are recorded below; author/internal replay is not this reviewer's fresh rebuild. Known theorem route: only genuinely new faithful useful reusable formal delta can qualify.

\textbf{Research or reuse implication:} Propose this exact scope and family once, subject to the stated source, formal, novelty and deadline checks.

{\small \url{https://github.com/alejandrozu/openmath-2026-judging/blob/d0ed487f487fa7ec814ff705ab55249983fe8707/OpenMath-Judging/review/entries/E02-apery-odd-formalization.txt}\par}

\begin{samepage}
\subsection{E02-m2-research-base3-reverse-add  -  Base three reverse and add carry obstruction}
{\small Contributors: Rachel Teo, Wenyi Wang, Hamdi Abderrahmene | Family: M2-BASE 3\par}

\textbf{Recorded scope:} If x plus its canonical base-3 digit reversal has palindromic base-3 digits, then x already has palindromic base-3 digits or every pair of mirrored input digits sums to less than 3.

\end{samepage}
\textbf{Source classification:} Known-mathematics formalization / formal-library novelty and usefulness require verification

\textbf{Prior comparison:} Classical carry/symmetry context is known Potential reusable uniform digit/carry formalization; exact target admission and formal-library freshness remain review questions

\textbf{Formal/source status:} Packet evidence: The additional exact-source Lean 4.33.1 replay passed all 235 own-source dependency modules and seven independent selected proof files. Individual source hashes and selected endpoint standard-axiom reports are retained. Original archival passes remain separate. The source replay continued across scheduling restarts using hash-matched successful dependency receipts and exact pinned external-library caches. Each proof ran in a separate process; no aggregate importability, independent compiler implementation or official checker acceptance is asserted. Endpoint signatures and declared axioms are recorded below; author/internal replay is not this reviewer's fresh rebuild. Known theorem route: only genuinely new faithful useful reusable formal delta can qualify.

\textbf{Research or reuse implication:} Propose this exact scope and family once, subject to the stated source, formal, novelty and deadline checks.

{\small \url{https://github.com/alejandrozu/openmath-2026-judging/blob/d0ed487f487fa7ec814ff705ab55249983fe8707/OpenMath-Judging/review/entries/E02-m2-research-base3-reverse-add.txt}\par}

\begin{samepage}
\subsection{E02-m2-research-theta-half-sum  -  Theta half sum negative inverse cancellation}
{\small Contributors: Rachel Teo, Wenyi Wang, Hamdi Abderrahmene | Family: M2-THETA\par}

\textbf{Recorded scope:} For an odd prime k and odd natural a,b with 1 \ensuremath{\leq} a,b < k and ab mod k = k\ensuremath{-}1, define H(k,a)=\ensuremath{\Sigma} from j=1 to (k\ensuremath{-}1)/2 of (\ensuremath{-}1)\textasciicircum{}floor(aj/k). Then H(k,a)+H(k,b)=0.

\end{samepage}
\textbf{Source classification:} Known-mathematics formalization / formal-library novelty and usefulness require verification

\textbf{Prior comparison:} Classical Hardy/Dedekind identities explain the mathematics; source-based derivation documented Direct centered-residue permutation and parity proof supplies the formalization under review

\textbf{Formal/source status:} Packet evidence: The additional exact-source Lean 4.33.1 replay passed all 235 own-source dependency modules and seven independent selected proof files. Individual source hashes and selected endpoint standard-axiom reports are retained. Original archival passes remain separate. The source replay continued across scheduling restarts using hash-matched successful dependency receipts and exact pinned external-library caches. Each proof ran in a separate process; no aggregate importability, independent compiler implementation or official checker acceptance is asserted. Endpoint signatures and declared axioms are recorded below; author/internal replay is not this reviewer's fresh rebuild. Known theorem route: only genuinely new faithful useful reusable formal delta can qualify.

\textbf{Research or reuse implication:} Propose this exact scope and family once, subject to the stated source, formal, novelty and deadline checks.

{\small \url{https://github.com/alejandrozu/openmath-2026-judging/blob/d0ed487f487fa7ec814ff705ab55249983fe8707/OpenMath-Judging/review/entries/E02-m2-research-theta-half-sum.txt}\par}

\begin{samepage}
\subsection{E02-m2-research-wolstenholme  -  Reciprocal square numerator arithmetic}
{\small Contributors: Rachel Teo, Wenyi Wang, Hamdi Abderrahmene | Family: M2-WOLSTENHOLME\par}

\textbf{Recorded scope:} For every m \ensuremath{\geq} 1, the rational sum \ensuremath{\Sigma} from k=1 to m of 1/k\ensuremath{^2} has nonpositive 2-adic and 3-adic valuations. The selected bridge excludes a prime from its numerator when the valuation is nonpositive; for every prime p>3, gcd(p,A007406(p\ensuremath{-}1))=p.

\end{samepage}
\textbf{Source classification:} Known-mathematics formalization / formal-library novelty and usefulness require verification

\textbf{Prior comparison:} Pinned Mathlib has ordinary harmonic 2-adic valuation theory, not the same reciprocal-square result The equal-minimum 3-adic noncancellation and numerator/prime bridge are the principal selected formal content

\textbf{Formal/source status:} Packet evidence: The additional exact-source Lean 4.33.1 replay passed all 235 own-source dependency modules and seven independent selected proof files. Individual source hashes and selected endpoint standard-axiom reports are retained. Original archival passes remain separate. The source replay continued across scheduling restarts using hash-matched successful dependency receipts and exact pinned external-library caches. Each proof ran in a separate process; no aggregate importability, independent compiler implementation or official checker acceptance is asserted. Endpoint signatures and declared axioms are recorded below; author/internal replay is not this reviewer's fresh rebuild. Known theorem route: only genuinely new faithful useful reusable formal delta can qualify.

\textbf{Research or reuse implication:} Propose this exact scope and family once, subject to the stated source, formal, novelty and deadline checks.

{\small \url{https://github.com/alejandrozu/openmath-2026-judging/blob/d0ed487f487fa7ec814ff705ab55249983fe8707/OpenMath-Judging/review/entries/E02-m2-research-wolstenholme.txt}\par}

\begin{samepage}
\subsection{E02-m2-research-a060841  -  A060841 reciprocal determinant integrality}
{\small Contributors: Rachel Teo, Wenyi Wang, Hamdi Abderrahmene | Family: M2-LCM-MATRIX\par}

\textbf{Recorded scope:} For n \ensuremath{\geq} 1, the inverse determinant of the n by n matrix with entries 1/lcm(i+1,j+1) is an integer exactly when n belongs to \{1,...,34,36,38\}.

\end{samepage}
\textbf{Source classification:} Known-mathematics formalization / formal-library novelty and usefulness require verification

\textbf{Prior comparison:} Substantive determinant-to-product and valuation bridge beyond located public product-only development Exact informal classification was already public; product/valuation helpers have prior formal overlap

\textbf{Formal/source status:} Packet evidence: The additional exact-source Lean 4.33.1 replay passed all 235 own-source dependency modules and seven independent selected proof files. Individual source hashes and selected endpoint standard-axiom reports are retained. Original archival passes remain separate. The source replay continued across scheduling restarts using hash-matched successful dependency receipts and exact pinned external-library caches. Each proof ran in a separate process; no aggregate importability, independent compiler implementation or official checker acceptance is asserted. Endpoint signatures and declared axioms are recorded below; author/internal replay is not this reviewer's fresh rebuild. Known theorem route: only genuinely new faithful useful reusable formal delta can qualify.

\textbf{Research or reuse implication:} Propose this exact scope and family once, subject to the stated source, formal, novelty and deadline checks.

{\small \url{https://github.com/alejandrozu/openmath-2026-judging/blob/d0ed487f487fa7ec814ff705ab55249983fe8707/OpenMath-Judging/review/entries/E02-m2-research-a060841.txt}\par}

\begin{samepage}
\subsection{E02-m2-research-erdos887  -  Erdős and Rosenfeld four nearby divisors}
{\small Contributors: Rachel Teo, Wenyi Wang, Hamdi Abderrahmene | Family: M2-ERDOS 887\par}

\textbf{Recorded scope:} There are infinitely many natural n with at least four distinct divisors d in the open real interval (sqrt(n), sqrt(n)+40 n\textasciicircum{}(1/4)). The preserved statement uses equivalent floor/ceiling integer endpoints.

\end{samepage}
\textbf{Source classification:} Known-mathematics formalization / formal-library novelty and usefulness require verification

\textbf{Prior comparison:} Eight-consecutive-factor construction with explicit divisors and interval inequalities Located prior constant-16 benchmark statement has a proof hole; located axiomatized work is not a closed prior proof; bounded search is not a universal novelty certificate

\textbf{Formal/source status:} Packet evidence: The additional exact-source Lean 4.33.1 replay passed all 235 own-source dependency modules and seven independent selected proof files. Individual source hashes and selected endpoint standard-axiom reports are retained. Original archival passes remain separate. The source replay continued across scheduling restarts using hash-matched successful dependency receipts and exact pinned external-library caches. Each proof ran in a separate process; no aggregate importability, independent compiler implementation or official checker acceptance is asserted. Endpoint signatures and declared axioms are recorded below; author/internal replay is not this reviewer's fresh rebuild. Known theorem route: only genuinely new faithful useful reusable formal delta can qualify.

\textbf{Research or reuse implication:} Propose this exact scope and family once, subject to the stated source, formal, novelty and deadline checks.

{\small \url{https://github.com/alejandrozu/openmath-2026-judging/blob/d0ed487f487fa7ec814ff705ab55249983fe8707/OpenMath-Judging/review/entries/E02-m2-research-erdos887.txt}\par}

\begin{samepage}
\subsection{E02-m2-research-regular-polygon  -  Regular polygon chord and distance multiplicity}
{\small Contributors: Rachel Teo, Wenyi Wang, Hamdi Abderrahmene | Family: M2-POLYGON\par}

\textbf{Recorded scope:} For every even n \ensuremath{\geq} 5 and every distance d occurring between distinct vertices of the regular n-gon, the number of unordered vertex pairs at distance d is at most n.

\end{samepage}
\textbf{Source classification:} Known-mathematics formalization / formal-library novelty and usefulness require verification

\textbf{Prior comparison:} Generic two-circle-intersection uniqueness is already in pinned Mathlib The selected case is the complete regular-polygon geometry/counting family; the final inequality or proof length alone does not certify novelty

\textbf{Formal/source status:} Packet evidence: The additional exact-source Lean 4.33.1 replay passed all 235 own-source dependency modules and seven independent selected proof files. Individual source hashes and selected endpoint standard-axiom reports are retained. Original archival passes remain separate. The source replay continued across scheduling restarts using hash-matched successful dependency receipts and exact pinned external-library caches. Each proof ran in a separate process; no aggregate importability, independent compiler implementation or official checker acceptance is asserted. Endpoint signatures and declared axioms are recorded below; author/internal replay is not this reviewer's fresh rebuild. Known theorem route: only genuinely new faithful useful reusable formal delta can qualify.

\textbf{Research or reuse implication:} Propose this exact scope and family once, subject to the stated source, formal, novelty and deadline checks.

{\small \url{https://github.com/alejandrozu/openmath-2026-judging/blob/d0ed487f487fa7ec814ff705ab55249983fe8707/OpenMath-Judging/review/entries/E02-m2-research-regular-polygon.txt}\par}

\begin{samepage}
\subsection{E02-FOCUS-GROTH  -  Grothendieck: restricted sign-matrix bounds}
{\small Contributors: Rachel Teo, Wenyi Wang, Hamdi Abderrahmene | Family: GROTHENDIECK\par}

\textbf{Recorded scope:} For every n and every 3-by-n sign matrix A, Q\_V(A)<=sqrt(3/2) C(A) in every real inner-product space V. Additionally Q\_V(A)<=sqrt(2) C(A) for the four literal archived matrices A\_6x6\_106, A\_6x8\_690, A\_7x8\_53110 and A\_s8x8\_1705. Q\_V is the unit-vector objective supremum, with Q\_V(A)=0 for an empty feasible set, and C is its real scalar specialization. These supporting cases form one review family.

\end{samepage}
\textbf{Source classification:} Unresolved or missing-artifact record; nonacceptance does not prove nonnovelty

\textbf{Prior comparison:} No original universal-constant advance established. A new formalization may be M2 only after source and prior-library comparison.

\textbf{Formal/source status:} Independent pinned Lean 4.33.1 replay completed both exact source modules Row3Granular and SemanticTestGranular and all five selected endpoint prints, with standard kernel axioms or no axioms; no native, admitted or unknown axiom is present in the selected scope. Row 3 body remains byte-identical apart from the documented import adjustment; the four-case file retains statements and certificate data with eight documented arithmetic tactic repairs and explicit imports. Restricted upper-bound scope and originality hold remain unchanged.

\textbf{Research or reuse implication:} Open score 0; keep supporting formal material, with no automatic extra M2 family.

{\small \url{https://github.com/alejandrozu/openmath-2026-judging/blob/d0ed487f487fa7ec814ff705ab55249983fe8707/OpenMath-Judging/review/entries/E02-FOCUS-GROTH.txt}\par}

\begin{samepage}
\subsection{E02-FOCUS-BB6  -  Busy Beaver 6: exact halting witness}
{\small Contributors: Rachel Teo, Wenyi Wang, Hamdi Abderrahmene | Family: BB6\par}

\textbf{Recorded scope:} Starting on a bi-infinite blank tape in state A, the six-state two-symbol machine 1RB0LC\_1RD1RB\_0LH1LE\_1LA1LD\_0LA1LF\_0LE0RD executes exactly 1,235,211 transitions to H with 1,519 ones on its final tape and no earlier halted prefix. Every state A-F is visited by step 40. Separate exact simulators give final head -2001 and visited span 2,028 (minimum -2001, maximum 26, including the post-halting head); head and span are not kernel-certified claims.

\textbf{Source classification:} Baseline/known-result candidate, or incomplete claim with prior overlap; see exact evidence

\textbf{Prior comparison:} No new original Busy Beaver frontier demonstrated; a reusable newly formalized simulator would require an explicit M2 claim and prior-library check.

\textbf{Formal/source status:} A new standard-kernel certificate in Lean 4.33.1 successfully checks the literal machine, exact run, final ones count, absence of an earlier halt and all-six-states visitation. Claims.lean endpoints exact\_halting\_witness and visits\_all\_six\_states have only propext and Quot.sound in their printed axiom sets. The 129-source-module closure uses no native\_decide, custom axiom or admission, and requires no Mathlib. Semantics.lean proves the tape-cell equations for both move directions. Original dictionary- and set-tape simulator records are preserved as independent evidence.

\textbf{Research or reuse implication:} Open score 0. Do not count steps or ones as a fraction of resolving BB6.

{\small \url{https://github.com/alejandrozu/openmath-2026-judging/blob/d0ed487f487fa7ec814ff705ab55249983fe8707/OpenMath-Judging/review/entries/E02-FOCUS-BB6.txt}\par}

\end{samepage}
\end{document}
